\documentclass[a4paper]{article}
\usepackage{lmodern}
\usepackage[T1]{fontenc}
\usepackage[utf8]{inputenc}
\usepackage{array}
\usepackage{booktabs}
\usepackage{amssymb,amsmath,amsthm,mathtools}
\usepackage{enumitem}
\usepackage{xcolor}
\usepackage[margin=1in]{geometry}
\usepackage[colorlinks=true,linkcolor=blue!50!black,citecolor=blue!50!black,urlcolor=blue!50!black]{hyperref}
\usepackage{longtable}
\usepackage{float}
\usepackage{microtype}

\theoremstyle{plain}
\newtheorem{theorem}{Theorem}[subsection]
\newtheorem{thm}[theorem]{Theorem}
\newtheorem{proposition}{Proposition}[subsection]
\newtheorem{lemma}{Lemma}[subsection]
\newtheorem{corollary}{Corollary}[subsection]
\theoremstyle{definition}
\newtheorem{definition}{Definition}[subsection]
\newtheorem{example}{Example}[subsection]
\theoremstyle{remark}
\newtheorem{remark}{Remark}[subsection]
\newtheorem{observation}{Observation}[subsection]
\newtheorem{conjecture}{Conjecture}[subsection]
\newtheorem{question}{Question}[subsection]

\newcommand{\Z}{\mathbb{Z}}
\newcommand{\F}{\mathbb{F}}
\newcommand{\ord}{\operatorname{ord}}
\newcommand{\gen}[1]{\langle #1 \rangle}
\DeclareMathOperator{\lcm}{lcm}
\DeclareMathOperator{\popcount}{w}
\newcommand{\dvd}{\mid}
\newcommand{\Sa}{S_a}
\newcommand{\half}{\tfrac12}

% --- macros and theorem environments for the Avoidance paper (Jackopedia) ---
\newtheorem{construction}{Construction}[subsection]
\newtheorem{openproblem}{Open Problem}[subsection]
\DeclareMathOperator{\val}{val}
\DeclareMathOperator{\len}{len}
\newcommand{\base}[2]{(#1)_{#2}}
\newcommand{\good}{\textsf{good}}
\newcommand{\bad}{\textsf{bad}}
\newcommand{\NN}{\mathbb{N}}
\newcommand{\QQ}{\mathbb{Q}}
\newcommand{\ZZ}{\mathbb{Z}}
\newcommand{\Lbb}[1][b,B]{L_{#1}}
\newcommand{\Sbb}[1][b,B]{S_{#1}}
\newcommand{\hbb}[1][b,B]{h_{#1}}

% --- float placement: keep tables near their text, avoid dumping at end ---
\setcounter{topnumber}{4}
\setcounter{bottomnumber}{4}
\setcounter{totalnumber}{8}
\renewcommand{\topfraction}{0.9}
\renewcommand{\bottomfraction}{0.9}
\renewcommand{\textfraction}{0.06}
\renewcommand{\floatpagefraction}{0.7}

\title{Project Ladislaus}
\author{NOTXXDOG \and ROBIN ROVENSZKY}
\date{Since May 2026}

\begin{document}

\maketitle
\begin{abstract}
    Having a blast with precisely 2 interesting problems, plus interesting holdouts from the central domains of the \emph{Busy Beaver function}.
\end{abstract}

\tableofcontents

\section{Introduction}
\emph{Genuinely} working on random (all) math problems shared on \emph{bbchallenge}.

\section{The Problems}

\renewcommand{\arraystretch}{1.5}

\begin{center}
    \begin{tabular}{
        >{\centering\arraybackslash}m{4cm}
        !{\vrule}
        >{\centering\arraybackslash}m{1cm}
        !{\vrule}
        >{\centering\arraybackslash}m{3.75cm}
        !{\vrule}
        >{\centering\arraybackslash}m{1.5cm}
        !{\vrule}
        >{\centering\arraybackslash}m{1.5cm}
        !{\vrule}
        >{\centering\arraybackslash}m{1.5cm}
    }
    Problem & Seen & Progress & Solved & Formal & General \\
    \cline{1-6}
    BMO Solved & Yes & 3, 4 solved & 3, 4, not 7 & 3, 4, not 7 & NSI \\
    \cline{1-6}
    Binary String Rewriting Process (Iterative Processes) & Yes & Exhausted & No & Yes & General $L_0$ \\
    \cline{1-6} 
    Division Problem on Powers of Two (Modularity) & Yes & Exhausted & No & Yes & General $S_a$ \\
    \cline{1-6}
    BMO Unsolved & No & --- & --- & --- & --- \\
    \cline{1-6}
    Genuinely Unappreciated BB(6) Holdouts & No & --- & --- & --- & --- \\
    \cline{1-6}
    BB(3x3) Holdouts & No & --- & --- & --- & --- \\
    \cline{1-6}
    BB(2x5) Holdouts & No & --- & --- & --- & ---
    \end{tabular}
\end{center}
\renewcommand{\arraystretch}{1}
\emph{seen means looked at, appreciated or unappreciated its existence, then in the second case logically concluded not diving deeply, not scrutinizing would be the best choice in regards to mental health.} \\\\
\emph{note that a check on solved usually means ongoing formalisation, and a check on formalisation usually means ongoing generalisation.} \\\\
\emph{NSI means No Strong Incentive, while SI means Strong Incentive}

\section{Motivation, Modularity and Iterative Processes Bias}
Since I'm incredibly biased towards Iterative Processes and Modularity, the two particular problems regarding these are of high interest to generalise, as another project of mine is to scrutinize what makes some of these problems notoriously hard to solve while easy to describe, and why in particular the gap between super easy to solve and intractable is so tight.

I have half-assed this during a Mazsó yap class, and reached the following conclusions:

\subsection{Modularity}
\textbf{Any problem involving division.}\\

First step: generalise to $S_q$, not just $S_{-3}$.
Second step: generalise to $k^n$ instead of just $2^n$. 
Third step: generalise to (mod n + a), not just (mod n + 1).
Now we are at $k^n$ congruent to q (mod n + a) (all variables in Z)
Fourth step: generalise to $k^n$ congruent to q (mod b + a) (all variables in Z)
Now we have reached the far general case that analyses exponentiation of general variables vs. addition of general variables, but restricted to Z.
Finally, generalise to C.

We can then create a complete taxonomy of division-related problems: restrict only to C (or perhaps some -adic family? good question.)
With this problem, we will achieve exponentiation vs. addition equaling a general number. We can then further generalise to all combinations of the followings: addition, subtraction, multiplication, division, exponentiation, square roots (+ logarithms?) and then we will have to study how to generalise further.
We may also see if its WLOG to not equal a general number, but also some order of operations. The hope is that then combining more operations is WLOG (i.e. addition and subtraction vs. division equaling general number).

\subsection{Iterative Processes}
\textbf{Any evolving system.}\\
\emph{Perhaps it would be best to call it Evolving Systems, Iterative Processes for ones that just iterate and not evolve in the meantime and then Functional Iterative Processes for ones like Collatz?}\\

Taxonomy of iterative processes: defining them in most general form in first go would be a difficult question, therefore we proceed as we did with modularity:
Start from Collatz. We can generalise $f(n)_{a,b} = n/2$ if even 3n+1 if odd to $f(n)_{a,b} = n/2$ if even, an+b if odd.
We can proceed with generalising further: `Let n/q instead of n/2 and let there be 2q params to f : Z -> Z where params come in pairs $a_1 b_1, a_2 b_2, a_3 b_3$. They let each rule be defined when n is congruent to some k < q != 0.` This is the last that I would call a member of the Collatz iterative family.

Then we have reached a point where we are taking an iterative function which has a dependency (please come up with a better name i beg you) that the function decides based on to apply one of the rules which are of form an+b, or one other rule which is of the form n/q. 

.... continuing this here until the conclusion should be imagined by some schizophrenic person that has fucking too much time. im limited by Discord anyway so fuckyeah

We reach the conclusion therefore that iterative processes are any function f : S -> S (for some ultimate sets S that we don't conclude, but could be C, i am legitimately too fuclking stupid for this atm) that we observe during several repetitions, i.e, $f^{(t)}$. We will see if I can reach any better conclusion (i.e. whether observing two different functions f : S -$> f_{2}(S)$ or  the fractal special case f : S -> f(S) etc., or something else that im way too tired to realize atp)

But I think it may be more general to think of iterative processes as any arbitrary process, iterated, literally... therefore any system under evolution including dynamical maps and whatnots, neural nets etc. belong to this rich family of beauty, and then we can name the above conclusion algebraic or functional iterative processes, because these names are genuinely so cool. fucklking


\section{Binary-string rewriting problem}

\subsection{Paper 1, Special Case ($L_0 = 10$)}

\begin{abstract}
We study a deterministic rewriting process $F$ on finite binary strings that carries an integer counter $s$ and halts precisely when $s$ falls below $2$; the central question is whether the trajectory generated from the seed $L_0=10$ is infinite.  We reduce this question to a scalar inequality on the counter sequence $(s_n)$ and analyse it from several directions.  The halting value $s=0$ is excluded unconditionally at every even index by an exact prefix invariant $(r_1(L_n)=1\iff n\text{ even})$, and at odd indices under the mild bound $s_{n-2}\ge 4$.  We prove the two-step monotonicity $s_n\ge s_{n-2}$ at every even-parity step whose predecessor is also of even parity, via an even-step floor $s_n\ge 3\lfloor(s_{n-2}-3)/2\rfloor$, and we give an exact recurrence for the $1$-count $\nu(L_n)$ that fixes the parity sequence in closed form.  A counter-inert tail lemma yields an exact decomposition of the even-step counter into a count term and an alignment-filtered size surplus, and shows that the count alone provably cannot close the residual gap.  The map is recast as a faithful integer conjugate of Collatz type, as a gap-vector rewrite, and as a scalar process with no closed recurrence.  Two independent from-scratch engines yield the counter through $s_{26}=88\,613\,778$ and the parity through $\nu(L_{26})=222\,569\,758$, confirm two-step monotonicity throughout $3\le n\le 26$, and localize the single remaining obstruction to even-parity steps with odd-parity predecessors.  A dedicated section explains, with the forward and backward routes side by side, why the residual inequality resists elementary methods.  The central theorem and the full monotonicity conjecture remain open.
\end{abstract}


\subsubsection{Introduction}
The process studied here is easy to define and stubbornly resistant to complete analysis. One iterates a fixed rewriting rule on binary strings; the rule carries a small integer counter and is engineered to stop the iteration the instant that counter drops below the threshold~$2$. From the seed $L_0=10$ the iteration appears to run forever, with the counter exploding to infinity after a single, isolated brush with the threshold. Proving that it truly never stops is the problem.

The architecture of the analysis is the following: modulo one structural fact---that a particular halting value never occurs on the orbit of $10$---the whole question is equivalent to a single inequality about the size of the counter, driven by a transparent geometric mechanism. We establish that structural fact, prove the governing inequality in a structurally identified sub-case, push the computational frontier outward, and pin down sharply the one configuration on which a complete proof still founders.

\subsubsection*{Conventions} Strings are finite words over $\{0,1\}$; positions are numbered $1,2,3,\dots$ from the left. For $w\in\{0,1\}^*$ we write $\nu(w)$ for the number of $1$'s and $\omega(w)$ for the number of odd-length maximal runs of $1$'s. Concatenation is denoted by juxtaposition or a dot, and $u^k$ is the $k$-fold concatenation of $u$. We call a string \emph{even} (resp.\ \emph{odd}) when $\nu$ is even (resp.\ odd); this ``parity'' refers to $\nu$, not to the index $n$ (Table~\ref{b1:tab:traj}).

\begin{definition}[The map $F$]\label{b1:def:F}
For $L\in\{0,1\}^*$, the value $F(L)$, when defined, is computed in three steps.

\emph{Step 1 (pad; seed the counter).} Let $k=\nu(L)$. If $k$ is even, append $00$ to $L$ and set $s:=0$; if $k$ is odd, append $0001$ to $L$ and set $s:=-1$. Call the result $L^{(1)}$; note $\nu(L^{(1)})$ is even.

\emph{Step 2 (pair the $1$'s; fill and count).} Let $p_1<\dots<p_{2m}$ be the positions of the $1$'s in $L^{(1)}$, grouped into consecutive pairs $(p_1,p_2),\dots,(p_{2m-1},p_{2m})$. For $j=1,\dots,m$ in turn: set positions $p_{2j-1},\dots,p_{2j}-1$ to $1$; set position $p_{2j}$ to $0$; and add $p_{2j}-p_{2j-1}-1$ to $s$. Call the result $L^{(2)}$.

\emph{Step 3 (halt or extend).} If $s<2$, the process halts and $F(L)$ is undefined. Otherwise $F(L):=L^{(2)}\cdot(0110)^{s-2}$.
\end{definition}

Informally, Step 2 takes the $1$'s two at a time, fills the gap between each pair into a solid block of $1$'s, deletes the block's right endpoint, and tallies the filled gap; the counter $s$ is the total filling. Step 1 makes the count of $1$'s even so that the pairing is exhaustive, and Step 3 reads the total filling and, when there was enough, glues on a tail whose length grows with $s$.

\begin{definition}[Trajectory]\label{b1:def:traj}
Set $L_0:=10$ and $L_{n+1}:=F(L_n)$ whenever the right-hand side is defined. Write $s_n$ for the counter computed during $F(L_n)$.
\end{definition}

\begin{thm}[Non-halting; the central claim, open]\label{b1:thm:central}
$L_n$ is defined for every $n\ge 0$; equivalently, $s_n\ge 2$ for all $n$.
\end{thm}

Theorem~\ref{b1:thm:central} is not proved here, and we believe it is not within reach of elementary methods. What we provide is a reduction to a single scalar inequality and a sharply delimited residual obstruction; \S\ref{b1:sec:status} states the status with precision.

\subsubsection{Summary of results} The main results are the following.
\begin{enumerate}[label=(\roman*)]
\item A prefix-alternation lemma (Lemma~\ref{b1:lem:prefix}): $r_1(L_n)=1\iff n$ even, where $r_1$ is the length of the first maximal $1$-run.
\item Exclusion of the all-even halting profile $s_n=0$ (Proposition~\ref{b1:prop:alleven}): unconditionally at every even index and, at odd indices, whenever $s_{n-2}\ge 4$.
\item An even-step floor and a proven sub-case of two-step monotonicity (Theorem~\ref{b1:thm:floor}): $s_n\ge 3\lfloor(s_{n-2}-3)/2\rfloor$, whence $s_n\ge s_{n-2}$ once $s_{n-2}\ge 12$.
\item An exact recurrence for $\nu(L_n)$ (Proposition~\ref{b1:prop:nu}) which closes the parity dynamics.
\item A counter-inert tail lemma (Lemma~\ref{b1:lem:inert}) and an exact body localization of the even-step counter (Observation~\ref{b1:obs:decomp}): $s_n$ is a functional of the body $L^{(2)}_{n-1}$ alone, and decomposes as a count term plus an alignment-filtered size surplus.
\item Three exact reformulations (\S\ref{b1:sec:equiv}): a faithful integer conjugate $T$ (Theorem~\ref{b1:thm:T}) built on the identity that Step~2 is the doubled alternating bit-sum (Lemma~\ref{b1:lem:S}); a natural-scale gap-vector conjugate (Proposition~\ref{b1:prop:gap}); and the scalar-counter form (Proposition~\ref{b1:prop:scalar}).
\item An extended, independently cross-validated computation. The record reaches $s_{26}=88\,613\,778$, corroborated by $\nu(L_{26})=222\,569\,758$; two-step monotonicity $s_n\ge s_{n-2}$ holds for $3\le n\le 26$; a run of six consecutive even-parity steps occurs across which the counter nevertheless grows; the open obstruction is localized to even steps with odd-parity predecessors (\S\ref{b1:sec:12.1}).
\item A discussion (\S\ref{b1:sec:resist}) of why the problem resists elementary methods, placing the forward and backward routes side by side.
\end{enumerate}

\subsubsection{Elementary structure and the worked trajectory}\label{b1:sec:elem}
Running the rule once by hand fixes ideas. From $L_0=10$ the single $1$ makes $k=1$ odd, so Step 1 appends $0001$ and sets $s=-1$, giving $L^{(1)}=100001$. The $1$'s lie at positions $1$ and $6$; filling positions $1$--$5$ and zeroing position $6$ yields $111110$, and the filled gap is $4$, so $s=3$. As $3\ge 2$ we append $(0110)^1$ and read off $L_1=1111100110$. Continuing,
\begin{align*}
L_2 &= 10101110111110011001100110,\\
L_3 &= 1100101101010001000100010000,\\
L_4 &= 100011011001111000011110000000011001100110011001100110,
\end{align*}
and the counter takes the values $3,5,2,8,31,66,158,\dots$ (Table~\ref{b1:tab:traj}).

Two features organize everything. First, the sequence is \emph{explosive}: from $s_3$ onward it roughly multiplies each step, and $|L_n|$ exceeds $10^8$ before $n=22$. Second, it \emph{grazes the threshold exactly once}: $s_2=2$, with no margin. Any proof must survive that single near miss at $n=2$ and then explain why it never recurs.

\begin{lemma}[$F$ strictly lengthens]\label{b1:lem:length}
Whenever $F(L)$ is defined, $|F(L)|\ge|L|+2$.
\end{lemma}
\begin{proof}
Step 1 appends two or four symbols; Step 2 preserves length; Step 3 appends $4(s-2)\ge 0$ symbols.
\end{proof}

A first consequence, central to the backward viewpoint (\S\ref{b1:sec:backward}): every string has only finitely many iterated preimages under $F$, since a preimage is strictly shorter; orbit membership is decidable.

\begin{lemma}[Outputs end in $0$]\label{b1:lem:end0}
Whenever $F(L)$ is defined, its last symbol is $0$.
\end{lemma}
\begin{proof}
Step 2 sets the position $p_{2m}$ of the last $1$ of $L^{(1)}$ to $0$; no later position holds a $1$, so $L^{(2)}$ ends in $0$. If $s=2$ then $F(L)=L^{(2)}$; if $s>2$ then $F(L)$ ends in the final $0$ of the last appended block $0110$.
\end{proof}

\subsubsection{Odd parity forces survival}\label{b1:sec:odd}
\begin{lemma}[Odd-parity lemma]\label{b1:lem:odd}
If $\nu(L)$ is odd, then the counter computed during $F(L)$ satisfies $s\ge 2$.
\end{lemma}
\begin{proof}
Write $k=\nu(L)$ odd, $b\ge 0$ the number of trailing $0$'s of $L$, so $L=L_f\cdot 1\cdot 0^b$. Step 1 appends $0001$, producing $L^{(1)}=L_f\cdot 1\cdot 0^{b+3}\cdot 1$. Now $\nu(L^{(1)})=k+1$ is even, and the final two $1$'s of $L^{(1)}$ are separated by exactly $b+3$ zeros, so the last pair contributes $b+3\ge 3$. Every pair contributes $\ge 0$ and $s$ began at $-1$, whence $s\ge -1+(b+3)\ge 2$.
\end{proof}

From here on only even-parity strings can halt. With $s$ starting at $0$,
\begin{equation}\label{b1:eq:sumeven}
s=\sum_{j=1}^{m}(p_{2j}-p_{2j-1}-1)\qquad(\nu(L)\text{ even}),
\end{equation}
a sum of nonnegative within-pair gaps. Nothing forbids $s\in\{0,1\}$; locating those occurrences is the next task.

\subsubsection{The run-length calculus}\label{b1:sec:rl}
Write the run-length encoding of $L$ (with $\nu(L)\ge 1$) as
\begin{equation}\label{b1:eq:rle}
L=0^{c_0}1^{r_1}0^{c_1}1^{r_2}0^{c_2}\cdots 1^{r_R}0^{c_R},
\end{equation}
with $R\ge 1$, each $r_i\ge 1$, each internal $c_i\ge 1$. Put $C_i=r_1+\cdots+r_i$, so $C_0=0$ and $C_R=\nu(L)$.

\begin{proposition}[Counter formula]\label{b1:prop:counter}
If $\nu(L)$ is even, the counter computed during $F(L)$ is
\begin{equation}\label{b1:eq:counter}
s=\sum_{\substack{1\le i\le R-1\\C_i\ \text{odd}}}c_i.
\end{equation}
\end{proposition}
\begin{proof}
Step 1 adds no $1$'s, so \eqref{b1:eq:sumeven} applies. Number the $1$'s $1,\dots,N$; Step 2 pairs them $\{1,2\},\{3,4\},\dots$, and each pair contributes the number of $0$'s between its two $1$'s. The boundary between run $i$ and run $i+1$ is spanned by a pair iff $C_i$ is odd, contributing $c_i$; same-run pairs contribute $0$. Summing gives~\eqref{b1:eq:counter}.
\end{proof}

\begin{corollary}[Halting dichotomy]\label{b1:cor:halting}
Let $\nu(L)$ be even. Then $F(L)$ halts iff exactly one of: (i) every run length is even ($s=0$, equivalently $\omega(L)=0$); or (ii) the run lengths contain exactly two odd values at consecutive positions $j_0,j_0+1$, with $c_{j_0}=1$ ($s=1$).
\end{corollary}
\begin{proof}
By~\eqref{b1:eq:counter} and $c_i\ge 1$ internally, $s$ counts internal indices with $C_i$ odd, weighted by $c_i$. The parity sequence $C_0,\dots,C_R$ starts and ends even and toggles at odd-length runs; the two stated profiles are exactly $s=0$ and $s=1$.
\end{proof}

Combined with Lemma~\ref{b1:lem:odd}, the dichotomy reduces Theorem~\ref{b1:thm:central} to a statement about shape.

\subsubsection{Structural lemmas}\label{b1:sec:struct}
Fix $L$, let $\Lambda=L\cdot(00\text{ or }0001)$, and let the consecutive pairs of $1$'s in $\Lambda$ have within-pair gaps $g_j\ge 0$ for $j=1,\dots,m=\tfrac12\nu(\Lambda)$. Write $C_i,\gamma_i$ for the cumulative $1$-counts and internal $0$-runs of $\Lambda$, and set $X(L)=\#\{i:C_i\text{ odd and }\gamma_i\text{ odd}\}$.

\begin{lemma}[Reshaping]\label{b1:lem:reshape}
The image $L^{(2)}$ has exactly $m$ runs of $1$'s, the $j$-th of length $g_j+1$; hence $F(L)$ has $1$-runs $g_1+1,\dots,g_m+1$ followed by $s-2$ runs of length $2$. A body run is odd iff its pair has even gap; the appended tail contributes only even-length runs.
\end{lemma}
\begin{proof}
Pair $j$ sets $p_{2j-1},\dots,p_{2j}-1$ to $1$ (a block of length $g_j+1$) and zeros $p_{2j}$; distinct blocks are separated by the zeroed position. Each factor $0110$ contributes one run $11$ of length $2$. Finally $g_j+1$ is odd iff $g_j$ is even.
\end{proof}

\begin{lemma}[Reshaping count]\label{b1:lem:reshapecount}
For every $L$, $\omega(F(L))=m-X(L)$.
\end{lemma}

\begin{lemma}[Odd-run floor from intra-run pairs]\label{b1:lem:oddrunfloor}
With $\mu$ the number of pairs both of whose $1$'s lie in a single $1$-run of $\Lambda$ and $R'$ the number of $1$-runs of $\Lambda$, $F(L)$ has at least $\mu$ odd-length $1$-runs, and $\mu\ge m-(R'-1)$.
\end{lemma}

\begin{lemma}[The tail manufactures odd runs]\label{b1:lem:tailodd}
If $L=W\cdot(0110)^u$ with $u\ge 1$ and $W$ ending in $0$, then $\omega(F(L))\ge u-\varepsilon$, where $\varepsilon=0$ if $\nu(L)$ even and $\varepsilon=1$ if $\nu(L)$ odd.
\end{lemma}

\begin{lemma}[Counter floor]\label{b1:lem:cfloor}
If $\nu(L)$ is even, then $s(F(L))\ge\lceil\omega(L)/2\rceil$.
\end{lemma}

\begin{lemma}[Growth]\label{b1:lem:growth}
If $L=W\cdot(0110)^u$ with $u\ge 1$, $W$ ending in $0$, and $\nu(L)$ odd, then $s(F(L))\ge 2u+2$.
\end{lemma}

Proofs of Lemmas~\ref{b1:lem:reshapecount}--\ref{b1:lem:growth} are the pairing analysis of \S\ref{b1:sec:rl} applied run by run; they are routine and we omit them, having verified each against the engines of Appendix~\ref{b1:app:alg}.

\subsubsection{Prefix dynamics and the all-even profile}\label{b1:sec:prefix}
We now address the first fragile profile, $s_n=0$. By Corollary~\ref{b1:cor:halting} this is exactly $\omega(L_n)=0$. We exclude it unconditionally at even indices and reduce it to a weak counter bound at odd indices.

For a string $L$ with $\nu(L)\ge 1$, let $r_1(L)$ denote the length of its first maximal $1$-run, and (when $L$ has at least two $1$-runs) $c_1(L)$ the length of its first internal $0$-run.

\begin{lemma}[Prefix alternation]\label{b1:lem:prefix}
For all $n\ge 0$,
\[ r_1(L_{n+1})=\begin{cases}1,& r_1(L_n)\ge 2,\\ c_1(L_n)+1\ (\ge 2),& r_1(L_n)=1.\end{cases} \]
Since $r_1(L_0)=1$, it follows that $r_1(L_n)=1\iff n$ even.
\end{lemma}
\begin{proof}
By Lemma~\ref{b1:lem:reshape}, the first $1$-run of $L_{n+1}=F(L_n)$ has length $g_1+1$, where $g_1$ is the within-pair gap of the first pair of $\Lambda_n=L_n\cdot(\mathrm{pad})$. The pad is appended at the right, so the first $1$-run of $\Lambda_n$ equals $r_1(L_n)$. If $r_1(L_n)\ge 2$ the first two $1$'s are adjacent, $g_1=0$ and $r_1(L_{n+1})=1$. If $r_1(L_n)=1$ and $L_n$ has at least two $1$-runs, the first pair straddles the first boundary, giving $g_1=c_1(L_n)$ and $r_1(L_{n+1})=c_1(L_n)+1\ge 2$. (If $L_n$ has a single $1$-run, then $\nu(L_n)=1$ is odd; Step 1 appends $0001$, creating a second $1$-run at gap $\ge 3$, so again $r_1(L_{n+1})\ge 2$.) With $a_n:=[\,r_1(L_n)\ge 2\,]$, $a_{n+1}=\lnot a_n$ and $r_1(L_0)=1$ give the closed form.
\end{proof}

\begin{proposition}[All-even profile excluded]\label{b1:prop:alleven}
Recall $s_n=0\iff\omega(L_n)=0$.
\begin{enumerate}[label=(\alph*)]
\item For every even $n$, $L_n$ begins with a $1$-run of length $1$; hence $\omega(L_n)\ge 1$ and $s_n\ne 0$. This is unconditional.
\item For every odd $n$: if $\nu(L_n)$ is odd then $s_n\ge 2$ by Lemma~\ref{b1:lem:odd}; if $\nu(L_n)$ is even and $s_{n-2}\ge 4$, then $\omega(L_n)\ge 1$, so $s_n\ne 0$.
\end{enumerate}
\end{proposition}
\begin{proof}
(a) By Lemma~\ref{b1:lem:prefix}, $r_1(L_n)=1$ for even $n$, so $\omega(L_n)\ge 1$ and $s_n\ne 0$ by Corollary~\ref{b1:cor:halting}. (b) The odd-$\nu$ sub-case is Lemma~\ref{b1:lem:odd}. For the even-$\nu$ sub-case, $L_{n-1}=L^{(2)}_{n-2}\cdot(0110)^u$ with $u=s_{n-2}-2$ and $L^{(2)}_{n-2}$ ending in $0$. If $s_{n-2}\ge 4$ then $u\ge 2$, and Lemma~\ref{b1:lem:tailodd} gives $\omega(L_n)\ge u-1\ge 1$.
\end{proof}

\begin{remark}\label{b1:rem:rem17}
The asymmetry is genuine: at an odd index with $\nu(L_n)$ even, the leading run can be even (Lemma~\ref{b1:lem:prefix} gives $r_1=c_1+1$, of unconstrained parity), so part~(a)'s direct argument does not transfer, and one falls back on the tail count of part~(b). On the actual trajectory $s_{n-2}\ge 4$ holds at every relevant step (Table~\ref{b1:tab:traj}); the only index with $s_{n-2}<4$ is $n-2=2$ (where $s_2=2$), and $n=4$ has $\nu(L_4)$ odd, so the hypothesis of part~(b) is never invoked at a vulnerable step. Thus on the trajectory the all-even profile is excluded throughout.
\end{remark}

\begin{remark}[The all-even profile is reachable as an output]\label{b1:rem:rem18}
One might hope to prove the impossibility of $s_n=0$ through the backward range analysis, by arguing that every output of $F$ carries a $0$ strictly between its first and last $1$, so that no iterate is a single $1$-run $1^a0^b$ and the simplest all-even shape would be unreachable. \emph{This is not so.} The odd pad $0001$ places the (subsequently zeroed) last $1$ at the very end of $L^{(1)}$, leaving no internal $0$ in $L^{(2)}$; explicitly,
\[ F(1)=11110=1^40\quad(\text{an output of }F,\ \omega=0),\qquad F(11110)=\text{halt with }s=0. \]
Thus the all-even profile \emph{is} reachable as an output of $F$, and the exclusion of $s=0$ is a property of particular orbits, not of the range of $F$. The exclusion on the orbit of $10$ rests on Proposition~\ref{b1:prop:alleven}, which uses only the prefix invariant (Lemma~\ref{b1:lem:prefix}), together with the finite backward decision procedure of \S\ref{b1:sec:backward}, which confirms independently that all $4565$ halting strings of length $\le 14$ are absent from the orbit of $10$; neither establishes unreachability of all-even profiles in general. (A second, formally independent certificate of the orbit-of-$10$ all-even exclusion is given in the companion seed-independent treatment by encoding the process as a $2$-state, $5$-symbol Turing machine and citing a Finite Automata Reduction non-halting certificate from the Busy Beaver Challenge; we do not reproduce it here.)
\end{remark}

For the remainder, ``halting'' may be read as ``$s_n=1$.''

\begin{proposition}[Reduction to a single profile]\label{b1:prop:single}
The trajectory halts iff some $s_n=1$.
\end{proposition}
\begin{proof}
Halting at step $n$ means $s_n\in\{0,1\}$; Proposition~\ref{b1:prop:alleven} excludes $s_n=0$ on the trajectory.
\end{proof}

\subsubsection{The reduction to a scalar inequality}\label{b1:sec:reduction}
For $n\ge 1$ the iterate has the form $L_{n-1}=L^{(2)}_{n-1}\cdot(0110)^u$ with $u=s_{n-2}-2$ and $L^{(2)}_{n-1}$ ending in $0$ (Lemma~\ref{b1:lem:reshape}).

\begin{proposition}[Reduction]\label{b1:prop:firsthalt}
If the process halts and $N$ is the first index at which it does, then $s_{N-2}\le 5$.
\end{proposition}
\begin{proof}
Since $N$ is the first halt, $s_{N-2}\ge 2$, so $u=s_{N-2}-2\ge 0$. If $u=0$ then $s_{N-2}=2\le 5$. If $u\ge 1$, apply Lemma~\ref{b1:lem:tailodd} to $L_{N-1}=L^{(2)}_{N-1}\cdot(0110)^u$: $\omega(L_N)\ge u-1=s_{N-2}-3$. By Proposition~\ref{b1:prop:single} the halt is $s_N=1$, so $\omega(L_N)=2$; hence $s_{N-2}-3\le 2$.
\end{proof}

\begin{corollary}[The remaining content of the theorem]\label{b1:cor:remaining}
The process never halts provided $s_n\ge 6$ for all $n\ge 3$; a fortiori it never halts if the two-step monotonicity $s_n\ge s_{n-2}$ holds for all $n\ge 3$.
\end{corollary}
\begin{proof}
If first halt at $N$, Proposition~\ref{b1:prop:firsthalt} gives $N-2\le 2$, i.e.\ $N\le 4$; but $s_0,\dots,s_4=3,5,2,8,31$ are all $\ge 2$. Monotonicity propagates $s_0,s_1,s_2$ on the two parity subsequences.
\end{proof}

Thus Theorem~\ref{b1:thm:central} is equivalent to the scalar statement that the counter never returns to $\{2,3,4,5\}$ after the third step, for which two-step monotonicity is sufficient. The mechanism is explicit. At an odd-parity step, Lemma~\ref{b1:lem:growth} gives
\begin{equation}\label{b1:eq:double}
s_n\ge 2s_{n-1}-2\qquad(\nu(L_n)\text{ odd}),
\end{equation}
so odd steps at least double the counter. At an even-parity step, Lemmas~\ref{b1:lem:cfloor} and~\ref{b1:lem:tailodd} give
\begin{equation}\label{b1:eq:half}
s_n\ge\lceil\omega(L_n)/2\rceil\ge\lceil(s_{n-2}-3)/2\rceil\qquad(\nu(L_n)\text{ even}),
\end{equation}
at most a near-halving. Estimate~\eqref{b1:eq:half} alone does not yield $s_n\ge s_{n-2}$; a long run of consecutive even-parity steps could compound several near-halvings. The next conjecture asserts this never happens.

\begin{conjecture}[Two-step monotonicity; open]\label{b1:conj:mono}
For all $n\ge 3$, $s_n\ge s_{n-2}$.
\end{conjecture}

\subsubsection{An improved even-step floor and a proven sub-case}\label{b1:sec:7.1}
We beat the factor $\tfrac12$ of~\eqref{b1:eq:half} in a structurally identified case. The key is that when the predecessor step is of even parity, the tail it carries is reshaped into a regular comb of single $1$'s separated by $0$-gaps of length $3$.

\begin{thm}[Even-step floor with even-parity predecessor]\label{b1:thm:floor}
Suppose step $n$ is of even parity, step $n-1$ is also of even parity, and $u:=s_{n-2}-2\ge 1$. Then the last $u$ body $1$-runs of $L_n$ all have length $1$ and are separated by internal $0$-runs of length exactly $3$, and consequently
\begin{equation}\label{b1:eq:floor}
s_n\ge 3\Bigl\lfloor\tfrac{u-1}{2}\Bigr\rfloor\ge\tfrac32 s_{n-2}-6.
\end{equation}
In particular $s_n\ge s_{n-2}$ whenever $s_{n-2}\ge 12$, so Conjecture~\ref{b1:conj:mono} holds at every even-parity step whose predecessor is of even parity, with finitely many small cases to check directly.
\end{thm}
\begin{proof}
Write $L_{n-1}=Y\cdot(0110)^u$ with $Y=L^{(2)}_{n-2}$ ending in $0$, $u=s_{n-2}-2$. The tail carries $2u$ ones (even), and $\nu(L_{n-1})$ is even, so $\nu(Y)$ is even. Forming $L_n$: Step 1 appends $00$, the $\nu(Y)$ ones of $Y$ pair among themselves, and the $2u$ tail ones pair within the tail as $(1,2),(3,4),\dots$.

The tail attached to $Y$ reads $0\underbrace{11\,00\,11\,00\cdots 00\,11}_{u\ \text{blocks}}0$. By Lemma~\ref{b1:lem:reshape} each block becomes a $1$-run of length $1$, Step 2 zeros its right one, and between two consecutive blocks the symbols are the zeroed right one plus the separating $00$---three $0$'s. So the image of the tail region is $(1\,000)^{u-1}\,1$: $u$ single-$1$ runs separated by $u-1$ internal $0$-runs of length $3$, lying strictly inside $L_n$.

Across these $u$ runs the cumulative count $C$ rises by $1$ per run, so the $u-1$ separators sit at $u-1$ consecutive values of $C$, at least $\lfloor(u-1)/2\rfloor$ of them odd, each contributing $c_i=3$ via~\eqref{b1:eq:counter}. Hence $s_n\ge 3\lfloor(u-1)/2\rfloor\ge\tfrac32(u-2)=\tfrac32 s_{n-2}-6$, which is $\ge s_{n-2}$ iff $s_{n-2}\ge 12$. The only even-after-even step with $s_{n-2}<12$ on the trajectory is $n=3$, where $s_3=8\ge s_1=5$ holds directly.
\end{proof}

\begin{remark}\label{b1:rem:floor-consistent}
The bound~\eqref{b1:eq:floor} is consistent with the data: at every even-after-even step ($n=3,11,12,22,23,24,25,26$) the observed ratios $s_n/s_{n-2}$ range from $1.60$ to $7.33$. The comb structure (last $u=s_{n-2}-2$ body $1$-runs of length $1$, separators of length $3$) was confirmed exactly at $n=3,11,12$. The complementary case---an even step whose predecessor is of odd parity---is exactly where the tight instances of Conjecture~\ref{b1:conj:mono} occur; see \S\ref{b1:sec:12.1}.
\end{remark}

\subsubsection{The counter-inert tail and body localization}\label{b1:sec:inert}
The even-step floor of \S\ref{b1:sec:7.1} lower-bounds $s_n$ through $\omega(L_n)$. We now isolate a cleaner and exact statement which unifies the proved and the open cases. The principle is that the counter at an even step is a functional of the body alone: the freshly appended tail contributes nothing.

\begin{lemma}[Counter-inert tail]\label{b1:lem:inert}
Let step $n$ be of even parity, and write $L_n=B\cdot T$ with $T=(0110)^{s_{n-1}-2}$ the tail appended in forming $L_n=F(L_{n-1})$ and $B=L^{(2)}_{n-1}$ the body. Then $T$ contributes exactly $0$ to $s_n$, and $s_n$ equals the within-pair-gap sum over $B$ alone.
\end{lemma}
\begin{proof}
Step 1 appends $00$ and the $1$'s of $L_n$ pair consecutively from the left. The tail carries $2(s_{n-1}-2)$ ones, even, and $\nu(L_n)$ is even, so $\nu(B)$ is even; $B$-ones pair within $B$ and tail ones pair within the tail as blocks $11$ of gap $0$.
\end{proof}

Combining Lemma~\ref{b1:lem:inert} with~\eqref{b1:eq:counter} yields an exact decomposition of the even-step counter.

\begin{observation}[Body localization and the count/surplus split]\label{b1:obs:decomp}
At an even-parity step $n$, the nonzero within-pair gaps occupy exactly the first $\nu(L^{(2)}_{n-1})/2$ pair-slots---the body---and the trailing pairs sum to $0$. Each contributing body pair carries a within-pair gap $g\ge 1$; writing $g=1+(g-1)$ partitions the counter exactly as
\begin{equation}\label{b1:eq:decomp}
s_n = \underbrace{\#\{\text{body boundaries with odd cumulative $1$-count}\}}_{\text{count term}} + \underbrace{\sum_{\text{those boundaries}}(g-1)}_{\text{size surplus}}.
\end{equation}
The count term is bounded below by the floor $\lceil\omega(L_n)/2\rceil$ of Lemma~\ref{b1:lem:cfloor}, with equality only when each parity-raised interval of the cumulative-$1$-count contributes a single boundary---an extremal case that does not occur at tight instances of Conjecture~\ref{b1:conj:mono}. The size surplus is carried by the contributing within-pair gaps of size $\ge 2$.
\end{observation}

The decomposition~\eqref{b1:eq:decomp} was confirmed quantitatively at every odd-predecessor even step within reach (Appendix~\ref{b1:app:evidence}, Table~\ref{b1:tab:decomp}). In each case the freshly appended tail contributes $0$, in exact agreement with Lemma~\ref{b1:lem:inert}. At the binding case $n=15$ the contributing gaps are dominated by $1$'s---the signature of the cross-paired comb produced by an odd-parity predecessor---the count term is $23{,}273$ (strictly exceeding the lower bound $\lceil\omega(L_{15})/2\rceil=22{,}455$), and the surplus over that count, namely $5810=5613\cdot 1+ 1\cdot 2+ 65\cdot 3$, is supplied entirely by the gaps of size $\ge 2$, which originate, recursively, in the deeper body $L^{(2)}_{n-2}$. The true margin is $s_{15}-s_{13}=349$.

This is the precise sense in which the residual difficulty has been relocated: by~\eqref{b1:eq:decomp} the whole question is whether the size surplus is large enough; \S\ref{b1:sec:12.1} shows that the count term, even with its strongest provable lower bound $\lceil\omega/2\rceil$, can never reach it.

\subsubsection{The $\nu$-recurrence and the parity dynamics}\label{b1:sec:nu}
The parity of $\nu(L_n)$ decides which regime governs step $n$ and is the pivotal discrete datum of the whole development. It obeys an exact recurrence.

\begin{proposition}[$\nu$-recurrence]\label{b1:prop:nu}
Whenever $F(L)$ is defined, with $\nu=\nu(L)$ and $s$ the counter,
\[ \nu(F(L))=\begin{cases}\tfrac{\nu}{2}+3s-4,& \nu\text{ even},\\ \tfrac{\nu+1}{2}+3s-3,& \nu\text{ odd}.\end{cases} \]
Consequently the parity of $\nu(L_{n+1})$ is determined by $(\nu(L_n)\bmod 4,\,s_n\bmod 2)$; for even parity, $\nu(L_{n+1})\equiv\nu(L_n)/2+s_n\pmod 2$.
\end{proposition}
\begin{proof}
$\nu(F(L))=\nu(L^{(2)})+2(s-2)$. By Lemma~\ref{b1:lem:reshape}, $L^{(2)}$ has $m$ runs of total $1$-length $m+\sum_j g_j$. For even parity, $\sum_j g_j=s$ and $m=\nu/2$, so $\nu(L^{(2)})=\nu/2+s$; for odd parity, $\sum_j g_j=s+1$ and $m=(\nu+1)/2$, so $\nu(L^{(2)})=(\nu+1)/2+s+1$. Substitute.
\end{proof}

Iterating Proposition~\ref{b1:prop:nu} from $\nu(L_0)=1$ with the computed counters reproduces the observed parity string exactly (Appendix~\ref{b1:app:evidence}), and yields the parity at $n=26$ without materializing $L_{26}$. The counter $s_n$ itself admits no scalar recurrence: by Proposition~\ref{b1:prop:counter} it depends on the full pattern of odd-cumulative gaps of $L_n$.

\subsubsection{Equivalent formulations}\label{b1:sec:equiv}
We record three reformulations of the process. Throughout, $\popcount(x)$ denotes the popcount of $x$ and $V(L):=\mathrm{int}(L,2)$ the binary value.

\subsubsection{The pair-and-fill identity}\label{b1:sec:10.1}
For an integer $M$ with $\popcount(M)=2m$ and $1$-bit positions $b_1>b_2>\dots>b_{2m}\ge 0$, set
\[ S(M):=\sum_{i=1}^{2m}(-1)^{i+1}2^{b_i}. \]

\begin{lemma}[Step 2 is doubling of the alternating bit-sum]\label{b1:lem:S}
If $\nu(L)$ is even and $M=V(L^{(1)})$, then $V(L^{(2)})=2S(M)$. The pair $(b_{2j-1},b_{2j})$ effects the local change $2^{b_{2j-1}}-3\cdot 2^{b_{2j}}$ to $M$.
\end{lemma}
\begin{proof}
The left/right members of pair $j$ sit at heights $A_j=b_{2j-1}>B_j=b_{2j}$. Step 2 replaces $2^{A_j}+2^{B_j}$ by $\sum_{i=B_j+1}^{A_j}2^i=2^{A_j+1}-2^{B_j+1}$; the local change is $2^{A_j}-3\cdot 2^{B_j}$. Summing, $V(L^{(2)})=\sum_j(2^{A_j+1}-2^{B_j+1})=2\sum_i(-1)^{i+1}2^{b_i}=2S(M)$.
\end{proof}

The constant $3$ here is the genuine arithmetic kernel of the process, the structural counterpart of the $3$ in the Collatz map.

\subsubsection{The integer form}\label{b1:sec:10.2}
By Lemma~\ref{b1:lem:prefix} the leading $0$-run of every iterate is empty, so each $L_n$ begins with $1$; each ends in $0$ (Lemma~\ref{b1:lem:end0}). Hence $V$ is a bijection from the reachable strings onto $\mathbb Z^+$, with $V(L_0)=2$.

\begin{thm}[Integer conjugate]\label{b1:thm:T}
Through $V$, $F$ is conjugate to $T:\mathbb Z^+\dashrightarrow\mathbb Z^+$ given by
\[ T(N)=2S(M)\cdot 16^{s-2}+6\cdot\tfrac{16^{s-2}-1}{15},\qquad M=\begin{cases}4N,& \popcount(N)\text{ even},\\ 16N+1,& \popcount(N)\text{ odd},\end{cases} \]
\[ s=s_{\mathrm{init}}+\popcount(S(M))-\tfrac12\popcount(M),\qquad s_{\mathrm{init}}=\begin{cases}0,& \popcount(N)\text{ even},\\ -1,& \popcount(N)\text{ odd},\end{cases} \]
defined exactly when $s\ge 2$. In hexadecimal, $T(N)$ is the digit string of $2S(M)$ followed by $s-2$ copies of the digit $6$. The seed is $N_0=2$.
\end{thm}
\begin{proof}
Step 1 multiplies by $4$ (pad $00$) or $16$ then $+1$ (pad $0001$), giving $M$. Step 2 gives $V(L^{(2)})=2S(M)$ (Lemma~\ref{b1:lem:S}). The counter is $s=s_{\mathrm{init}}+\nu(L^{(2)})-m=s_{\mathrm{init}}+\popcount(S(M))-\tfrac12\popcount(M)$. Step 3, when $s\ge 2$, appends $(0110)^{s-2}$; since $0110_2=6$ fills one hex digit, this is multiplication by $16^{s-2}$ followed by adding $\underbrace{6\cdots 6}_{s-2}=6(16^{s-2}-1)/15$.
\end{proof}

\begin{remark}[How close to Collatz?]\label{b1:rem:collatz}
The shape is Collatz-like: a parity test on $\popcount(N)$ selects one of two affine pads, the kernel carries the constant $3$ (Lemma~\ref{b1:lem:S}), and the whole question reduces to one inequality on a derived scalar. Two features are not Collatz-like, and both are intrinsic. First, the orbit grows doubly exponentially,
\[ N_0,N_1,N_2,N_3,\dots=2,\,998,\,45\,868\,646,\,213\,192\,976,\dots, \]
because the binary value carries $\Theta(|L_n|)$ digits and $|L_n|\sim 4^{s_n}$ grows geometrically. Second, $S$ is non-local: it depends on the global ordering of all set bits rather than on $N$ through a fixed affine rule.
\end{remark}

\subsubsection{The gap-vector form}\label{b1:sec:10.3}
Encode a string that begins with $1$ and ends in $0$ by the gaps between its consecutive $1$'s together with its trailing-zero count,
\[ (e;t),\qquad e=(e_1,\dots,e_{N-1})\in\mathbb Z_{\ge 0}^{N-1},\ t\ge 1, \]
where $N=\nu(L)$. The seed is $\bigl((\,);1\bigr)$.

\begin{proposition}[Gap-vector conjugate]\label{b1:prop:gap}
Write $N=\nu$. One step of $F$ is:
\begin{itemize}
\item (\emph{pad; seed the counter}) If $N$ is even, set $s_{\mathrm{init}}:=0$ and leave $(e;t)$ unchanged. If $N$ is odd, set $s_{\mathrm{init}}:=-1$, append $e_N:=t+3$, reset $t:=0$, and increment $N$. Now $N$ is even.
\item (\emph{counter}) $s=s_{\mathrm{init}}+(e_1+e_3+e_5+\cdots+e_{N-1})$; halt if $s<2$.
\item (\emph{rewrite}) The new word of $1$-runs and separators is
\[ 1^{e_1+1}\,0^{1+e_2}\,1^{e_3+1}\,0^{1+e_4}\cdots 1^{e_{N-1}+1}\ \text{followed by}\ (0110)^{s-2}. \]
The new trailing-zero count is $t'=1$ when $s>2$, and $t'=t+3$ when $s=2$.
\end{itemize}
\end{proposition}
\begin{proof}
This is the run-length calculus of \S\ref{b1:sec:rl} in gap coordinates. Pair $j$ has within-pair gap $e_{2j-1}$ and yields a block of $e_{2j-1}+1$ ones; its zeroed right endpoint together with the between-pair gap $e_{2j}$ makes the separator $1+e_{2j}$. The counter is $s_{\mathrm{init}}+\sum_j e_{2j-1}$. The odd-$N$ case is the $0001$ pad. When $s>2$ the appended factor ends in $0$, $t'=1$; when $s=2$ the word ends at the zeroed right endpoint of the last pair, after which lie the $t+2$ trailing zeros of the padded string (two from the $00$ pad), giving $t'=t+3$.
\end{proof}

\begin{remark}[Two subtleties of the encoding]\label{b1:rem:enc-subtleties}
Two points are easy to get wrong and worth isolating. First, the additive seed $s_{\mathrm{init}}=-1$ of the odd-parity pad must be carried in the counter line; without it the seed step returns $s=4$ in place of $s_0=3$. Second, the trailing-zero count is not always $1$: the unique $s=2$ step ($n=2$) produces $L_3$ with four trailing zeros ($t=1$, so $t'=t+3=4$), so $t$ must be tracked honestly. With both handled as stated, the gap-vector engine reproduces the entire trajectory exactly.
\end{remark}

\subsubsection{The scalar form}\label{b1:sec:10.4}
Collapsing to the counter alone: by Proposition~\ref{b1:prop:alleven} the all-even profile never occurs on the trajectory, and then (Proposition~\ref{b1:prop:single}) a halt is exactly some $s_n=1$. With Corollary~\ref{b1:cor:remaining} this gives the following.

\begin{proposition}[Scalar form]\label{b1:prop:scalar}
On the trajectory, $F$ is non-halting if and only if the integer sequence $(s_n)_{n\ge 0}$ never re-enters $\{2,3,4,5\}$ after $n=2$; a clean sufficient condition is two-step monotonicity $s_n\ge s_{n-2}$ for $n\ge 3$ (Conjecture~\ref{b1:conj:mono}). The sequence satisfies $s_n\ge 2s_{n-1}-2$ at popcount-odd steps and $s_n\ge\lceil(s_{n-2}-3)/2\rceil$ at popcount-even steps, but admits no closed recurrence.
\end{proposition}

This last point is the precise sense in which the problem is of Collatz type, and it is the wall on which both the forward inequality and the backward search (\S\ref{b1:sec:backward}) rest.

\subsubsection{The backward viewpoint}\label{b1:sec:backward}
A complementary strategy reasons from a hypothetical halt. Its foundation is the finiteness noted after Lemma~\ref{b1:lem:length}: preimages are strictly shorter, so the backward tree of any string is finite and orbit membership is decidable. In fact $F$ is exactly invertible.

\begin{proposition}[Invertibility]\label{b1:prop:inv}
Let $T\in\{0,1\}^*$. Every $X$ with $F(X)=T$ arises as follows: choose $u\ge 0$ with $T=Y\cdot(0110)^u$ and $Y$ ending in $0$; from the maximal $1$-blocks $[a,b]$ of $Y$ form the candidate $L^{(1)}$ with $1$'s exactly at $\{a,b+1\}$ over the length of $Y$; delete a suffix $00$ or $0001$ to obtain $X$; and retain $X$ iff $F(X)=T$. There are finitely many preimages.
\end{proposition}
\begin{proof}
If $F(X)=T$ then $T=L^{(2)}\cdot(0110)^{s-2}$ with $L^{(2)}$ ending in $0$, so $T$ has the stated form with $u=s-2$ and $Y=L^{(2)}$. Step 2 turns each pair $(p_{2j-1},p_{2j})$ into the block $[p_{2j-1},p_{2j}-1]$ and zeros $p_{2j}$, so the blocks $[a,b]$ of $Y$ recover the pair endpoints, determining $L^{(1)}$ uniquely. Verifying $F(X)=T$ discards spurious reconstructions. Finiteness is Lemma~\ref{b1:lem:length}.
\end{proof}

Proposition~\ref{b1:prop:inv} turns the backward search into a finite, exact computation. Suppose the process first halts at index $N\ge 2$. By Proposition~\ref{b1:prop:firsthalt}, $s_{N-2}\le 5$. We now ask what is forced about the predecessor counter $s_{N-1}$ and the shape of $L_N$.

\begin{proposition}[Shape of a first halt]\label{b1:prop:halt-corr}
If the process first halts at $N\ge 2$, then $s_{N-2}\le 5$ (Proposition~\ref{b1:prop:firsthalt}), and by Lemma~\ref{b1:lem:inert} the halt depends only on the body $B=L^{(2)}_{N-1}$: the tail $(0110)^{s_{N-1}-2}$ of $L_N$ contributes $0$ to $s_N$. Consequently the predecessor counter $s_{N-1}$ is not determined by the halt and may be arbitrarily large.
\end{proposition}
\begin{proof}
Proposition~\ref{b1:prop:firsthalt} gives $s_{N-2}\le 5$. At the halt $\nu(L_N)$ is even, so Lemma~\ref{b1:lem:inert} applies to the step computing $s_N$, which therefore depends on $B$ alone, independently of the tail length $s_{N-1}-2$.
\end{proof}

\begin{remark}[The predecessor counter is not pinned at a first halt]\label{b1:rem:rem-pinning}
It is tempting to think that a first halt forces $s_{N-1}=2$, so that the halting string has an empty tail $L_N=L^{(2)}_{N-1}$ and consists of $m$ image blocks summing to $m+2$ ones. \emph{This is not so}: the pinning $s_{N-1}=2$ would contradict Lemma~\ref{b1:lem:inert}, which makes the halt independent of $s_{N-1}$. Two explicit counterexamples:
\begin{itemize}
\item the seed $1100111$ first halts at $N=2$ with $s_{N-1}=3$, $s_N=1$;
\item the seed $110011111110$ first halts at $N=8$ with $s_{N-1}=102$, $s_N=1$.
\end{itemize}
The subtlety is that the tail's $11$ blocks are even runs of $L_N$, not odd runs, so Lemma~\ref{b1:lem:cfloor} yields only $\lceil\omega/2\rceil=1$ rather than $s_N\ge 2$. The halting strings are therefore $B\cdot(0110)^k$ ($k\ge 0$) with $B$ fragile.
\end{remark}

The finite backward decision procedure is exact regardless. By Lemma~\ref{b1:lem:length} the iterated preimages of any fixed string form a finite tree, so ``does $w$ ever occur on the trajectory?'' is decided by tracing that tree to its roots and testing membership of $L_0=10$. Running this decision procedure wholesale, every halting string of length at most $14$---there are $4565$ of them, of which $3137$ have three or more $1$-runs---is found to be absent from the trajectory. The procedure does not extend to a closed-form argument for unbounded lengths: backward chains lengthen with size, so the backward route demands a uniform descent of strength equal to the two-step monotonicity $s_n\ge s_{n-2}$. As \S\ref{b1:sec:resist} makes precise, this is the same wall the forward route meets.

\subsubsection{Computational results}\label{b1:sec:comp}
All computations were carried out by two independent implementations of Definition~\ref{b1:def:F}, described in Appendix~\ref{b1:app:alg}: a direct bit-level simulator and a run-length engine acting on the gap sequence of the $1$'s. The two agree on every output where both are feasible---full strings through $n=12$, and the scalar quantities $s_n$ and $\nu(L_n)\bmod 2$ through $n=25$---and the run-length engine was further cross-checked against the $\nu$-recurrence (Proposition~\ref{b1:prop:nu}). Every figure reported here was reproduced by at least two independent computations; the external verification is described in Appendix~\ref{b1:app:ai}.

\paragraph{Counter and parity.} Table~\ref{b1:tab:traj} lists $s_n$ and the parity of $\nu(L_n)$ for $0\le n\le 26$. Every $s_n\ge 2$; the unique near miss is $s_2=2$.

\paragraph{Two-step monotonicity.} The inequality $s_n\ge s_{n-2}$ holds for every $3\le n\le 26$. The tightest instance is the odd$\to$even transition $n=15$, where $s_{15}/s_{13}=29083/28734\approx 1.0121$.

\paragraph{A long even-parity run.} The parity of $\nu(L_n)$ is even at $n=2,3,10,11,12,15,19,21,22,23,24,25,26$. In particular $n=21,\dots,26$ are six consecutive even-parity steps---precisely the configuration whose danger~\eqref{b1:eq:half} cannot a priori rule out. Yet across this run the counter strictly increases at every step,
\[ s_{21..26}=3\,732\,091,\,8\,693\,523,\,18\,136\,691,\,24\,703\,847,\,55\,831\,137,\,88\,613\,778, \]
so the feared compounding of near-halvings does not occur. As $n=26$ is even-after-even, it falls in the regime already settled by Theorem~\ref{b1:thm:floor}.

\subsubsection{The residual obstruction, localized}\label{b1:sec:12.1}
The even-parity steps split by the parity of their predecessor:
\begin{center}
\begin{tabular}{ll}
predecessor even (Theorem~\ref{b1:thm:floor} applies) & $n=3,11,12,22,23,24,25,26$\,; ratios $1.60$--$7.33$\\
predecessor odd (open) & $n=10,15,19,21$\,; ratios $1.36,1.01,1.32,5.66$
\end{tabular}
\end{center}
Every tight instance of Conjecture~\ref{b1:conj:mono}---$n=15$ ($1.012$), $n=19$ ($1.324$), $n=10$ ($1.361$)---has an odd-parity predecessor. The structural reason is that an odd-parity predecessor forces cross-pairing of its tail (Lemma~\ref{b1:lem:tailodd}, $\varepsilon=1$), producing odd runs separated by small gaps rather than the size-$3$ comb of Theorem~\ref{b1:thm:floor}; at $n=15$ the counted gaps average $1.25$, so the floor collapses to the factor $\tfrac12$ of~\eqref{b1:eq:half}.

By the decomposition~\eqref{b1:eq:decomp} of Observation~\ref{b1:obs:decomp}, the entire residual difficulty is the size-surplus term. At the binding case the strongest provable lower bound on the count term---the floor $\lceil\omega(L_n)/2\rceil$ of Lemma~\ref{b1:lem:cfloor}---is already insufficient: directly,
\[ \lceil\omega(L_{15})/2\rceil=22\,455<s_{13}=28\,734,\qquad \omega(L_{15})=44\,910<2\,s_{13}, \]
so any argument of the form ``$s_n\ge\tfrac12\omega(L_n)$ with $\omega(L_n)\ge 2s_{n-2}$'' is dead on arrival. Concretely, the contributing-gap multiset at $n=15$ is $\{1:17\,594,\,2:5613,\,3:1,\,4:65\}$ (Table~\ref{b1:tab:decomp}): the count term contributes the $23\,273$ boundaries---strictly exceeding the lower bound $22\,455$---and the surplus over that count, namely $5613\cdot 1+1\cdot 2+65\cdot 3=5810$, is supplied entirely by the gaps of size $\ge 2$, which originate, recursively, in the deeper body $L^{(2)}_{n-2}$. The true margin is $s_{15}-s_{13}=349$.

This is the entire residual content of Conjecture~\ref{b1:conj:mono}:
\[ \text{prove }s_n\ge s_{n-2}\text{ at an even-parity step $n$ whose predecessor is of odd parity,} \]
a statement requiring control of the alignment-filtered sizes of the cross-paired gaps, not merely of odd-run counts. The surplus is a sum of the predecessor's within-pair gaps restricted to the boundaries at which the running $1$-count is odd; that restriction is an alignment depending on the parities of all earlier run lengths, hence on the full history---by Proposition~\ref{b1:prop:scalar}, on a quantity with no scalar recurrence. No descent or Lyapunov function in the computable data ($s_n$, $s_n+s_{n-1}$, ratios, $\omega$, gap-moment sums) was found that is preserved across the $s_2=2$ near miss while forcing $s_n\ge 6$ thereafter. The inequality thus appears to lie beyond elementary methods; \S\ref{b1:sec:resist} develops why, and \S\ref{b1:sec:status} states the status precisely.

\subsubsection{Why the problem resists}\label{b1:sec:resist}
It is worth setting out, concretely, why a problem whose answer is in so little doubt should be so hard to prove. The obstruction is not a shortage of structure---the process is unusually transparent---but a specific mismatch between the quantity one must control and the quantities one can compute. We organise the discussion around the two natural proof strategies, which turn out to fail at the same point.

\subsubsection{Two routes, one obstruction}\label{b1:sec:13.1}
The forward route sweeps along the trajectory and asks how a single application of $F$ reshapes runs. By the reshaping identities (Lemmas~\ref{b1:lem:reshape}--\ref{b1:lem:oddrunfloor}) the odd-length runs of $L_{n+1}$ are exactly the images of the even-gap pairs of $L_n$, the appended tail contributes only even runs, and intra-run pairing manufactures odd runs in profusion. To forbid both fragile profiles at once it suffices to know that every even-parity iterate carries at least four odd-length runs (four, not two, because the count is necessarily even). Call this \emph{Invariant~(A)}. It instantly excludes $\omega=0$ and $\omega=2$, and with Lemma~\ref{b1:lem:odd} it would finish Theorem~\ref{b1:thm:central}.

The difficulty is that the only uniform lower bound available, $\mu\ge m-(R'-1)$ of Lemma~\ref{b1:lem:oddrunfloor}, collapses precisely when the string degenerates to a near-comb of length-one runs: for $L=1010\cdots 10$ one has $\nu(\Lambda)\approx R'$ and the bound gives no floor at all. Invariant~(A) therefore cannot be read off a single step; it is a statement about the orbit never degenerating to such a comb at an even step.

The backward route assumes a first halt at $N$ and analyses $L_N$. By Proposition~\ref{b1:prop:halt-corr}, $s_{N-2}\le 5$ and the halt is a functional of the body $B=L^{(2)}_{N-1}$ alone---the tail of $L_N$ is counter-inert. The analysis of Remark~\ref{b1:rem:rem-pinning} allows nonempty tails: the halting strings are $B\cdot(0110)^k$ with $B$ fragile, and the finite preimage decision procedure (Proposition~\ref{b1:prop:inv}) traces the backward tree of any $B$ to its roots and tests membership of $L_0$. This procedure dispatches every halting string of length $\le 14$, but it does not extend to a closed-form argument: backward chains lengthen with $|B|$, and what is required to rule out fragile $B$ uniformly is precisely Invariant~(A)---that no $L_n$ ever lands in a fragile shape. The forward route calls it Invariant~(A); the backward route needs it to retire the unbounded family of fragile bodies $B$. The two decisive steps coincide.

\subsubsection{Why that obstruction is a wall}\label{b1:sec:13.2}
Three features, each intrinsic, conspire to keep Invariant~(A) out of reach.

\emph{The counter is a parity functional, small only on fragile shapes.} By Proposition~\ref{b1:prop:counter}, $s$ adds up the internal gaps that sit after an odd cumulative count of $1$'s. This is a global parity selection, not a local or affine readout: which gaps are counted depends on the parities of all preceding run lengths. A parity functional can be small only on configurations of a very particular shape, so the entire content of the theorem is that the dynamics steer the orbit clear of those shapes---and no amount of cleverness about a single step controls a global parity alignment.

\emph{The decisive quantity is a size surplus, not a count.} The exact decomposition~\eqref{b1:eq:decomp} splits the even-step counter into a count term, bounded below by the floor $\lceil\omega/2\rceil$, and a size surplus carried by the within-pair gaps of size $\ge 2$ at odd boundaries. At the binding case $n=15$ the count floor $\lceil\omega(L_{15})/2\rceil=22\,455$ is already smaller than $s_{13}=28\,734$, so no argument bounding odd-run counts from below by $\lceil\omega/2\rceil$ can close the gap (\S\ref{b1:sec:12.1}); the whole margin comes from the sizes of a few gaps, and those sizes are themselves alignment-filtered, depending on the full history. One is asked to control not how many odd runs there are, but how large the right-aligned ones are---a strictly finer datum.

\emph{There is no scalar recurrence, and the single near miss forbids crude monotonicity.} By Proposition~\ref{b1:prop:scalar} the counter obeys no closed recurrence: $s_{n+1}$ depends on the entire odd-cumulative gap pattern of $L_n$, not on finitely many earlier counters. This is the precise sense in which the system is of Collatz type, and through the integer conjugate $T$ (Theorem~\ref{b1:thm:T}) the kinship is exact---a popcount-parity test selecting one of two affine pads, with the kernel carrying the Collatz constant $3$ (Lemma~\ref{b1:lem:S}). Worse for any monotone approach, the trajectory touches the threshold with equality exactly once, at $s_2=2$. A successful argument must therefore permit the counter to descend all the way to $2$ once and then explain---structurally---why it never descends again.

\emph{Growth makes finite verification powerless as a substitute.} The integer orbit grows doubly exponentially (Remark~\ref{b1:rem:collatz}), because the binary value carries $\Theta(|L_n|)$ digits and $|L_n|\sim 4^{s_n}$ grows geometrically; the run count itself roughly triples per step. No finite computation, however far pushed, can stand in for the missing uniform argument.

\subsubsection{What the evidence does and does not settle}\label{b1:sec:13.3}
The reductions are complete and rigorous: Theorem~\ref{b1:thm:central} is equivalent, modulo the proved Lemmas~\ref{b1:lem:odd},~\ref{b1:lem:reshape} and Proposition~\ref{b1:prop:counter}, to the single statement that no even-parity iterate wears a fragile profile, for which two-step monotonicity is a clean sufficient condition. That statement holds for every iterate within computational reach---the even members carry $4,6,1420,3704,1960,44910,\dots$ odd runs, never the $0$ or $2$ a halt would require---and for every halting string of length at most $14$ by the backward decision procedure. The counter's geometric recovery after its lone near miss, the strict growth even across six consecutive even-parity steps, two independent simulators in agreement, and the external verification of Appendix~\ref{b1:app:ai} all point the same way. What remains is a single, sharply isolated inequality whose proof would require a size-aware estimate of the cross-paired gap structure that the present methods do not supply.

\subsubsection{Status}\label{b1:sec:status}
\subsubsection*{Proved} $F$ strictly lengthens and outputs end in $0$ (Lemmas~\ref{b1:lem:length},~\ref{b1:lem:end0}); odd parity forces $s\ge 2$ (Lemma~\ref{b1:lem:odd}); for even parity the counter is the odd-cumulative gap sum~\eqref{b1:eq:counter}, whence the halting dichotomy (Corollary~\ref{b1:cor:halting}); the reshaping identities (Lemmas~\ref{b1:lem:reshape}--\ref{b1:lem:growth}); the prefix-alternation invariant $r_1(L_n)=1\iff n$ even (Lemma~\ref{b1:lem:prefix}); the exclusion of the all-even profile, unconditionally at even indices and---via $s_{n-2}\ge 4$---at odd indices (Proposition~\ref{b1:prop:alleven}); the Reduction, that a first halt forces $s_{N-2}\le 5$ (Proposition~\ref{b1:prop:firsthalt}); the $\nu$-recurrence (Proposition~\ref{b1:prop:nu}); the counter-inert tail and the exact count/surplus decomposition (Lemma~\ref{b1:lem:inert}, Observation~\ref{b1:obs:decomp}); the even-step floor giving two-step monotonicity at every even step with an even-parity predecessor (Theorem~\ref{b1:thm:floor}); invertibility of $F$ and the first-halt shape (Propositions~\ref{b1:prop:inv},~\ref{b1:prop:halt-corr}).

\subsubsection*{Reduced} Theorem~\ref{b1:thm:central} is equivalent (Corollary~\ref{b1:cor:remaining}) to the scalar statement that $(s_n)$ never re-enters $\{2,3,4,5\}$ after $n=2$, for which two-step monotonicity $s_n\ge s_{n-2}$ is sufficient. Monotonicity itself is now reduced to a single configuration: an even-parity step with an odd-parity predecessor (\S\ref{b1:sec:12.1}). The odd-parity steps ($s_n\ge 2s_{n-1}-2$) and the even-after-even steps (Theorem~\ref{b1:thm:floor}) are settled. Within the open configuration, the difficulty is localized further by~\eqref{b1:eq:decomp} onto the alignment-filtered size surplus, the count term being provably insufficient under its strongest available lower bound.

\subsubsection*{Open} The inequality $s_n\ge s_{n-2}$ at an even step with odd-parity predecessor, and hence Conjecture~\ref{b1:conj:mono} and Theorem~\ref{b1:thm:central}. The binding case $n=15$ has a margin of $1.2\%$. The computational evidence---the counter's geometric growth after its single near miss, the strict growth even across a run of six even-parity steps, two independent simulators in agreement, and the external verification of Appendix~\ref{b1:app:ai}---leaves little doubt that the theorem is true; what remains is the closing of one sharply delimited, structurally identified inequality.

\subsubsection{Algorithms}\label{b1:app:alg}
All computational claims rest on direct implementations of Definition~\ref{b1:def:F}. We summarize the two engines; both were re-implemented independently during the external verification of Appendix~\ref{b1:app:ai}, with agreement on all common outputs.

\paragraph{Bit-level simulator.} Given $L$: set $k=\nu(L)$; if $k$ is even append $00$ and $s\leftarrow 0$, else append $0001$ and $s\leftarrow -1$. List the $1$-positions $p_1<\dots<p_{2m}$; for $j=1,\dots,m$ set positions $p_{2j-1},\dots,p_{2j}-1$ to $1$, set $p_{2j}$ to $0$, and $s\leftarrow s+(p_{2j}-p_{2j-1}-1)$. If $s<2$ halt; otherwise return $L^{(2)}\cdot(0110)^{s-2}$. This is the literal ground truth, feasible while the string fits in memory (through $n=12$ for full strings).

\paragraph{Run-length / gap engine.} Represent $L$ by its run vector $[c_0,r_1,c_1,\dots,r_R,c_R]$, equivalently by the gap-vector $(e;t)$ of \S\ref{b1:sec:10.3}. One step is computed by Proposition~\ref{b1:prop:gap}: pad and seed $s_{\mathrm{init}}$; read $s$ as the seed plus the odd-indexed gap sum; emit blocks $e_{2j-1}+1$ and separators $1+e_{2j}$; append the tail and update $t$ honestly. This runs in time and space $O(\#\text{runs})$, advancing the trajectory well beyond the point where strings can be stored. The value $s_{26}$ was obtained by a memory-frugal pass that sums separators by cumulative-count parity, never materializing the $1.7\times 10^8$-entry gap vector; that pass was validated against the direct value at $n=10$ before its use at $n=26$.

\subsubsection{Computational evidence}\label{b1:app:evidence}
With $L_0=10$, the counter sequence $s_0,\dots,s_{26}$ is
\begin{multline*}
3,5,2,8,31,66,158,392,938,2305,1277,4885,9361,28734,67339,29083,\\
218336,498138,1332397,659539,4146749,3732091,8693523,18136691,\\
24703847,55831137,88613778,
\end{multline*}
and the parity string of $\nu(L_n)$ for $0\le n\le 26$ is (o $=$ odd, E $=$ even)
\[ \mathrm{o\ o\ E\ E\ o\ o\ o\ o\ o\ o\ E\ E\ E\ o\ o\ E\ o\ o\ o\ E\ o\ E\ E\ E\ E\ E\ E.} \]
The $\nu$-recurrence (Proposition~\ref{b1:prop:nu}), iterated from $\nu(L_0)=1$ with these counters, reproduces this parity string exactly, and predicts $\nu(L_{26})=222\,569\,758$, reproduced independently by the separator-parity construction of $s_{26}$. Two-step monotonicity $s_n\ge s_{n-2}$ holds for all $3\le n\le 26$; the tightest instance is $n=15$ ($\approx 1.0121$). For the structural claim of Theorem~\ref{b1:thm:floor}, the reshaped comb was confirmed exactly at $n=3,11,12$.

The odd-run counts $\omega(L_n)$ for $0\le n\le 15$ are
\[ 1,1,4,6,1,11,47,105,235,589,1420,3704,1960,7549,13983,44910; \]
the small values occur only at odd-parity iterates where survival is guaranteed by Lemma~\ref{b1:lem:odd}; at every even-parity iterate the count is well above the value $2$ that a halt would require. The backward decision procedure confirms that all $4565$ halting strings of length at most $14$---including the $3137$ with three or more $1$-runs---are absent from the trajectory.

The exact body/tail decomposition of \S\ref{b1:sec:inert} was verified at the odd-predecessor even steps within reach; Table~\ref{b1:tab:decomp} records the result.

\begin{table}[H]\centering
\begin{tabular}{rrr|l}
$n$ & $s_n$ & body $\Sigma$ / tail $\Sigma$ & nonzero contributing-gap multiset\\
\hline
10 & 1\,277 & 1\,277 / 0 & $\{1:551,\,2:230,\,3:87,\,5:1\}$\\
15 & 29\,083 & 29\,083 / 0 & $\{1:17\,594,\,2:5613,\,3:1,\,4:65\}$\\
19 & 659\,539 & 659\,539 / 0 & $\{1:273\,509,\,2:147\,958,\,3:17\,266,\,4:7769,\,5:810,\,10:319\}$
\end{tabular}
\caption{Body/tail decomposition of the counter at the open (odd-predecessor) even steps. The tail $\Sigma$ is identically $0$ (Lemma~\ref{b1:lem:inert}). For $n=15$, the count term is $23{,}273$ (strictly exceeding $\lceil\omega/2\rceil=22{,}455$) and the size surplus $5613\cdot 1+1\cdot 2+65\cdot 3=5810$, so $s_{15}=23{,}273+5810=29{,}083$.}
\label{b1:tab:decomp}
\end{table}

\begin{table}[H]\centering
\begin{tabular}{rrl|rrl}
$n$ & $s_n$ & par.\ $\nu(L_n)$ & $n$ & $s_n$ & par.\ $\nu(L_n)$\\
\hline
0 & 3 & odd & 14 & 67\,339 & odd\\
1 & 5 & odd & 15 & 29\,083 & \textbf{even}\\
2 & 2 & \textbf{even} & 16 & 218\,336 & odd\\
3 & 8 & \textbf{even} & 17 & 498\,138 & odd\\
4 & 31 & odd & 18 & 1\,332\,397 & odd\\
5 & 66 & odd & 19 & 659\,539 & \textbf{even}\\
6 & 158 & odd & 20 & 4\,146\,749 & odd\\
7 & 392 & odd & 21 & 3\,732\,091 & \textbf{even}\\
8 & 938 & odd & 22 & 8\,693\,523 & \textbf{even}\\
9 & 2\,305 & odd & 23 & 18\,136\,691 & \textbf{even}\\
10 & 1\,277 & \textbf{even} & 24 & 24\,703\,847 & \textbf{even}\\
11 & 4\,885 & \textbf{even} & 25 & 55\,831\,137 & \textbf{even}\\
12 & 9\,361 & \textbf{even} & 26 & 88\,613\,778 & \textbf{even}\\
13 & 28\,734 & odd & & &
\end{tabular}
\caption{The counter $s_n$ and parity of $\nu(L_n)$, $0\le n\le 26$. Even-parity (potentially halting) steps are set in bold. The value $s_{26}=88\,613\,778$ was obtained by a memory-frugal separator-parity pass and corroborated by the independently computed $\nu(L_{26})=222\,569\,758$ via Proposition~\ref{b1:prop:nu}.}
\label{b1:tab:traj}
\end{table}

\subsubsection{Computational programs}\label{b1:app:programs}
The two implementations of Definition~\ref{b1:def:F}---the bit-level simulator and the run-length/gap engine of Appendix~\ref{b1:app:alg}---together with the verification and statistics-gathering scripts, are available at
\begin{center}
\url{https://github.com/RobinCodes/Projects/tree/main/Collatz-type%20problem/Collatz-type%20problem%20v6}.
\end{center}

\subsubsection{Note on the use of artificial intelligence}\label{b1:app:ai}
This work was prepared with substantial assistance from an artificial-intelligence system, and most of the results presented here were created with such assistance. Machine assistance implemented the map of Definition~\ref{b1:def:F} in two independent ways---the bit-level and run-length engines of Appendix~\ref{b1:app:alg}---cross-validated them on every commonly feasible output, extended the computational record of Appendix~\ref{b1:app:evidence}, and helped formulate and prove the structural results: the prefix-alternation invariant (Lemma~\ref{b1:lem:prefix}), the exclusion of the all-even profile (Proposition~\ref{b1:prop:alleven}), the even-step floor (Theorem~\ref{b1:thm:floor}), the $\nu$-recurrence (Proposition~\ref{b1:prop:nu}), the counter-inert tail and its exact decomposition (Lemma~\ref{b1:lem:inert}, Observation~\ref{b1:obs:decomp}), the three reformulations of \S\ref{b1:sec:equiv}, and the localization of the residual obstruction (\S\ref{b1:sec:12.1},~\S\ref{b1:sec:resist}). Computation contributed $s_{26}=88\,613\,778$, $\nu(L_{26})=222\,569\,758$, and the body/tail decompositions of Table~\ref{b1:tab:decomp}.

We do not specify the particular artificial-intelligence system or systems used; readers seeking further detail about the methods are invited to contact the authors: Robin Rovenszky (\texttt{robincodes} on Discord) and notxxdog (\texttt{.notxxdog} on Discord).

Every proof and every computational claim in this paper was independently verified by two human readers and by independent machine re-derivations. An independent re-derivation, working only from the statements of this paper, reconstructed the map of Definition~\ref{b1:def:F} from scratch, re-implemented it in two further from-scratch engines (a literal bit-level engine and a numpy gap-vector engine), cross-validated the two engines against each other, and reproduced every computational and structural claim: the counter sequence and parity through $n=26$ with $s_{26}=88\,613\,778$ and $\nu(L_{26})=222\,569\,758$; the integer-conjugate values $N_0,\dots,N_3=2,998,45\,868\,646,213\,192\,976$ and the conjugacy $T(N_n)=N_{n+1}$ through $n=7$; all structural lemmas of \S\S\ref{b1:sec:rl}--\ref{b1:sec:struct}--\ref{b1:sec:prefix} by brute force on strings of length $\le 12$ beginning with $1$, with zero violations; the body/tail decomposition multisets of Table~\ref{b1:tab:decomp}; and the count/surplus decomposition of Observation~\ref{b1:obs:decomp}, with the count term identified as the number of contributing boundaries (which strictly exceeds $\lceil\omega/2\rceil$ at tight instances). This re-derivation reached the same conclusion regarding the central inequality---that it lies beyond elementary methods.

No result here is accepted on the basis of machine assertion. Every definition, lemma, and proof is stated so as to be checkable by hand, and every computational claim is reproducible by re-running the procedures of Appendix~\ref{b1:app:alg} and the programs of Appendix~\ref{b1:app:programs}. The status of the central open inequality is reported exactly as it stands.


\subsection{Paper 2, Master paper}

\begin{abstract}
We study a deterministic rewriting map $F$ on finite binary strings that carries an integer counter $s$ and halts precisely when $s$ falls below $2$, and we ask, for an arbitrary seed $L_0$, whether the trajectory it generates is infinite.  The map strictly lengthens its argument, so it has no periodic points; every orbit is a simple path that either halts in finite time or diverges, and the entire cycle-elimination half of a Collatz-type problem is absent.  We develop the structural theory at the level of the map: the counter is an alternating, run-length functional (Proposition~\ref{b2:prop:counter}) whose vanishing is a fragile shape (Corollary~\ref{b2:cor:halting}); a single application reshapes runs by explicit identities (\S\ref{b2:sec:struct}); odd parity forces survival and doubling (Lemma~\ref{b2:lem:odd}); and a first halt at index $N$ forces $s_{N-2}\le 5$ (Proposition~\ref{b2:prop:firsthalt}).  We characterize the seeds that halt immediately (Proposition~\ref{b2:prop:step0}), give the prefix-alternation invariant with its parity governed by the seed (Proposition~\ref{b2:prop:prefix}), recast $F$ as a faithful integer conjugate of Collatz type (\S\ref{b2:sec:equiv}), and analyse the backward, predecessor-counting route (\S\ref{b2:sec:backward}).  The distinguished seed $L_0=10$ appears here as one example (\S\ref{b2:sec:seed10}): its counter grazes the halting threshold exactly once and then explodes; we record an independent, machine-verified exclusion of its all-even ($s=0$) halting mode by encoding the process as a $2$-state, $5$-symbol Turing machine and citing a Finite Automata Reduction non-halting certificate from the Busy Beaver Challenge.  We then show that the seed $10$ is atypically tame: for most non-halting seeds two-step monotonicity fails and a positive fraction graze the threshold value $2$ repeatedly and survive (\S\ref{b2:sec:general}).  The residual inequality on which a complete proof founders is the same for every infinite orbit and is a property of $F$, not of any seed (\S\ref{b2:sec:wall}); two independent, cross-validated engines support the empirical claims.  We close with conjectures and open questions.  The central non-halting theorem remains open.
\end{abstract}


\subsubsection{Introduction}\label{b2:sec:intro}
The process studied here is easy to define and stubbornly resistant to complete analysis. A fixed rewriting rule is iterated on binary strings; the rule carries a small integer counter and is engineered to stop the iteration the instant the counter drops below the threshold~$2$. From a seed $L_0$ one forms $L_{n+1}=F(L_n)$ and asks whether the orbit is infinite. We take the seed as a variable and study the map itself, with the distinguished seed $L_0=10$ appearing as a worked example (\S\ref{b2:sec:seed10}).

The architecture of the analysis is the following. Almost every structural fact is a property of one application of $F$, hence of the map and of every seed alike; the seed enters only through a base case, a finite list of small counter values, and the shape of the orbit's encounter with the threshold. Modulo one structural input---that a particular all-even halting value is not attained---the question for a given seed is equivalent to a single inequality on the size of the counter, driven by a transparent geometric mechanism. We establish the seed-independent structure, prove the governing inequality in a structurally identified sub-case, characterize the seeds that halt immediately, and pin down sharply the one configuration on which a complete proof still founders---and which is the same for every seed.

\subsubsection*{Conventions} Strings are finite words over $\{0,1\}$; positions are numbered $1,2,3,\dots$ from the left. For $w\in\{0,1\}^*$ we write $\nu(w)$ for the number of $1$'s and $\omega(w)$ for the number of odd-length maximal runs of $1$'s. Concatenation is juxtaposition or a dot, and $u^k$ is the $k$-fold concatenation. We call a string \emph{even} (resp.\ \emph{odd}) when $\nu$ is even (resp.\ odd); this ``parity'' refers to $\nu$. Leading $0$'s are never modified by $F$ and ride along as an inert prefix, so without loss of generality every seed begins with $1$; trailing $0$'s affect the odd pad and the counter and are tracked honestly. For $L$ with $\nu(L)\ge 1$, $r_1(L)$ is the length of the first maximal $1$-run and $c_1(L)$ the length of the first internal $0$-run.

\begin{definition}[The map $F$]\label{b2:def:F}
For $L\in\{0,1\}^*$, the value $F(L)$, when defined, is computed in three steps.

\emph{Step 1 (pad; seed the counter).} Let $k=\nu(L)$. If $k$ is even, append $00$ and set $s:=0$; if $k$ is odd, append $0001$ and set $s:=-1$. Call the result $L^{(1)}$; note $\nu(L^{(1)})$ is even.

\emph{Step 2 (pair the $1$'s; fill and count).} Let $p_1<\dots<p_{2m}$ be the positions of the $1$'s in $L^{(1)}$, grouped into consecutive pairs. For $j=1,\dots,m$: set positions $p_{2j-1},\dots,p_{2j}-1$ to $1$; set $p_{2j}$ to $0$; add $p_{2j}-p_{2j-1}-1$ to $s$. Call the result $L^{(2)}$.

\emph{Step 3 (halt or extend).} If $s<2$, the process halts and $F(L)$ is undefined; otherwise $F(L):=L^{(2)}\cdot(0110)^{s-2}$.
\end{definition}

\begin{definition}[Trajectory]\label{b2:def:traj}
Fix a seed $L_0$ and set $L_{n+1}:=F(L_n)$ whenever the right-hand side is defined. Write $s_n$ for the counter computed during $F(L_n)$. The seed $L_0$ is \emph{non-halting} if $L_n$ is defined for all $n\ge 0$, equivalently $s_n\ge 2$ for all $n$.
\end{definition}

\begin{thm}[Central question, open]\label{b2:thm:central}
There exist non-halting seeds. We conjecture that $L_0=10$ is one; we do not prove it, and we believe it is not within reach of elementary methods. What we provide is the seed-independent structure, a reduction to a single scalar inequality, and a sharply delimited residual obstruction common to all seeds.
\end{thm}

\subsubsection{Summary of results}
\begin{enumerate}[label=(\roman*)]
\item No cycles, hence a clean dichotomy (Theorem~\ref{b2:thm:nocycles}): $F$ strictly lengthens, so it has no periodic points; every orbit is a simple path, finite (halting) or diverging.
\item The run-length calculus (Proposition~\ref{b2:prop:counter}, Corollary~\ref{b2:cor:halting}) and the reshaping identities (\S\ref{b2:sec:struct}), all properties of the map.
\item Immediate halts characterized (Proposition~\ref{b2:prop:step0}).
\item Prefix alternation with seed-controlled parity (Proposition~\ref{b2:prop:prefix}) and the all-even exclusion (\S\ref{b2:sec:prefix}); $s=0$ is not excluded for all seeds ($F(1)=11110$), but is excluded on the orbit of $10$ by two independent routes, including a Turing-machine encoding certified non-halting by Finite Automata Reduction (Remark~\ref{b2:rem:TM}).
\item The reduction to $s_{N-2}\le 5$ at a first halt (Proposition~\ref{b2:prop:firsthalt}) and an even-step floor giving two-step monotonicity at even-after-even steps (Theorem~\ref{b2:thm:floor}).
\item The counter-inert tail and body localization (Lemma~\ref{b2:lem:inert}, Observation~\ref{b2:obs:decomp}), the $\nu$-recurrence (Proposition~\ref{b2:prop:nu}), and three exact conjugacies (\S\ref{b2:sec:equiv}).
\item The backward viewpoint and the first-halt analysis (\S\ref{b2:sec:backward}).
\item The distinguished seed $10$ (\S\ref{b2:sec:seed10}) and the demonstration that it is atypical: two-step monotonicity and a single threshold graze are not generic (\S\ref{b2:sec:general}).
\item The wall (\S\ref{b2:sec:wall}): the residual even-after-odd inequality is a property of $F$; every infinite orbit meets it. Conjectures and open questions follow (\S\ref{b2:sec:conj}--\S\ref{b2:sec:open}).
\end{enumerate}

\subsubsection{Elementary structure and no cycles}\label{b2:sec:elem}
\begin{lemma}[$F$ strictly lengthens]\label{b2:lem:length}
Whenever $F(L)$ is defined, $|F(L)|\ge|L|+2$.
\end{lemma}
\begin{proof}
Step 1 appends two or four symbols; Step 2 preserves length; Step 3 appends $4(s-2)\ge 0$ symbols.
\end{proof}

\begin{lemma}[Outputs end in $0$]\label{b2:lem:end0}
Whenever $F(L)$ is defined, its last symbol is $0$; hence every defined iterate ends in $0$.
\end{lemma}
\begin{proof}
Step 2 sets the position of the last $1$ of $L^{(1)}$ to $0$ and no later position holds a $1$; if $s=2$ then $F(L)=L^{(2)}$ ends in that $0$, and if $s>2$ it ends in the final $0$ of the last block $0110$.
\end{proof}

\begin{thm}[No cycles; the dichotomy]\label{b2:thm:nocycles}
For every seed, the iterates $L_0,L_1,\dots$ are pairwise distinct. Hence $F$ has no periodic points and no nontrivial cycles, and each orbit is a simple path that is either finite (it reaches a string on which $F$ is undefined) or infinite with $|L_n|\to\infty$ strictly.
\end{thm}
\begin{proof}
By Lemma~\ref{b2:lem:length} the lengths strictly increase along any orbit, so no string recurs; in particular $F^k(X)=X$ ($k\ge 1$) is impossible. A deterministic, non-repeating orbit is a simple path, terminating exactly at a halt and otherwise continuing forever.
\end{proof}

A first consequence, central to the backward viewpoint (\S\ref{b2:sec:backward}): every string has only finitely many iterated preimages under $F$, since a preimage is strictly shorter; thus orbit-membership is decidable. Theorem~\ref{b2:thm:nocycles} also removes from the present problem the cycle-elimination task that is the heart of the Collatz problem: here only halt versus diverge is at issue. We write $\mathcal H$ for the set of halting seeds.

\subsubsection{Odd parity forces survival}\label{b2:sec:odd}
\begin{lemma}[Odd-parity lemma]\label{b2:lem:odd}
If $\nu(L)$ is odd, then the counter computed during $F(L)$ satisfies $s\ge 2$. Moreover if $L=W\cdot(0110)^u$ with $u\ge 1$ and $W$ ending in $0$, then $s\ge 2u+2$.
\end{lemma}
\begin{proof}
Write $k=\nu(L)$ odd and $L=L_f\cdot 1\cdot 0^b$ with $b\ge 0$. Step 1 appends $0001$, giving $L^{(1)}=L_f\cdot 1\cdot 0^{b+3}\cdot 1$; now $\nu(L^{(1)})$ is even and the last pair contributes $p_{2m}-p_{2m-1}-1=b+3\ge 3$. Since every pair contributes $\ge 0$ and $s$ began at $-1$, $s\ge 2$. The growth bound follows by counting the cross pair joining $W$ to the tail and the $u-1$ interior tail pairs of gap $2$.
\end{proof}

From here only even-parity strings can halt. For even parity, $s=\sum_j(p_{2j}-p_{2j-1}-1)\ge 0$, a sum of nonnegative within-pair gaps; nothing forbids $s\in\{0,1\}$, and locating those is the next task.

\subsubsection{The run-length calculus}\label{b2:sec:rl}
Write the run-length encoding of $L$ (with at least one $1$) as $L=0^{c_0}1^{r_1}0^{c_1}\cdots 1^{r_R}0^{c_R}$, with $R\ge 1$, $r_i\ge 1$, internal $c_i\ge 1$. Put $C_i=r_1+\cdots+r_i$, so $C_0=0$, $C_R=\nu(L)$.

\begin{proposition}[Counter formula]\label{b2:prop:counter}
If $\nu(L)$ is even, the counter computed during $F(L)$ is
\[ s=\sum_{1\le i\le R-1,\ C_i\text{ odd}} c_i. \]
\end{proposition}
\begin{proof}
Step 1 adds no $1$'s, so the $1$'s of $L$ pair as $\{1,2\},\{3,4\},\dots$. The boundary between runs $i$ and $i+1$ is spanned by a pair iff $C_i$ is odd, contributing $c_i$; same-run pairs contribute $0$; the index $i=R$ does not arise. Summing gives the claim.
\end{proof}

\begin{corollary}[Halting dichotomy]\label{b2:cor:halting}
Let $\nu(L)$ be even. Then $F(L)$ halts iff exactly one of: (i) every run length is even ($s=0$, equivalently $\omega(L)=0$); or (ii) the run lengths contain exactly two odd values, at consecutive positions $j_0,j_0+1$, with $c_{j_0}=1$ ($s=1$).
\end{corollary}
\begin{proof}
By Proposition~\ref{b2:prop:counter} and $c_i\ge 1$ internally, $s$ counts internal indices with $C_i$ odd, weighted by $c_i$. The parity sequence $C_0,\dots,C_R$ starts and ends even and toggles at odd-length runs; the two stated profiles are exactly $s=0$ and $s=1$.
\end{proof}

\subsubsection{Structural lemmas}\label{b2:sec:struct}
Fix $L$, let $\Lambda=L\cdot(\text{pad})$, let the consecutive pairs of $1$'s in $\Lambda$ have within-pair gaps $g_j\ge 0$, $j=1,\dots,m=\tfrac12\nu(\Lambda)$, and set $X(L)=\#\{i:C_i\text{ odd and }\gamma_i\text{ odd}\}$ where $\gamma_i$ are the internal $0$-runs of $\Lambda$.

\begin{lemma}[Reshaping]\label{b2:lem:reshape}
The image $L^{(2)}$ has exactly $m$ runs of $1$'s, the $j$-th of length $g_j+1$; hence $F(L)$ has $1$-runs $g_1+1,\dots,g_m+1$ followed by $s-2$ runs of length $2$. A body run is odd iff its pair has even gap, and the tail contributes only even-length runs.
\end{lemma}

\begin{lemma}[Reshaping count]\label{b2:lem:reshapecount}
For every $L$, $\omega(F(L))=m-X(L)$.
\end{lemma}

\begin{lemma}[Odd-run floor]\label{b2:lem:oddrunfloor}
With $\mu$ the number of pairs lying inside a single $1$-run of $\Lambda$ and $R'$ the number of $1$-runs of $\Lambda$, $F(L)$ has at least $\mu$ odd-length $1$-runs, and $\mu\ge m-(R'-1)$.
\end{lemma}

\begin{lemma}[Tail manufactures odd runs]\label{b2:lem:tailodd}
If $L=W\cdot(0110)^u$ with $u\ge 1$ and $W$ ending in $0$, then $\omega(F(L))\ge u-\varepsilon$ with $\varepsilon=0$ if $\nu(L)$ even and $\varepsilon=1$ if $\nu(L)$ odd.
\end{lemma}

\begin{lemma}[Counter floor]\label{b2:lem:cfloor}
If $\nu(L)$ is even, then $s(F(L))\ge\lceil\omega(L)/2\rceil$.
\end{lemma}

Proofs are the pairing analysis of \S\ref{b2:sec:rl} applied run by run; they are routine and we omit them, having verified each against the engines of Appendix~\ref{b2:app:alg}.

\subsubsection{Immediate halts}\label{b2:sec:step0}
\begin{proposition}[Step-$0$ halts]\label{b2:prop:step0}
$F(L_0)$ is undefined iff $\nu(L_0)$ is even and either (i) every maximal $1$-run of $L_0$ has even length, or (ii) $L_0$ has exactly two odd-length $1$-runs, adjacent and separated by a single $0$.
\end{proposition}
\begin{proof}
A halt forces $\nu(L_0)$ even (Lemma~\ref{b2:lem:odd}); Step 1's pad $00$ leaves the $1$-run structure unchanged, so Corollary~\ref{b2:cor:halting} applies to $L_0$ directly.
\end{proof}

The step-$0$ halts form a regular language. Among the $16383$ seeds of length $\le 14$ (beginning with $1$), $4859$ ($\approx 29.7\%$) halt, the halt being early (Table~\ref{b2:tab:stats}); the fraction declines with length (Table~\ref{b2:tab:bylength}). Membership in $\mathcal H$ is recursively enumerable by forward simulation; its decidability is open (Question~\ref{b2:q:dec}).

\subsubsection{Prefix dynamics and the all-even profile}\label{b2:sec:prefix}
\begin{proposition}[Prefix alternation]\label{b2:prop:prefix}
For all $n\ge 0$, $r_1(L_{n+1})=1$ if $r_1(L_n)\ge 2$, and $r_1(L_{n+1})=c_1(L_n)+1\ (\ge 2)$ if $r_1(L_n)=1$. With $a_n:=[\,r_1(L_n)\ge 2\,]$ one has $a_{n+1}=\lnot a_n$, hence $a_n=a_0\oplus[n\text{ odd}]$ and
\[ r_1(L_n)=1\iff \begin{cases} n\text{ even}, & r_1(L_0)=1,\\ n\text{ odd}, & r_1(L_0)\ge 2. \end{cases} \]
\end{proposition}
\begin{proof}
The recurrence follows from Lemma~\ref{b2:lem:reshape} applied to the first pair: if $r_1(L_n)\ge 2$ the first two $1$'s are adjacent (gap $0$, run length $1$); if $r_1(L_n)=1$ the first pair straddles the first boundary, giving run length $c_1(L_n)+1$. The closed form solves $a_{n+1}=\lnot a_n$.
\end{proof}

\begin{proposition}[All-even exclusion]\label{b2:prop:alleven}
Recall $s_n=0\iff\omega(L_n)=0$. At every index $n$ with $r_1(L_n)=1$ one has $\omega(L_n)\ge 1$, hence $s_n\ne 0$; by Proposition~\ref{b2:prop:prefix} these are the even indices when $r_1(L_0)=1$ and the odd indices when $r_1(L_0)\ge 2$. The exclusion is \emph{not universal} across seeds: the all-even profile does occur as an output of $F$ for some inputs.
\end{proposition}
\begin{proof}
$r_1(L_n)=1$ is an odd-length leading run, so $\omega(L_n)\ge 1$ and $s_n\ne 0$ by Corollary~\ref{b2:cor:halting}. For non-universality, see Remark~\ref{b2:rem:rem18}.
\end{proof}

\begin{remark}[The all-even profile is reachable as an output]\label{b2:rem:rem18}
One might hope that every output of $F$ carries a $0$ strictly between its first and last $1$---so that no iterate is a single $1$-run $1^a0^b$, and the all-even profiles would be unreachable from the range of $F$. This is not so. The odd pad $0001$ places the (subsequently zeroed) last $1$ at the very end, leaving no internal $0$; explicitly,
\[ F(1)=11110=1^40\quad(\text{an output of }F,\ \omega=0),\qquad F(11110)=\text{halt with }s=0. \]
Thus the all-even profile is reachable as an output, and the exclusion of $s=0$ is a property of particular orbits, not of the range of $F$. The exclusion therefore rests on the prefix argument of Proposition~\ref{b2:prop:alleven}, together with the orbit-of-$10$ certificate of Remark~\ref{b2:rem:TM}.
\end{remark}

\begin{proposition}[Reduction to a single profile]\label{b2:prop:single}
On any orbit that excludes the all-even profile, the trajectory halts iff some $s_n=1$.
\end{proposition}
\begin{proof}
A halt has $s_n\in\{0,1\}$ (Lemma~\ref{b2:lem:odd}, Corollary~\ref{b2:cor:halting}); excluding $s_n=0$ leaves $s_n=1$.
\end{proof}

\subsubsection{The reduction to a scalar inequality}\label{b2:sec:reduction}
\begin{proposition}[First-halt bound]\label{b2:prop:firsthalt}
If the orbit of any seed first halts at index $N\ge 2$, then $s_{N-2}\le 5$.
\end{proposition}
\begin{proof}
At the halt $\nu(L_N)$ is even and $s_N\in\{0,1\}$, so $\omega(L_N)\in\{0,2\}\le 2$ (Corollary~\ref{b2:cor:halting}). Since $N\ge 2$, $L_{N-1}=L^{(2)}_{N-1}\cdot(0110)^u$ with $u=s_{N-2}-2\ge 0$ and $L^{(2)}_{N-1}$ ending in $0$. If $u=0$, $s_{N-2}=2$; if $u\ge 1$, Lemma~\ref{b2:lem:tailodd} gives $\omega(L_N)\ge u-1=s_{N-2}-3$, so $s_{N-2}-3\le 2$.
\end{proof}

\begin{corollary}[Remaining content]\label{b2:cor:remaining}
A seed is non-halting provided $s_n\ge 6$ for all $n\ge 3$ and $s_0,\dots,s_2\ge 2$; a fortiori it is non-halting if the counter satisfies the two-step monotonicity $s_n\ge s_{n-2}$ for all $n\ge 3$ together with $s_1,s_2\ge 2$.
\end{corollary}
\begin{proof}
By Proposition~\ref{b2:prop:firsthalt} a first halt forces $s_{N-2}\le 5$, so $N\le 4$ once $s_n\ge 6$ for $n\ge 3$; checking $s_0,s_1,s_2$ excludes early halts. Monotonicity propagates the small even/odd subsequence values.
\end{proof}

The mechanism behind the inequality is explicit. At an odd-parity step, Lemma~\ref{b2:lem:odd} gives $s_n\ge 2s_{n-1}-2$, so odd steps at least double the counter. The counter can lose ground only at an even-parity step, where Lemmas~\ref{b2:lem:cfloor} and~\ref{b2:lem:tailodd} give $s_n\ge\lceil(s_{n-2}-3)/2\rceil$, at most a near-halving. A long run of even steps could in principle compound several near-halvings; the next results control the benign case.

\begin{thm}[Even-step floor with even-parity predecessor]\label{b2:thm:floor}
Suppose step $n$ is even parity, step $n-1$ is also even parity, and $u:=s_{n-2}-2\ge 1$. Then the last $u$ body $1$-runs of $L_n$ have length $1$ and are separated by internal $0$-runs of length exactly $3$, whence
\[ s_n\ge 3\Bigl\lfloor\frac{u-1}{2}\Bigr\rfloor\ge \tfrac32 s_{n-2}-6, \]
so $s_n\ge s_{n-2}$ whenever $s_{n-2}\ge 12$. Thus two-step monotonicity holds at every even-after-even step, with finitely many small cases to check directly.
\end{thm}
\begin{proof}
The tail $(0110)^u$ of $L_{n-1}$ carries $2u$ ones; since both steps are even parity, the body $1$'s pair among themselves and the tail $1$'s pair within the tail, each block $11$ reshaping to a single $1$ followed by the zeroed endpoint and the separating $00$---three $0$'s. The image of the tail region is $(1\,000)^{u-1}1$: $u$ single-$1$ runs separated by $0$-runs of length $3$. Across these the cumulative count rises by $1$ per run, so at least $\lfloor(u-1)/2\rfloor$ of the $u-1$ separators sit at odd cumulative count and contribute $3$ each to the counter formula (Proposition~\ref{b2:prop:counter}). The arithmetic bound and threshold $s_{n-2}\ge 12$ follow.
\end{proof}

\subsubsection{The counter-inert tail and body localization}\label{b2:sec:inert}
\begin{lemma}[Counter-inert tail]\label{b2:lem:inert}
Let step $n$ be of even parity, and write $L_n=B\cdot T$ with $T=(0110)^{s_{n-1}-2}$ the tail appended in forming $L_n$ and $B=L^{(2)}_{n-1}$ the body (so $B$ ends in $0$). Then $T$ contributes exactly $0$ to $s_n$, and $s_n$ equals the within-pair-gap sum over $B$ alone.
\end{lemma}
\begin{proof}
Step 1 appends $00$ and the $1$'s of $L_n$ pair consecutively from the left. The tail carries $2(s_{n-1}-2)$ ones, an even number, and $\nu(L_n)$ is even, so $\nu(B)$ is even; the $B$-ones pair among themselves and the tail ones pair within the tail, each tail pair being a block $11$ of gap $0$.
\end{proof}

\begin{observation}[Body localization]\label{b2:obs:decomp}
At an even-parity step $n$ the nonzero within-pair gaps occupy exactly the body's pair-slots. Each contributing pair has within-pair gap $e_{2j-1}\ge 1$; decomposing this gap as $1+(e_{2j-1}-1)$ partitions the counter into a count term and a size-surplus term:
\begin{equation}\label{b2:eq:decomp}
s_n \;=\; \underbrace{\#\{\text{body boundaries with odd cumulative $1$-count}\}}_{\text{count floor}}
\;+\;\underbrace{\sum_{\text{those boundaries}}\bigl(e_{2j-1}(\Lambda_{n-1})-1\bigr)}_{\text{size surplus}}.
\end{equation}
The count floor is bounded below by $\lceil\omega(L_n)/2\rceil$ (Lemma~\ref{b2:lem:cfloor}), with equality only when each parity-raised interval of the cumulative-$1$-count contributes a single boundary; at tight instances the actual count strictly exceeds this bound. The size surplus is carried by the contributing within-pair gaps of size $\ge 2$. As an illustration, at $n=15$ the count floor is $23{,}273$, the lower bound $\lceil\omega(L_{15})/2\rceil=22{,}455$, the surplus is $5810$, and $s_{15}=23{,}273+5810=29{,}083$ (Appendix~\ref{b2:app:evidence}, Table~\ref{b2:tab:decomp}).
\end{observation}

\subsubsection{The $\nu$-recurrence and parity dynamics}\label{b2:sec:nu}
\begin{proposition}[$\nu$-recurrence]\label{b2:prop:nu}
Whenever $F(L)$ is defined, with $\nu=\nu(L)$ and counter $s$,
\[ \nu(F(L))=\begin{cases}\tfrac{\nu}{2}+3s-4, & \nu\text{ even},\\ \tfrac{\nu+1}{2}+3s-3, & \nu\text{ odd}.\end{cases} \]
Consequently the parity of $\nu(L_{n+1})$ is determined by $\nu(L_n)\bmod 4$ and $s_n\bmod 2$; for even parity, $\nu(L_{n+1})\equiv\nu(L_n)/2+s_n\pmod 2$.
\end{proposition}
\begin{proof}
$\nu(F(L))=\nu(L^{(2)})+2(s-2)$; by Lemma~\ref{b2:lem:reshape}, $\nu(L^{(2)})=\nu/2+s$ (even) or $(\nu+1)/2+s+1$ (odd). Substitute.
\end{proof}

The parity of $\nu(L_n)$ decides which regime---odd-step doubling or even-step near-halving---governs step $n$, yet the counter $s_n$ obeys no scalar recurrence: by Proposition~\ref{b2:prop:counter} it depends on the full pattern of odd-cumulative gaps of $L_n$. This non-locality is the structural heart of the difficulty.

\subsubsection{Equivalent formulations}\label{b2:sec:equiv}
Let $\popcount(x)$ be the popcount of $x$ and $V(L)=\mathrm{int}(L,2)$ the binary value (most significant bit first). For $M$ with $\popcount(M)=2m$ and $1$-bit positions $b_1>\dots>b_{2m}\ge 0$, set $S(M)=\sum_i(-1)^{i+1}2^{b_i}$.

\begin{lemma}[Step 2 is doubling of the alternating bit-sum]\label{b2:lem:S}
If $\nu(L)$ is even and $M=V(L^{(1)})$, then $V(L^{(2)})=2S(M)$. Each pair $(b_{2j-1},b_{2j})$ effects the local change $2^{b_{2j-1}}-3\cdot 2^{b_{2j}}$ to $M$.
\end{lemma}

The constant $3$ here is the genuine arithmetic kernel of the process---the counterpart of the $3$ in the Collatz map.

\begin{proposition}[Integer conjugate]\label{b2:prop:T}
Each defined iterate begins with $1$ (Proposition~\ref{b2:prop:prefix}) and ends in $0$ (Lemma~\ref{b2:lem:end0}), so $V$ is a bijection of the reachable strings onto a subset of $\mathbb Z^+$. Through $V$, $F$ is conjugate to a Collatz-type map $T:\mathbb Z^+\dashrightarrow\mathbb Z^+$ given by
\[ T(N)=2S(M)\cdot 16^{s-2}+6\cdot\frac{16^{s-2}-1}{15}, \]
where $M=4N$ if $\popcount(N)$ even and $M=16N+1$ if $\popcount(N)$ odd, and $s=s_{\mathrm{init}}+\popcount(S(M))-\tfrac12\popcount(M)$ with $s_{\mathrm{init}}=0$ (resp.\ $-1$) for $\popcount(N)$ even (resp.\ odd); $T$ is defined when $s\ge 2$. In hexadecimal $T(N)$ is the digit string of $2S(M)$ followed by $s-2$ copies of the digit $6$. The seed is $N_0=V(L_0)$, and the question is whether the $T$-orbit of $N_0$ is infinite.
\end{proposition}

A natural-scale gap-vector conjugate and the scalar-counter form are obtained likewise; the popcount-parity test selecting one of two affine pads, with the kernel constant $3$, is the precise sense in which the process is of Collatz type.

\begin{observation}[No mergers in range]\label{b2:obs:inj}
On all strings of length $\le 11$ beginning with $1$, $F$ is injective: no two share an image. With Theorem~\ref{b2:thm:nocycles} the reachable graph is a disjoint forest of simple chains; distinct orbits do not merge in range (Conjecture~\ref{b2:conj:inj}).
\end{observation}

\subsubsection{The backward viewpoint}\label{b2:sec:backward}
Preimages are strictly shorter (Lemma~\ref{b2:lem:length}), so the backward tree of any string is finite and $F$ is exactly invertible.

\begin{proposition}[Invertibility]\label{b2:prop:inv}
Every $X$ with $F(X)=T$ arises by choosing $u\ge 0$ with $T=Y\cdot(0110)^u$, $Y$ ending in $0$; reconstructing $L^{(1)}$ from the maximal $1$-blocks of $Y$; and deleting a Step-$1$ suffix; one retains $X$ iff $F(X)=T$. There are finitely many preimages.
\end{proposition}

\begin{proposition}[Shape of a first halt]\label{b2:prop:halt-shape}
Suppose the orbit of a seed first halts at $N\ge 2$. Then $s_{N-2}\le 5$ (Proposition~\ref{b2:prop:firsthalt}), and by Lemma~\ref{b2:lem:inert} the halt depends only on the body $B=L^{(2)}_{N-1}$: the tail $(0110)^{s_{N-1}-2}$ of $L_N$ contributes $0$ to $s_N$. Consequently the predecessor counter $s_{N-1}$ is not determined by the halt and may be arbitrarily large.
\end{proposition}
\begin{proof}
Proposition~\ref{b2:prop:firsthalt} gives the first claim. At the halt $\nu(L_N)$ is even, so Lemma~\ref{b2:lem:inert} applies to the step computing $s_N$, which therefore depends on $B$ alone, independently of the tail length $s_{N-1}-2$.
\end{proof}

\begin{remark}[The predecessor counter is not pinned at a first halt]\label{b2:rem:rem31}
One might expect a first halt to force $s_{N-1}=2$, so that the halting string has an empty tail $L_N=L^{(2)}_{N-1}$. This is not so, and would contradict Lemma~\ref{b2:lem:inert}, which makes the halt independent of $s_{N-1}$. The seed $110011111110$ first halts at $N=8$ with $s_{N-1}=102$ (and halt value $s_N=1$); the seed $1100111$ first halts at $N=2$ with $s_{N-1}=3$. The subtlety is that the tail's $11$ blocks are even runs of $L_N$, not odd runs, so Lemma~\ref{b2:lem:cfloor} gives only $\lceil\omega/2\rceil=1$ rather than $s_N\ge 2$; the halting strings are $B\cdot(0110)^k$ ($k\ge 0$) with $B$ fragile. The finite backward decision procedure---tracing the finite preimage tree of any fixed string to its roots and testing membership of the seed---is exact regardless and dispatches, for the orbit of $10$, every halting string of length $\le 14$.
\end{remark}

\subsubsection{The distinguished seed $L_0=10$}\label{b2:sec:seed10}
The seed $10$ deserves separate attention as the smallest interesting case.

Running the rule once: $\nu(10)=1$ is odd, so Step 1 appends $0001$, $s=-1$, $L^{(1)}=100001$; filling positions $1$--$5$ and zeroing $6$ gives $111110$ with filled gap $4$, so $s=3$, and $L_1=1111100110$. Continuing,
\[ L_2=10101110111110011001100110,\qquad L_3=1100101101010001000100010000, \]
and the counter takes the values $s_0,s_1,s_2,s_3,s_4,\dots=3,5,2,8,31,66,158,\dots$ (Table~\ref{b2:tab:traj}).

Two features organize this orbit. It is \emph{explosive}: from $s_3$ on the counter roughly multiplies each step, and $|L_n|$ exceeds $10^8$ before $n=22$. And it \emph{grazes the threshold exactly once}: $s_2=2$, with no margin, after which the counter never returns to $\{2,3,4,5\}$ through $n=26$. Here $r_1(L_0)=1$, so by Proposition~\ref{b2:prop:prefix} the all-even exclusion (Proposition~\ref{b2:prop:alleven}) is unconditional at even $n$; at the odd indices the small computed values discharge the remaining cases, so the orbit of $10$ excludes the all-even profile throughout. Two-step monotonicity $s_n\ge s_{n-2}$ holds for all $3\le n\le 26$; the tightest instance is the odd$\to$even transition $n=15$, with $s_{15}/s_{13}=29083/28734\approx 1.012$.

\begin{remark}[A Turing-machine certificate for the all-even exclusion]\label{b2:rem:TM}
The exclusion of the all-even ($s=0$) halt on the orbit of $10$ admits a rigorous, machine-checked proof independent of the prefix argument. Encode the process as the $2$-state, $5$-symbol Turing machine
\[ \texttt{1RB3LA1RA4LA2RA\quad 2LA---1LA0RA3RB} \]
(states $A,B$; symbols $0,\dots,4$; the entry \texttt{---} is the undefined, halting transition, on reading $1$ in state $B$), started on the blank tape. By construction, the machine simulates the trajectory and is arranged so that its halting transition is reached precisely when the simulated string attains an all-even configuration; thus, \emph{by construction}, the machine's halting is equivalent to the occurrence of an $s=0$ event. The non-halting of this machine was established by Finite Automata Reduction (FAR), a non-halting decision method developed within the Busy Beaver Challenge collaboration \cite{b2:bbc}: FAR exhibits a finite automaton over tape configurations that contains the initial configuration, is closed under the transition function, and contains no halting configuration, yielding a fully verifiable certificate. The FAR certificate is the formal object of record; an independent reproduction confirmed that the machine runs without halting for $10^7$ steps, consistent with the certificate. Conditional on the certificate---and hence on the constructed encoding equivalence---the orbit of $10$ never attains an $s=0$ configuration, with no length bound. Together with Proposition~\ref{b2:prop:single}, this reduces the halting question for the orbit of $10$ to the $s=1$ profile alone.
\end{remark}

\subsubsection{General seeds: the seed $10$ is atypical}\label{b2:sec:general}
The single threshold graze and the two-step monotonicity of the orbit of $10$ are not generic.

\begin{observation}[Failure of two-step monotonicity]\label{b2:obs:mono}
Among the $1136$ non-halting seeds of length $\le 11$, at least $586$ violate $s_n\ge s_{n-2}$ for some $n\ge 3$ within the first $80$ simulated steps, and at least $659$ do so within $150$ steps. The count is non-decreasing as the simulation is extended further, so the exact total is depth-dependent; the robust content is that two-step monotonicity (the clean sufficient condition of Corollary~\ref{b2:cor:remaining}) is failed by a substantial majority of non-halting seeds, hence is a feature of the orbit of $10$, not an invariant available for general seeds.
\end{observation}

\begin{observation}[Repeated threshold-grazing]\label{b2:obs:graze}
Define $g(L_0)=\#\{n\ge 1:s_n=2\}$. For $217$ of the $1136$ non-halting seeds of length $\le 11$, $g\ge 1$, with observed values up to $3$. The extreme seed $1100001110$ has
\[ s=3,10,2,19,2,42,2,87,340,77,836,\dots, \]
touching the threshold value $2$ at $n=2,4,6$ before exploding; many such seeds return to $s_n\le 5$ after first exceeding $5$. Hence even ``$s_n\ge 6$ eventually'' fails generically, and the only robust criterion is the definitional $s_n\ge 2$ for all $n$.
\end{observation}

A seed whose orbit were eventually of all-odd parity would be provably non-halting, by the doubling bound $s_n\ge 2s_{n-1}-2$; the longest all-odd prefix we find is $13$ steps (seed $1100011010$), after which the parity turns even. No provably non-halting seed is known (Question~\ref{b2:q:provable}).

\subsubsection{Why the problem resists}\label{b2:sec:wall}
The obstruction is a mismatch between the quantity one must control and the quantities one can compute, and it is the same for every seed. To forbid both fragile profiles at an even step it suffices to know that every even-parity iterate carries at least four odd-length runs; but the only uniform lower bound, $\mu\ge m-(R'-1)$ (Lemma~\ref{b2:lem:oddrunfloor}), collapses for near-combs $1010\cdots10$. The decisive quantity is the size surplus of Observation~\ref{b2:obs:decomp}, not a count: at a tight instance the \emph{provable} count-based lower bound $\lceil\omega/2\rceil\approx\tfrac12 s_{n-2}$ on the count floor is too small to close the gap with $s_{n-2}$, and the whole margin is carried by within-pair gaps of size $\ge 2$ at odd boundaries---an alignment depending on the full history, a quantity with no scalar recurrence (Proposition~\ref{b2:prop:nu}). At $n=15$ this is concrete: $\lceil\omega(L_{15})/2\rceil=22{,}455<s_{13}=28{,}734$, so any argument bounding the count floor only from below by $\lceil\omega/2\rceil$ is dead on arrival, while the actual count floor $23{,}273$ together with the size surplus $5810$ exceeds $s_{13}$ by the margin $349$. The counter is a global parity functional, small only on configurations of a very particular shape; there is no scalar recurrence; and the integer orbit grows doubly exponentially, so finite verification cannot substitute for a uniform argument. Crucially, this residual inequality is a property of $F$ near fragile bodies, not of any seed: across every non-halting seed within reach, the even-parity step with an odd-parity predecessor eventually occurs (the all-odd escape hatch being unrealized), and that is exactly where the argument fails.

\begin{conjecture}[The wall, seed-independent form]\label{b2:conj:wall}
For every seed whose orbit is infinite, and every even-parity step $n$ whose predecessor is of odd parity, $s_n\ge s_{n-2}$.
\end{conjecture}

\subsubsection{Computational results}\label{b2:sec:comp}
Two independent implementations of Definition~\ref{b2:def:F}---a bit-level simulator and a memory-frugal gap-vector engine (Appendix~\ref{b2:app:alg})---agree on every common output: full strings through $n=12$, and the scalar quantities $s_n$ and $\nu(L_n)\bmod 2$ through $n=25$, with the gap engine cross-checked against Proposition~\ref{b2:prop:nu}. They were validated against each other on all $8191$ seeds of length $\le 13$ with zero discrepancies. For the orbit of $10$ the record reaches $s_{26}=88{,}613{,}778$, corroborated by $\nu(L_{26})=222{,}569{,}758$; two-step monotonicity holds for all $3\le n\le 26$; and a run of six consecutive even-parity steps occurs ($n=21,\dots,26$) across which the counter nevertheless grows.

\begin{table}[H]\centering
\begin{tabular}{rrl|rrl}
$n$ & $s_n$ & par.\ $\nu(L_n)$ & $n$ & $s_n$ & par.\ $\nu(L_n)$\\
\hline
0 & 3 & odd & 14 & 67\,339 & odd\\
1 & 5 & odd & 15 & 29\,083 & even\\
2 & 2 & even & 16 & 218\,336 & odd\\
3 & 8 & even & 17 & 498\,138 & odd\\
4 & 31 & odd & 18 & 1\,332\,397 & odd\\
5 & 66 & odd & 19 & 659\,539 & even\\
6 & 158 & odd & 20 & 4\,146\,749 & odd\\
7 & 392 & odd & 21 & 3\,732\,091 & even\\
8 & 938 & odd & 22 & 8\,693\,523 & even\\
9 & 2\,305 & odd & 23 & 18\,136\,691 & even\\
10 & 1\,277 & even & 24 & 24\,703\,847 & even\\
11 & 4\,885 & even & 25 & 55\,831\,137 & even\\
12 & 9\,361 & even & 26 & 88\,613\,778 & even\\
13 & 28\,734 & odd & & &
\end{tabular}
\caption{The counter $s_n$ and parity of $\nu(L_n)$ for the distinguished seed $L_0=10$, $0\le n\le 26$.}
\label{b2:tab:traj}
\end{table}

\begin{table}[H]\centering
\begin{tabular}{lr}
quantity & value\\
\hline
seeds of length $\le 14$ (beginning with $1$) & $16383$\\
halting / non-halting (presumed) & $4859$ ($29.7\%$) / $11524$\\
halt at step $0$ / $1$\,; max halt step observed & $1917$ / $2146$\,; $9$\\
counter at halt $=0$ / $=1$ & $1464$ / $3395$\\
first halts ($N\ge 2$) violating $s_{N-2}\le 5$ & $0$ of $796$\\
non-halting (len $\le 11$) violating two-step monotonicity & $\ge 586$ of $1136$ (depth-dependent)$^*$\\
non-halting (len $\le 11$) with $g(L_0)\ge 1$ & $217$ of $1136$\\
\end{tabular}
\caption{Classification of short seeds and structural statistics.\\
$^*${\small The count of monotonicity-violators is non-decreasing as the simulation depth grows. The figure $586$ is the count within $80$ simulation steps per seed; within $150$ steps the count is $\ge 659$.}}
\label{b2:tab:stats}
\end{table}

\begin{table}[H]\centering
\begin{tabular}{l|cccccccccccccc}
$\ell$ & 1 & 2 & 3 & 4 & 5 & 6 & 7 & 8 & 9 & 10 & 11 & 12 & 13 & 14\\
halt frac. & 1.00 & .50 & 1.00 & .75 & .69 & .69 & .61 & .57 & .50 & .44 & .39 & .34 & .29 & .25
\end{tabular}
\caption{Halting fraction by seed length $\ell$ (seeds beginning with $1$).}
\label{b2:tab:bylength}
\end{table}

\subsubsection{Status}\label{b2:sec:status}
\subsubsection*{Proved} $F$ strictly lengthens and outputs end in $0$ (Lemmas~\ref{b2:lem:length},~\ref{b2:lem:end0}); $F$ has no cycles and every orbit is a halt-or-diverge simple path (Theorem~\ref{b2:thm:nocycles}); odd parity forces survival and doubling (Lemma~\ref{b2:lem:odd}); the counter formula and halting dichotomy (Proposition~\ref{b2:prop:counter}, Corollary~\ref{b2:cor:halting}); the reshaping identities (\S\ref{b2:sec:struct}); the step-$0$ halt characterization (Proposition~\ref{b2:prop:step0}); the prefix-alternation invariant with seed-controlled parity (Proposition~\ref{b2:prop:prefix}) and the all-even exclusion (Proposition~\ref{b2:prop:alleven}, Remarks~\ref{b2:rem:rem18},~\ref{b2:rem:TM}); the first-halt bound $s_{N-2}\le 5$ (Proposition~\ref{b2:prop:firsthalt}); the even-step floor (Theorem~\ref{b2:thm:floor}); the counter-inert tail and body localization (Lemma~\ref{b2:lem:inert}, Observation~\ref{b2:obs:decomp}); the $\nu$-recurrence (Proposition~\ref{b2:prop:nu}); the conjugacies (\S\ref{b2:sec:equiv}); invertibility and the first-halt shape (Propositions~\ref{b2:prop:inv},~\ref{b2:prop:halt-shape}, Remark~\ref{b2:rem:rem31}).

\subsubsection*{Reduced} For a seed that excludes the all-even profile, non-halting is equivalent to $(s_n)$ never re-entering $\{2,3,4,5\}$ after its finite transient, for which two-step monotonicity is sufficient (Corollary~\ref{b2:cor:remaining}). Monotonicity is settled at odd steps (doubling) and at even-after-even steps (Theorem~\ref{b2:thm:floor}); the residue is the even-after-odd step (Conjecture~\ref{b2:conj:wall}), localized further onto the alignment-filtered size surplus (Observation~\ref{b2:obs:decomp}).

\subsubsection*{Open} Conjecture~\ref{b2:conj:wall} (hence Theorem~\ref{b2:thm:central}) for any non-halting seed, including $10$; the wall is intrinsic to $F$, every infinite orbit meets it, and the generic seed grazes the threshold repeatedly, making it harder than $10$ rather than easier.

\subsubsection{Conjectures}\label{b2:sec:conj}
\begin{conjecture}[General non-halting reduction]\label{b2:conj:reduction}
A seed is non-halting iff no iterate of its orbit wears a fragile profile ($s_n\ge 2$ for all $n$); Conjecture~\ref{b2:conj:wall} is a clean sufficient condition that the counter not re-enter $\{2,3,4,5\}$ after a finite transient.
\end{conjecture}

\begin{conjecture}[No eventually all-odd orbit]\label{b2:conj:noallodd}
No seed has an orbit that is eventually of all-odd parity; equivalently, every infinite orbit has infinitely many even-parity steps.
\end{conjecture}

\begin{conjecture}[Unbounded graze multiplicity]\label{b2:conj:graze}
$\sup_{L_0\in\mathcal H^c} g(L_0)=\infty$.
\end{conjecture}

\begin{conjecture}[Global injectivity]\label{b2:conj:inj}
$F$ is injective on the strings beginning with $1$ and ending with $0$.
\end{conjecture}

\begin{conjecture}[Vanishing halting density]\label{b2:conj:density}
$|\mathcal H_\ell|/2^{\ell-1}\to 0$ as $\ell\to\infty$, where $\mathcal H_\ell$ are the halting seeds of length $\ell$.
\end{conjecture}

\subsubsection{Open questions}\label{b2:sec:open}
\begin{question}[Decidability of halting]\label{b2:q:dec}
Is $\mathcal H$ decidable? It is recursively enumerable; orbit-membership is decidable (Proposition~\ref{b2:prop:inv}). Is non-halting recognizable via a computable certificate derived from Conjecture~\ref{b2:conj:wall}, or via FAR-type automata as in Remark~\ref{b2:rem:TM} applied uniformly?
\end{question}

\begin{question}[A provably non-halting seed]\label{b2:q:provable}
Is there a seed with an eventually all-odd orbit, hence provably non-halting? If none (Conjecture~\ref{b2:conj:noallodd}), what obstructs the parity dynamics of Proposition~\ref{b2:prop:nu} from sustaining odd parity, given that $s_n\bmod 2$ has no scalar recurrence?
\end{question}

\begin{question}[Density and language complexity of $\mathcal H$]\label{b2:q:density}
Does $|\mathcal H_\ell|/2^{\ell-1}$ converge, and to what (the observed fractions decrease past $\ell=3$, Table~\ref{b2:tab:bylength})? Is $\mathcal H$ regular, context-free, or context-sensitive? The step-$0$ halts are regular; the later halts are not obviously so.
\end{question}

\begin{question}[Uniform descent]\label{b2:q:descent}
Is there a Lyapunov function on the computable data, valid for all seeds, that survives the threshold grazes of Observation~\ref{b2:obs:graze}? None was found for the orbit of $10$, and the repeated grazes make the general task strictly harder.
\end{question}

\begin{question}[Arithmetic of the conjugate]\label{b2:q:arith}
Can the map $T$ of Proposition~\ref{b2:prop:T} and its alternating bit-sum kernel $2S$ (Lemma~\ref{b2:lem:S}) be analysed intrinsically on $\mathbb Z^+$, by analytic or $2$-adic methods?
\end{question}

\begin{question}[The Turing-machine bridge]\label{b2:q:TM}
The all-even exclusion on the orbit of $10$ is certified by FAR (Remark~\ref{b2:rem:TM}). Can the $s=1$ profile---the genuine residue---be encoded as a Turing machine whose non-halting is provable by FAR or a related decider, and if so, would such a certificate be uniform over seeds? A negative answer would locate the difficulty precisely at the boundary of automatic methods.
\end{question}

\subsubsection{Algorithms}\label{b2:app:alg}
The bit engine is a literal transcription of Definition~\ref{b2:def:F} on a character array; it is ground truth while the string fits in memory. The gap engine represents a string by the gaps $(e;t)$ between consecutive $1$'s and the trailing-zero count $t$, advancing in time and space $O(\#\text{runs})$. Three points are essential to a faithful gap engine: the additive seed $s_{\mathrm{init}}=-1$ of the odd pad; the trailing-zero count $t'=t+3$ at the unique $s=2$ even step; and the gap from the body's last $1$ into the appended tail, which equals (body trailing zeros)$+1$, with body trailing zeros $t+3$ after an even pad and $1$ after an odd pad. With these handled, the two engines agree on every common output.

\subsubsection{Computational evidence}\label{b2:app:evidence}
For the orbit of $10$, the counter sequence $s_0,\dots,s_{26}$ is
\[ 3,5,2,8,31,66,158,392,938,2305,1277,4885,9361,28734,67339,29083,\]\[218336,498138,1332397,659539,4146749,3732091,8693523,18136691,24703847,55831137,88613778, \]
with parity string of $\nu(L_n)$
\[ \mathrm{o\ o\ E\ E\ o\ o\ o\ o\ o\ o\ E\ E\ E\ o\ o\ E\ o\ o\ o\ E\ o\ E\ E\ E\ E\ E\ E}, \]
reproduced exactly by the $\nu$-recurrence (Proposition~\ref{b2:prop:nu}) iterated from $\nu(L_0)=1$, which also predicts $\nu(L_{26})=222{,}569{,}758$. The odd-run counts $\omega(L_n)$ for $0\le n\le 15$ are $1,1,4,6,1,11,47,105,235,589,1420,3704,1960,7549,13983,44910$; the small values occur only at odd-parity iterates, where survival is guaranteed by Lemma~\ref{b2:lem:odd}. The backward decision procedure confirms that all $4565$ halting strings of length $\le 14$ are absent from the orbit of $10$. For general seeds, the statistics of Tables~\ref{b2:tab:stats}--\ref{b2:tab:bylength} were obtained from the same engines; the first-halt counterexample to predecessor pinning is $L_0=110011111110$ ($N=8$, $s_{N-1}=102$, $s_N=1$), and the all-even output counterexample is $F(1)=11110$.

\begin{table}[H]\centering
\begin{tabular}{rrrr|l}
$n$ & $s_n$ & body $\Sigma$ & tail $\Sigma$ & nonzero contributing-gap multiset\\
\hline
10 & 1\,277 & 1\,277 & 0 & $\{1:551,\,2:230,\,3:87,\,5:1\}$\\
15 & 29\,083 & 29\,083 & 0 & $\{1:17\,594,\,2:5613,\,3:1,\,4:65\}$\\
19 & 659\,539 & 659\,539 & 0 & $\{1:273\,509,\,2:147\,958,\,3:17\,266,\,4:7769,\,5:810,\,10:319\}$
\end{tabular}
\caption{Body/tail decomposition of the counter at the open (odd-predecessor) even steps. The tail $\Sigma$ is identically $0$ (Lemma~\ref{b2:lem:inert}). For $n=15$ the count floor (total multiplicity) is $23{,}273$, the bound $\lceil\omega/2\rceil=22{,}455$, and the size surplus $5613\cdot 1+1\cdot 2+65\cdot 3=5810$, so $s_{15}=23{,}273+5810=29{,}083$.}
\label{b2:tab:decomp}
\end{table}

\subsubsection{Computational programs}\label{b2:app:programs}
The two implementations of Definition~\ref{b2:def:F}---the bit-level simulator and the gap-vector engine of Appendix~\ref{b2:app:alg}---together with the verification and statistics-gathering scripts, are available at
\begin{center}
\url{https://github.com/RobinCodes/Projects/tree/main/Collatz-type%20problem/Collatz-type%20problem%20v6}.
\end{center}

\subsubsection{Note on the use of artificial intelligence}\label{b2:app:ai}
This work was prepared with substantial assistance from an artificial-intelligence system, and most of the results presented here were created with such assistance: the map of Definition~\ref{b2:def:F} was implemented in two independent ways (Appendix~\ref{b2:app:alg}) and cross-validated on every commonly feasible output, the computational record was extended, and the structural results were formulated and proved with machine assistance. An independent re-derivation, working only from the statements of this paper, reconstructed the map from scratch, re-implemented it in two further engines (a literal bit-level engine and a numpy gap-vector engine), and reproduced every computational and structural claim, including: the counter sequence and parity through $n=26$ with $s_{26}=88{,}613{,}778$ and $\nu(L_{26})=222{,}569{,}758$; the integer-conjugate values $N_0,\dots,N_3=2,998,45\,868\,646,213\,192\,976$ and the conjugacy $T(N_n)=N_{n+1}$ through $n=7$; the structural lemmas of \S\S\ref{b2:sec:rl}--\ref{b2:sec:struct}--\ref{b2:sec:step0} by brute force on all strings of length $\le 12$ beginning with $1$, with zero violations; the body/tail decomposition multisets of Table~\ref{b2:tab:decomp}; the general-seed statistics of Tables~\ref{b2:tab:stats}--\ref{b2:tab:bylength}; and the count-floor identification of Observation~\ref{b2:obs:decomp}, including the depth-dependent monotonicity-violation figure of Observation~\ref{b2:obs:mono} and Table~\ref{b2:tab:stats}. This re-derivation reached the same conclusion regarding the central inequality---that it lies beyond elementary methods.

We do not specify the particular artificial-intelligence system or systems used; readers seeking further detail about the methods are invited to contact the authors: Robin Rovenszky (\texttt{robincodes} on Discord) and notxxdog (\texttt{.notxxdog} on Discord).

The Turing-machine encoding of Remark~\ref{b2:rem:TM} and its Finite Automata Reduction certificate are due to the authors and the Busy Beaver Challenge collaboration \cite{b2:bbc}; the simulation check reported there ($10^7$ steps without halting) was reproduced independently. The equivalence between the machine's halting and the $s=0$ event is, by construction, an assertion about the encoding; the FAR certificate is the formal object of record.

No result here is accepted on the basis of machine assertion. Every definition, lemma, and proof is stated so as to be checkable by hand, and every computational claim is reproducible by re-running the procedures of Appendix~\ref{b2:app:alg} and the programs of Appendix~\ref{b2:app:programs}. The status of the central open inequality (Conjecture~\ref{b2:conj:wall}) is reported exactly as it stands.

\begin{thebibliography}{9}
\bibitem{b2:bbc} The Busy Beaver Challenge, \emph{The determination of $BB(5)$ and non-halting deciders, including Finite Automata Reduction}, \url{https://bbchallenge.org}.
\end{thebibliography}


\section{Division problem with powers of two}

\subsection{Paper 1, Special Cases ($S_{-1} \; \And \; S_{-3}$)}

\begin{abstract}
We study which integers $m$ satisfy $m\mid 2^{m-1}+3$, equivalently which $n$ satisfy
$n+1\mid 2^{n}+3$ (OEIS A245728). Only two are known---$61\cdot 228806497$ and
$67\cdot 44210291368986343$, each a product of two primes---and none lies below $10^{15}$. Our
central result ties the question to the prime factorisation of $m$: a squarefree
$m=p_{1}\cdots p_{\omega}$ is a solution if and only if, for every $i$, $-3$ is a power of $2$
modulo $p_{i}$ and $m\equiv k_{i}+1\pmod{d_{i}}$, where $d_{i}=\ord_{p_{i}}(2)$ and
$2^{k_{i}}\equiv-3\pmod{p_{i}}$. Equivalently, the largest prime factor of $m$ must divide
$2^{P-1}+3$ with $P$ the product of the remaining primes, so locating solutions reduces to
factoring a single, rapidly growing integer. These local conditions exclude a positive
proportion of primes as factors and cannot be sharpened by prime-power moduli; from them we show
that every solution is squarefree at each prime below $1093$, and that the only two-prime
solutions with smaller prime at most $499$ are the two above. Heuristics suggest that this set is infinite (\S\ref{d1:sec:finite})---though extraordinarily
sparse, dominated by many-prime $m$ beyond any feasible factorisation, which is precisely why
only two-prime examples have been found. By contrast, $m\mid 2^{m-1}+1$ has no solutions, by a
short $2$-adic argument.
\end{abstract}


\subsubsection{Introduction}

For an integer $a$, write $S_{a}=\{\,n\ge 1:(n+1)\mid 2^{n}-a\,\}$. Setting $m=n+1$, membership
$n\in S_{a}$ is equivalent to
\begin{equation}\label{d1:eq:main}
m\mid 2^{m-1}-a,\qquad m\ge 2.
\end{equation}
This note concerns $a=-1$ and $a=-3$. We prove $S_{-1}=\varnothing$ and assemble, both
structurally and computationally, what is presently known about $S_{-3}$.

\paragraph{Notation.}
For an odd prime $p$ with $p\nmid 2$, write $d_{p}=\ord_{p}(2)$, the multiplicative order of $2$
in $\F_{p}^{\times}$. For $a\in\Z$ with $\gcd(a,p)=1$ and $a\in\gen{2}\subseteq\F_{p}^{\times}$,
write $\log_{2}a$ for the smallest non-negative integer $k$ with $2^{k}\equiv a\pmod{p}$. We use
$\widetilde{\log_{2}}a$ analogously for the discrete logarithm in $(\Z/p^{2})^{\times}$. The
$2$-adic valuation of an integer $N$ is $v_{2}(N)$.

\subsubsection{The case \texorpdfstring{$a=-1$}{a=-1}}

\begin{theorem}\label{d1:thm:minus1}
There is no integer $m>1$ with $m\mid 2^{m-1}+1$. Equivalently, $S_{-1}=\varnothing$.
\end{theorem}

\begin{proof}
Suppose such an $m$ exists. The integer $2^{m-1}+1$ is odd, so $m$ is odd. Fix any prime
$p\mid m$ and set $d=d_{p}$. From $p\mid 2^{m-1}+1$,
\[
2^{m-1}\equiv-1\pmod{p},\qquad 2^{2(m-1)}\equiv 1\pmod{p}.
\]
Hence $d\mid 2(m-1)$ but $d\nmid m-1$ (since $p$ is odd, $-1\not\equiv 1\pmod{p}$). It follows
that $d/\gcd(d,m-1)=2$, so $\gcd(d,m-1)=d/2$. Comparing $2$-adic valuations,
$v_{2}(d)-1=v_{2}(\gcd(d,m-1))\le v_{2}(m-1)$ on the one hand, while $d\nmid m-1$ together with
$d\mid 2(m-1)$ forces $v_{2}(d)>v_{2}(m-1)$ on the other. Therefore
\begin{equation}\label{d1:eq:val}
v_{2}(d)=v_{2}(m-1)+1.
\end{equation}
Now $d\mid p-1$, so $v_{2}(p-1)\ge v_{2}(d)=v_{2}(m-1)+1$. Setting $\nu=v_{2}(m-1)+1$, we have
shown $p\equiv 1\pmod{2^{\nu}}$ for every prime $p\mid m$. Since $m$ is a product of primes in
$1+2^{\nu}\Z$, we get $m\equiv 1\pmod{2^{\nu}}$, i.e.\ $2^{\nu}\mid m-1$. But
$\nu=v_{2}(m-1)+1$, contradicting the definition of $v_{2}(m-1)$.
\end{proof}

\begin{remark}\label{d1:rem:roots}
The proof gives slightly more: for any $m\ge 2$ and any odd $a$ with $\gcd(a,m)=1$ and
$a\equiv-1\pmod{p}$ for every $p\mid m$, the same $2$-adic argument applies. The obstruction is
genuinely about square roots of unity, and this is exactly why it does not transfer to $a=-3$:
$-3$ is not a root of unity modulo a typical prime, so no global multiplicative collapse occurs.
\end{remark}

\subsubsection{The case \texorpdfstring{$a=-3$}{a=-3}: elementary restrictions}

We turn to \eqref{d1:eq:main} with $a=-3$, that is,
\begin{equation}\label{d1:eq:star}\tag{$\star$}
m\mid 2^{m-1}+3.
\end{equation}

\begin{proposition}\label{d1:prop:elem}
Let $m\ge 2$ satisfy \eqref{d1:eq:star}. Then:
\begin{enumerate}[label=\textup{(\roman*)}]
\item $m$ is odd, equivalently $n=m-1$ is even;
\item $3\nmid m$;
\item $m$ is composite.
\end{enumerate}
\end{proposition}

\begin{proof}
(i) $2^{m-1}+3$ is odd, so $m$ is odd. (ii) If $3\mid m$ then $3\mid 2^{m-1}+3$, hence
$3\mid 2^{m-1}$, impossible. (iii) If $m$ is prime, Fermat's little theorem gives
$2^{m-1}\equiv 1\pmod m$, so $m\mid 4$, forcing $m\in\{2\}$, contradicting (i).
\end{proof}

\subsubsection{The known solutions}

The integers $n$ satisfying $(n+1)\mid 2^{n}+3$ form OEIS sequence A245728~\cite{d1:oeis}. Only two
terms are known:
\[
m_{1}=n_{1}+1=13957196317=61\cdot 228806497,
\]
\[
m_{2}=n_{2}+1=2962089521722084981=67\cdot 44210291368986343,
\]
and no further $n$ exists with $n+1<10^{15}$. Both factorisations are products of two distinct
primes; in each case the smaller prime is small ($61$ or $67$) and the cofactor is large. Both
were verified directly to satisfy \eqref{d1:eq:star}, and all four factors are prime. The second
solution exceeds $10^{18}$ and was found structurally rather than by enumeration; see
Section~\ref{d1:sec:comp}.

\subsubsection{Admissibility: a sieve on prime factors}\label{d1:sec:sieve}

The next lemma gives a strong necessary condition that every prime divisor of $m$ must satisfy;
it underlies all subsequent analysis.

\begin{lemma}[Local condition at a prime divisor]\label{d1:lem:local}
Let $m\ge 5$ satisfy \eqref{d1:eq:star}, let $p$ be a prime dividing $m$, and write $m=p^{a}t$ with
$\gcd(p,t)=1$. Then $p\notin\{2,3\}$, the element $-3$ lies in $\gen{2}\subseteq\F_{p}^{\times}$,
and, writing $d=d_{p}$ and $k=\log_{2}(-3)\bmod p$,
\begin{equation}\label{d1:eq:loc}
t\equiv k+1\pmod{d}.
\end{equation}
\end{lemma}

\begin{proof}
By Proposition~\ref{d1:prop:elem}, $p$ is odd and $p\ne 3$, so $\gcd(p,6)=1$. Since $d\mid p-1$,
$p\equiv 1\pmod d$, hence $p^{a}\equiv 1\pmod d$ and $m\equiv t\pmod d$. Therefore
$m-1\equiv t-1\pmod d$ and $2^{m-1}\equiv 2^{t-1}\pmod p$. Since $p\mid 2^{m-1}+3$, this gives
$2^{t-1}\equiv-3\pmod p$, which both shows $-3\in\gen{2}$ and yields \eqref{d1:eq:loc}.
\end{proof}

\begin{definition}\label{d1:def:adm}
A prime $p\notin\{2,3\}$ is \emph{admissible} if $-3\in\gen{2}\subseteq\F_{p}^{\times}$ and,
writing $d=d_{p}$ and $k=\log_{2}(-3)$:
\begin{enumerate}[label=\textup{(\Alph*)}]
\item if $d$ is even, then $k$ is even;
\item if $3\mid d$, then $k\not\equiv 2\pmod 3$.
\end{enumerate}
\end{definition}

\begin{theorem}\label{d1:thm:adm}
Every prime divisor of any $m$ satisfying \eqref{d1:eq:star} is admissible.
\end{theorem}

\begin{proof}
Fix $p\mid m$ and use the notation of Lemma~\ref{d1:lem:local}, with $t=m/p^{v_{p}(m)}$. Since $m$
is odd, $t$ is odd; since $3\nmid m$, $3\nmid t$. Reduce \eqref{d1:eq:loc} modulo $2$ and modulo
$3$. If $d$ is even, then $t\bmod d$ has the same parity as $t$, namely odd, so $k+1$ is odd,
i.e.\ $k$ is even. If $3\mid d$, then $t\bmod d$ has the same residue modulo $3$ as $t$, which is
$1$ or $2$ but not $0$, so $k+1\not\equiv 0\pmod 3$, i.e.\ $k\not\equiv 2\pmod 3$.
\end{proof}

\paragraph{Numerical effect of the sieve.}
Among the primes $p<200$ for which $-3\in\gen{2}$ in $\F_{p}^{\times}$ (a necessary condition
without which $p$ cannot divide any $m$), Theorem~\ref{d1:thm:adm} eliminates substantially more
than half. Table~\ref{d1:tab:sieve} records the data; admissible primes are marked $\checkmark$.
The admissible primes below $200$ are
\[
13,\ 19,\ 61,\ 67,\ 79,\ 103,\ 193,\ 199.
\]
The two known solutions use the admissible primes $61$ and $67$ together with admissible large
cofactors (one verifies directly that $228806497$ and $44210291368986343$ are admissible).
Condition (A) eliminates roughly half of the eligible primes (those with $d$ even and $k$ odd,
which by equidistribution is half), and condition (B) eliminates a further third of those with
$3\mid d$. Empirically, over primes below $2000$, about $56\%$ are $-3$-eligible and the sieve
removes about $61.5\%$ of those, leaving roughly $21.6\%$ of all primes admissible. Thus
admissible primes form a sparse but \emph{positive-proportion} subsequence (conjecturally of
positive density, by an Artin-type heuristic); in particular the supply of candidate small
factors is infinite, which already signals that no finite search can be conclusive.

\begin{table}[H]
\centering
\small
\begin{tabular}{rrrccc}
\toprule
$p$ & $d_{p}$ & $k$ & (A) $k$ even if $2\mid d$ & (B) $k\not\equiv 2\ (3)$ if $3\mid d$ & adm. \\
\midrule
5   & 4  & 1  & fail & ---  & $\times$ \\
7   & 3  & 2  & ---  & fail & $\times$ \\
11  & 10 & 3  & fail & ---  & $\times$ \\
13  & 12 & 10 & ok   & ok   & $\checkmark$ \\
19  & 18 & 4  & ok   & ok   & $\checkmark$ \\
29  & 28 & 19 & fail & ---  & $\times$ \\
37  & 36 & 8  & ok   & fail & $\times$ \\
53  & 52 & 43 & fail & ---  & $\times$ \\
59  & 58 & 21 & fail & ---  & $\times$ \\
61  & 60 & 36 & ok   & ok   & $\checkmark$ \\
67  & 66 & 6  & ok   & ok   & $\checkmark$ \\
79  & 39 & 10 & ---  & ok   & $\checkmark$ \\
83  & 82 & 31 & fail & ---  & $\times$ \\
97  & 48 & 43 & fail & ok   & $\times$ \\
101 & 100& 19 & fail & ---  & $\times$ \\
103 & 51 & 9  & ---  & ok   & $\checkmark$ \\
107 & 106& 17 & fail & ---  & $\times$ \\
131 & 130& 7  & fail & ---  & $\times$ \\
139 & 138& 110& ok   & fail & $\times$ \\
149 & 148& 13 & fail & ---  & $\times$ \\
163 & 162& 20 & ok   & fail & $\times$ \\
173 & 172& 113& fail & ---  & $\times$ \\
179 & 178& 19 & fail & ---  & $\times$ \\
181 & 180& 146& ok   & fail & $\times$ \\
193 & 96 & 90 & ok   & ok   & $\checkmark$ \\
197 & 196& 83 & fail & ---  & $\times$ \\
199 & 99 & 85 & ---  & ok   & $\checkmark$ \\
\bottomrule
\end{tabular}
\caption{Primes $p<200$ with $-3\in\gen{2}\pmod p$. Column (A) is vacuous when $d_{p}$ is odd;
column (B) is vacuous when $3\nmid d_{p}$.}
\label{d1:tab:sieve}
\end{table}

\subsubsection{Interlocking congruences}\label{d1:sec:interlock}

The local condition \eqref{d1:eq:loc} globalises by the Chinese Remainder Theorem.

\begin{proposition}\label{d1:prop:crt}
Let $m$ satisfy \eqref{d1:eq:star} with prime factorisation $m=\prod_{i=1}^{r}p_{i}^{a_{i}}$. For
each $i$, set $d_{i}=d_{p_{i}}$ and $k_{i}=\log_{2}(-3)\bmod p_{i}$, and put $t_{i}=m/p_{i}^{a_{i}}$.
Then $t_{i}\equiv k_{i}+1\pmod{d_{i}}$ for every $i$, and consequently
\[
m\equiv p_{i}^{a_{i}}(k_{i}+1)\pmod{p_{i}^{a_{i}}d_{i}}
\]
for every $i$. By CRT, $m$ lies in a single residue class modulo
$\lcm_{i}\!\big(p_{i}^{a_{i}}d_{i}\big)$.
\end{proposition}

\begin{proof}
The first statement is Lemma~\ref{d1:lem:local}. For the second, $m=p_{i}^{a_{i}}t_{i}$ and
$t_{i}\equiv k_{i}+1\pmod{d_{i}}$, so $m\equiv p_{i}^{a_{i}}(k_{i}+1)\pmod{p_{i}^{a_{i}}d_{i}}$.
\end{proof}

For both known solutions one verifies Proposition~\ref{d1:prop:crt} explicitly:
\[
m_{1}=61\cdot 228806497:\quad 228806497 \bmod 60 = 37 = k_{61}+1,
\]
\[
m_{2}=67\cdot 44210291368986343:\quad 44210291368986343 \bmod 66 = 7 = k_{67}+1.
\]

\subsubsection{Two-prime structure}\label{d1:sec:twoprime}

Both known $m$ are products of two distinct primes. Specialising Proposition~\ref{d1:prop:crt}
gives a tight description of all such solutions, and the next remark shows that one half of the
description is automatic.

\begin{corollary}\label{d1:cor:twoprime}
Suppose $m=pq$ with primes $p<q$, $p,q\ne 2,3$. Then $m$ satisfies \eqref{d1:eq:star} if and only
if both
\begin{equation}\label{d1:eq:pair}
q\equiv k_{p}+1\pmod{d_{p}},\qquad p\equiv k_{q}+1\pmod{d_{q}}
\end{equation}
hold, where $d_{p},d_{q}$ are the orders of $2$ and $k_{p},k_{q}$ the discrete logarithms of
$-3$ to base $2$ in $\F_{p}^{\times}$ and $\F_{q}^{\times}$ respectively. In particular $q$ is a
prime divisor of $2^{p-1}+3$ and $p$ is a prime divisor of $2^{q-1}+3$.
\end{corollary}

\begin{remark}[One-sidedness]\label{d1:rem:onesided}
The two conditions in \eqref{d1:eq:pair} are not independent once $q$ is drawn from the prime
factors of $N_{p}:=2^{p-1}+3$. Indeed, if $q\mid N_{p}$ then $2^{p-1}\equiv-3\pmod q$, and a
short computation with Fermat's little theorem shows
\[
2^{pq-1}\equiv (2^{p-1})^{q}\cdot 2^{q-1}\equiv(-3)^{q}\cdot 1\equiv -3\pmod q,
\]
so the congruence modulo $q$ holds automatically. Equivalently, the second condition
$p\equiv k_{q}+1\pmod{d_{q}}$ holds for \emph{every} prime factor $q$ of $N_{p}$. Hence
\[
\boxed{\,m=pq\ (p<q)\text{ solves }\eqref{d1:eq:star}\iff q\mid 2^{p-1}+3\ \text{ and }\ q\equiv k_{p}+1\pmod{d_{p}}.\,}
\]
This makes the two-prime problem for a fixed smaller prime $p$ entirely a question about the
prime factorisation of the single integer $N_{p}$, together with one residue test modulo
$d_{p}$ ($<p$).
\end{remark}

This is the basis of an efficient algorithm: for each admissible small prime $p$, factor
$N_{p}=2^{p-1}+3$, and for each prime factor $q>p$ test whether $q\equiv k_{p}+1\pmod{d_{p}}$.

\subsubsection{Solutions with more prime factors}\label{d1:sec:omega}

Corollary~\ref{d1:cor:twoprime} is the case $\omega(m)=2$. The same reduction works for any number
of prime factors; it both clarifies the structure of hypothetical longer solutions and explains
why only two-prime ones are known. Throughout, $m=p_{1}p_{2}\cdots p_{\omega}$ is squarefree with
$p_{1}<p_{2}<\cdots<p_{\omega}$ distinct primes, $d_{i}=d_{p_{i}}$ and $k_{i}=\log_{2}(-3)\bmod
p_{i}$. (By Section~\ref{d1:sec:sqfree} every solution is squarefree at all primes below $1093$, so
a repeated prime factor, if one ever occurs, must be $\ge 1093$; we return to this at the end.)

\begin{proposition}[Squarefree solutions of any length]\label{d1:prop:omega}
Let $m=p_{1}\cdots p_{\omega}$ be squarefree with each $p_{i}\notin\{2,3\}$. Then
$m\mid 2^{m-1}+3$ if and only if, for every $i$,
\[
-3\in\gen{2}\subseteq\F_{p_{i}}^{\times}\qquad\text{and}\qquad m\equiv k_{i}+1\pmod{d_{i}},
\]
the latter being equivalent to $\displaystyle\prod_{j\ne i}p_{j}\equiv k_{i}+1\pmod{d_{i}}$. In
particular every $p_{i}$ is admissible.
\end{proposition}

\begin{proof}
Since $m$ is squarefree, the CRT gives $m\mid 2^{m-1}+3$ iff $p_{i}\mid 2^{m-1}+3$ for every
$i$. For fixed $i$, $2^{m-1}\equiv-3\pmod{p_{i}}$ holds iff $-3\in\gen{2}$ in $\F_{p_{i}}^{\times}$
and $m-1\equiv k_{i}\pmod{d_{i}}$. As $d_{i}\mid p_{i}-1$ we have $p_{i}\equiv 1\pmod{d_{i}}$,
hence $m\equiv\prod_{j\ne i}p_{j}\pmod{d_{i}}$, giving the reformulation. Conditions (A),(B) of
Definition~\ref{d1:def:adm} follow as in Theorem~\ref{d1:thm:adm}.
\end{proof}

Thus a squarefree solution is exactly a set of admissible primes whose product lands
simultaneously in the prescribed class $k_{i}+1$ modulo each $d_{i}$. By CRT this confines $m$ to
one residue class modulo $\lcm_{i}d_{i}$ (Proposition~\ref{d1:prop:crt}); the difficulty is that the
unknown $m$ is forced to be the \emph{product} of the very primes defining the moduli. The
two-prime reduction of Remark~\ref{d1:rem:onesided} generalises verbatim: the condition at the
largest prime is a divisibility, and the remaining ones are residue tests.

\begin{corollary}[Recursive criterion]\label{d1:cor:omega}
Let $p_{1}<\cdots<p_{\omega}$ be admissible primes, put $P=p_{1}\cdots p_{\omega-1}=m/p_{\omega}$,
and write $N_{P}=2^{P-1}+3$. Then $m=Pp_{\omega}$ solves \eqref{d1:eq:star} if and only if
\[
p_{\omega}\mid N_{P}\qquad\text{and}\qquad m\equiv k_{i}+1\pmod{d_{i}}\quad(1\le i\le\omega-1).
\]
The congruence $m\equiv k_{\omega}+1\pmod{d_{\omega}}$ is automatic once $p_{\omega}\mid N_{P}$.
\end{corollary}

\begin{proof}
The congruence at $p_{\omega}$ reads $P\equiv k_{\omega}+1\pmod{d_{\omega}}$ (since $m\equiv P
\pmod{d_{\omega}}$), i.e.\ $2^{P-1}\equiv-3\pmod{p_{\omega}}$, i.e.\ $p_{\omega}\mid N_{P}$; this
holds for every prime factor of $N_{P}$. The remaining $\omega-1$ congruences are
Proposition~\ref{d1:prop:omega}.
\end{proof}

This is an algorithm: to find every $\omega$-prime solution extending a fixed tuple
$(p_{1},\dots,p_{\omega-1})$, factor the single integer $N_{P}=2^{P-1}+3$ and test each prime
factor $p_{\omega}>p_{\omega-1}$ against the $\omega-1$ residue conditions modulo
$d_{1},\dots,d_{\omega-1}$. Two points deserve emphasis. First, the smaller primes need
\emph{not} themselves form a solution: solutions do not nest, so one cannot grow a long solution
by extending a short one. Second, everything reduces to factoring $N_{P}$, and that is precisely
where the method fails.

\subsubsection{The computational barrier}
The integer $N_{P}=2^{P-1}+3$ has about $0.301\,P$ decimal digits, and $P=p_{1}\cdots
p_{\omega-1}$ is a product of admissible primes, hence grows super-exponentially with $\omega$.
Even the smallest admissible primes $13,19,61,67,\dots$, which minimise $P$, give the sizes of
Table~\ref{d1:tab:growth}.

\begin{table}[H]
\centering
\small
\begin{tabular}{cccl}
\toprule
$\omega$ & minimal $P=p_{1}\cdots p_{\omega-1}$ & digits of $N_{P}$ & feasibility \\
\midrule
$2$ & $13$ & $4$ & routine (broad search) \\
$3$ & $13\cdot 19=247$ & $75$ & at the edge (one tuple) \\
$4$ & $13\cdot 19\cdot 61=15067$ & $\approx 4{,}500$ & infeasible \\
$5$ & $13\cdot 19\cdot 61\cdot 67=1{,}009{,}489$ & $\approx 3\times 10^{5}$ & infeasible \\
\bottomrule
\end{tabular}
\caption{Smallest integer $N_{P}=2^{P-1}+3$ that must be factored to search for an
$\omega$-prime solution.}
\label{d1:tab:growth}
\end{table}

So two prime factors is the only regime open to a broad search, three is reachable for a single
tuple, and four or more is permanently beyond factorisation. Concretely, for three primes only
the smallest admissible pair $(p_{1},p_{2})=(13,19)$ keeps $N_{P}$ tractable:
\[
N_{247}=2^{246}+3=307\cdot 41047\cdot r,\qquad r\ \text{a 67-digit prime},
\]
and none of $13\cdot 19\cdot p_{3}$ with $p_{3}\in\{307,41047,r\}$ satisfies the two residue
conditions, so there is no three-prime solution with smaller pair $(13,19)$. Every other pair
already pushes $N_{P}$ past $200$ digits, so a complete three-prime search---let alone
$\omega\ge 4$---is not currently possible.

\subsubsection{How many solutions, and how large?}
Because the multi-prime regime is computationally invisible, the only guide is heuristic.
Modelling $2^{m-1}+3\bmod m$ as equidistributed makes a squarefree $m$ with all prime factors
admissible a solution with ``probability'' about $1/m$; when $d_{i}\approx p_{i}$ this agrees
with the structural estimate $\prod_{i}1/d_{i}$. Summing over candidates with exactly $\omega$
prime factors gives, for the number of $\omega$-prime solutions up to $X$,
\[
E_{\omega}(X)\ \approx\ \frac{1}{\omega!}\Bigl(\sum_{\substack{p\le X\\\text{admissible}}}
\frac{1}{p}\Bigr)^{\!\omega}\ \approx\ \frac{(\delta\,\log\log X)^{\omega}}{\omega!},
\]
where $\delta$ is the (logarithmic) density of admissible primes, empirically
$\delta\approx 0.15$. Hence
\[
\sum_{\omega\ge 2}E_{\omega}(X)\ \approx\ e^{\delta\log\log X}-1-\delta\log\log X
\ =\ (\log X)^{\delta}-1-\delta\log\log X\ \xrightarrow[X\to\infty]{}\ \infty.
\]
Three conclusions follow, of which the third is the point. (i)~The total diverges, so the
heuristic predicts \emph{infinitely many} solutions---but only as a minuscule power
$(\log X)^{\delta}$ of $\log X$, so they are astronomically sparse. (ii)~The estimate is of the
right order: at the scale of the known terms ($X\approx m_{2}$) it gives $\sum_{\omega\ge 2}
E_{\omega}$ of order $10^{-1}$, so finding two solutions is consistent with the model to within
its (considerable) crudeness; the structural refinement $\prod_{i}1/d_{i}$, which exceeds $1/m$
when some order $d_{i}$ is small, raises the estimate toward the observed count. (iii)~At every
reachable scale $\delta\log\log X<1$, so $E_{\omega}$ \emph{decreases} with $\omega$: two-prime
solutions are heuristically the most common, and a solution with three or more prime factors is
not expected until $X$ is of the order $\exp\!\bigl(\exp(c/\delta)\bigr)$---a number with
thousands of digits or more. The two known solutions are two-prime both because two-prime
solutions genuinely dominate at small heights and because, even were a longer solution to lie in
range, exhibiting it would require factoring an $N_{P}$ of thousands of digits. Longer solutions
are thus doubly hidden: too rare to meet in any searchable range, and walled off by
factorisation when they do occur.

\subsubsection{Non-squarefree \texorpdfstring{$m$}{m}}
If $p^{a}\,\|\,m$ with $a\ge 2$, Proposition~\ref{d1:prop:omega} must be read modulo $p^{a}$: the
condition at $p$ becomes $2^{m-1}\equiv-3\pmod{p^{a}}$, governed by $\ord_{p^{a}}(2)$ and the
discrete logarithm of $-3$ there, in the manner of Lemma~\ref{d1:lem:sqfree}. Corollary
\ref{d1:cor:sqfree} forbids $a\ge 2$ for every prime $p<1093$, so a non-squarefree solution would
have to involve a repeated prime $\ge 1093$; none is known, and the heuristic above, dominated by
squarefree $m$, is unaffected.

\subsubsection{Squarefreeness at small primes}\label{d1:sec:sqfree}

Both known solutions are squarefree. We now prove that this is forced at every non-Wieferich
prime, modulo a single discrete-logarithm coincidence that we then verify never occurs below
$1093$. Recall that $p$ is a \emph{Wieferich prime} (base $2$) if $2^{p-1}\equiv 1\pmod{p^{2}}$;
the only such primes below $4\times 10^{3}$ are $1093$ and $3511$. For $p$ non-Wieferich one has
$\ord_{p^{2}}(2)=p\,d_{p}$ (if $2^{d_{p}}\equiv 1\pmod{p^{2}}$ then $2^{p-1}\equiv 1\pmod{p^{2}}$,
since $d_{p}\mid p-1$; the contrapositive gives the order $p\,d_{p}$).

\begin{lemma}[Squarefreeness at a non-Wieferich prime]\label{d1:lem:sqfree}
Let $m$ satisfy \eqref{d1:eq:star} and let $p$ be a prime with $p^{2}\mid m$. Suppose $p$ is not a
Wieferich prime base $2$, so $D:=\ord_{p^{2}}(2)=p\,d_{p}$. Then $-3\in\gen{2}\subseteq
(\Z/p^{2})^{\times}$ and, writing $\widetilde{k}_{p}:=\widetilde{\log_{2}}(-3)$ for its discrete
logarithm modulo $p^{2}$,
\[
\widetilde{k}_{p}\equiv -1\pmod{p}.
\]
\end{lemma}

\begin{proof}
From $p^{2}\mid m\mid 2^{m-1}+3$ we get $2^{m-1}\equiv-3\pmod{p^{2}}$, so $-3\in\gen{2}$ modulo
$p^{2}$ and
\[
m-1\equiv\widetilde{k}_{p}\pmod{D}.
\]
Since $p$ is non-Wieferich, $D=p\,d_{p}$, so $p\mid D$ and we may reduce the previous congruence
modulo $p$:
\[
m-1\equiv\widetilde{k}_{p}\pmod{p}.
\]
But $p^{2}\mid m$ gives $m\equiv 0\pmod p$, hence $m-1\equiv-1\pmod p$. Therefore
$\widetilde{k}_{p}\equiv-1\pmod p$.
\end{proof}

\begin{remark}
The mechanism here differs from that of the sieve. Conditions (A),(B) of
Definition~\ref{d1:def:adm} arise from the cofactor $t=m/p^{a}$ being odd and coprime to $3$, which
constrains $t$ only modulo $2$ and $3$; the classes modulo $4,9,\dots$ remain free
(Section~\ref{d1:sec:nosharp}). In Lemma~\ref{d1:lem:sqfree} the modulus is the prime $p$ itself, and
the order $D=p\,d_{p}$ ties the residue $m\bmod p$, which is $0$ since $p^{2}\mid m$, to
$\widetilde{k}_{p}\bmod p$. This is a genuine constraint rather than a free residue class.
\end{remark}

\begin{corollary}\label{d1:cor:sqfree}
For every admissible prime $p<1093$ one has $\widetilde{k}_{p}\not\equiv-1\pmod p$. Consequently
no solution $m$ of \eqref{d1:eq:star} is divisible by $p^{2}$ for any prime $p<1093$; that is, every
$m\in S_{-3}$ is squarefree at all primes below $1093$.
\end{corollary}

\begin{proof}
If $p^{2}\mid m$ then $p\mid m$, so $p$ is admissible by Theorem~\ref{d1:thm:adm}, and every
admissible $p<1093$ is non-Wieferich. By Lemma~\ref{d1:lem:sqfree} (and noting that if $-3\notin
\gen{2}$ modulo $p^{2}$ then $p^{2}\nmid m$ outright), divisibility by $p^{2}$ would force
$\widetilde{k}_{p}\equiv-1\pmod p$. A direct computation of the discrete logarithm
$\widetilde{k}_{p}$ for each of the $39$ admissible primes $p<1093$ shows
$\widetilde{k}_{p}\not\equiv-1\pmod p$ in every case, a contradiction.
\end{proof}

Table~\ref{d1:tab:sqfree} records $\widetilde{k}_{p}\bmod p$ for the small admissible primes; the
value $-1\bmod p=p-1$ never appears.

\begin{table}[H]
\centering
\small
\begin{tabular}{rrrr}
\toprule
$p$ & $d_{p}$ & $\widetilde{k}_{p}\bmod p$ & $-1\bmod p$ \\
\midrule
13  & 12  & 7   & 12  \\
19  & 18  & 6   & 18  \\
61  & 60  & 19  & 60  \\
67  & 66  & 35  & 66  \\
79  & 39  & 7   & 78  \\
103 & 51  & 66  & 102 \\
193 & 96  & 7   & 192 \\
199 & 99  & 35  & 198 \\
211 & 210 & 186 & 210 \\
\bottomrule
\end{tabular}
\caption{For small admissible primes, the reduced discrete logarithm
$\widetilde{k}_{p}=\widetilde{\log_{2}}(-3)\bmod p^{2}$ taken modulo $p$, against $-1\bmod p$.
Equality (which would permit $p^{2}\mid m$) never occurs; verified for all admissible $p<1093$.}
\label{d1:tab:sqfree}
\end{table}

A full proof of squarefreeness would also require a uniform treatment of large primes and a
separate analysis of the Wieferich primes $1093$ and $3511$; the present argument settles only
the primes below $1093$.

\subsubsection{No prime-power sharpening of the sieve}\label{d1:sec:nosharp}

It is natural to ask whether replacing the moduli $2,3$ in conditions (A),(B) by higher prime
powers $4,8,\dots$ or $9,27,\dots$ yields additional exclusions. It does not.

\begin{proposition}\label{d1:prop:nosharp}
The only constraints on $t=m/p^{a}$ available in isolation are $t$ odd and $3\nmid t$. Reducing
the local congruence \eqref{d1:eq:loc} modulo any power $2^{e}$ or $3^{f}$ produces no exclusion
beyond conditions \textup{(A)} and \textup{(B)}.
\end{proposition}

\begin{proof}
Suppose $2^{e}\mid d$. Then \eqref{d1:eq:loc} forces $t\equiv k+1\pmod{2^{e}}$. The constraint
$t$ odd permits exactly the odd residues modulo $2^{e}$, and the forced value $k+1$ fails to be
odd precisely when $k+1$ is even, i.e.\ when $k$ is odd. This is condition (A) (which already
applies as soon as $2\mid d$). Higher $e$ add nothing, since oddness is a condition modulo $2$
only. Similarly, if $3^{f}\mid d$ then $t\equiv k+1\pmod{3^{f}}$, and $3\nmid t$ permits exactly
the residues coprime to $3$; the forced value $k+1$ fails to be coprime to $3$ precisely when
$3\mid k+1$, i.e.\ $k\equiv 2\pmod 3$, which is condition (B). Higher $f$ add nothing.
\end{proof}

We verified computationally that no admissible prime below $3000$ is removed by any mod-$4$,
mod-$8$, mod-$9$, or mod-$27$ refinement, consistent with Proposition~\ref{d1:prop:nosharp}. Any
genuine sharpening of the sieve must therefore come from \emph{interactions between distinct
prime factors}, i.e.\ from the interlocking congruences of Proposition~\ref{d1:prop:crt}, rather
than from a single local modulus.

\subsubsection{Computational evidence}\label{d1:sec:comp}

\subsubsection{Direct search}
The bound $m\le 10^{15}$ is the range exhaustively verified in OEIS A245728~\cite{d1:oeis}; the only
solution there is $m_{1}$. An independent brute-force test over all odd $m\le 10^{7}$ with
$3\nmid m$ reproduces no solution in that range, agreeing with the database.

\subsubsection{Two-prime search via factorisation of \texorpdfstring{$2^{p-1}+3$}{2^(p-1)+3}}
By Remark~\ref{d1:rem:onesided}, for each admissible smaller prime $p$ it suffices to factor
$N_{p}=2^{p-1}+3$ and test the prime factors $q>p$ against $q\equiv k_{p}+1\pmod{d_{p}}$. We
carried this out for every admissible $p\le 499$:
\begin{itemize}[nosep]
\item For $p\in\{13,19,61,67,79,103,193,199,211,271,307,373,463,499\}$ the number $N_{p}$ is
fully factored (by trial division, Pollard $\rho$, $p-1$, and ECM). The complete list of
two-prime solutions with smaller prime in this set is exactly $\{m_{1},m_{2}\}$, arising at
$p=61$ and $p=67$.
\item For $p\in\{367,439,487\}$ a composite cofactor of $109$, $132$, and $147$ decimal digits
respectively remains unfactored; in each case all prime factors found so far yield no solution,
and (each residual cofactor having been checked to be composite, not a single large prime) the
possibility of a further two-prime $m$ involving these $p$ is left open.
\end{itemize}

\begin{table}[H]
\centering
\small
\begin{tabular}{lll}
\toprule
smaller prime $p$ & $N_{p}=2^{p-1}+3$ status & solution \\
\midrule
$13,19$ & full & none \\
$61$ & full & $m_{1}=61\cdot 228806497$ \\
$67$ & full & $m_{2}=67\cdot 44210291368986343$ \\
$79,103,193,199,211,271,307,373,463,499$ & full & none \\
$367,439,487$ & partial ($109$--$147$-digit cofactor) & open \\
\bottomrule
\end{tabular}
\caption{Two-prime search over admissible smaller primes $p\le 499$. By Theorem~\ref{d1:thm:adm}
only admissible primes can occur as a smaller factor, so non-admissible primes require no
factorisation.}
\label{d1:tab:search}
\end{table}

A search for solutions with three or more prime factors is governed by the recursive criterion
of Section~\ref{d1:sec:omega}; only the single tuple $(13,19)$ is within reach, and it yields no
three-prime solution.

\subsubsection{Is \texorpdfstring{$S_{-3}$}{S(-3)} finite?}\label{d1:sec:finite}

This is the heart of ``fully solving'' the problem, and it is open. The available heuristics---a
structural estimate for two-prime solutions and a global density argument---point the same way,
toward infinitude; we give both and then say what a proof would require.

\paragraph{The two-prime estimate.}
For admissible $p$, the number $N_{p}=2^{p-1}+3$ has about $\log\log N_{p}\asymp\log p$ distinct
prime factors, and a given one lies in the class required modulo $d_{p}$ with ``probability''
about $1/d_{p}$. The expected number of two-prime solutions with smaller prime below $X$ is thus
of order $\sum_{p<X,\ \mathrm{adm}}(\log p)/d_{p}$, and this \emph{diverges}. Since $d_{p}\mid p-1$
we have $d_{p}<p$, hence $(\log p)/d_{p}>(\log p)/p$; assuming, as the data strongly support, that
the admissible primes have positive density $\delta$, Mertens' first theorem gives
\[
\sum_{\substack{p<X\\\text{admissible}}}\frac{\log p}{d_{p}}\ \ge\
\sum_{\substack{p<X\\\text{admissible}}}\frac{\log p}{p}\ \sim\ \delta\log X\ \xrightarrow[X\to\infty]{}\ \infty .
\]
The factor $\log p$ in the numerator is what fixes the rate at $\log X$: the bare sum $\sum
1/d_{p}$, lower-bounded by $\sum 1/p$, diverges only at the slower rate $\log\log X$ (numerically
$\sum_{p<X,\ \mathrm{adm}}1/d_{p}$ creeps from about $0.24$ at $X=200$ to $0.53$ at $X=5\cdot
10^{4}$). So the two-prime heuristic predicts \emph{infinitely many} two-prime solutions, not
finiteness. It is moreover well calibrated: $\sum_{p<X,\ \mathrm{adm}}(\log p)/d_{p}$ is about
$1.7$ at $X=2000$ and $2.4$ at $X=10^{4}$, consistent with the two known solutions, whose smaller
primes are $61$ and $67$.

\paragraph{A global density heuristic.}
Model $2^{m-1}+3\bmod m$ as roughly equidistributed; then $m$ is a solution with ``probability''
about $1/m$, and $\sum_{m\le X}1/m\sim\log X\to\infty$. A divergent expectation predicts
\emph{infinitely many} solutions. This is precisely the (correct) heuristic for the closely
analogous Fermat-pseudoprime problem $m\mid 2^{m-1}-2$, for which infinitely many solutions are
known (Erd\H{o}s). The inhomogeneous shift ($-3$ in place of $-1$) changes the combinatorics of
how prime factors glue together but not the divergence of the expectation.

Both estimates thus point the same way. The two-prime sum already diverges, and the global
heuristic is dominated by solutions with \emph{many} prime factors, whose expected distribution
(Section~\ref{d1:sec:omega}) also diverges; together they favour an infinite $S_{-3}$ whose elements
beyond $m_{2}$ are many-prime integers far beyond any search. On this view the two known terms
being two-prime reflects the limits of factorisation, not the true distribution of solutions.

A proof of \emph{finiteness} would need an upper bound, uniform in $p$, on the number of prime
factors of $2^{p-1}+3$ lying in prescribed arithmetic progressions, which is well beyond current
analytic number theory. A proof of \emph{infinitude} would need a Carmichael-style construction
producing solutions to order; the homogeneous case $a=-1$ admits the relevant root-of-unity
structure (and is settled in the opposite direction, $S_{-1}=\varnothing$), but for $a=-3$ the
varying targets $k_{q}$ obstruct the standard construction and none is known. Both directions are
presently out of reach.

\subsubsection{Open questions}

\begin{enumerate}[label=\arabic*.]
\item \emph{Is $S_{-3}$ finite?} This remains open; the heuristics of Section~\ref{d1:sec:finite}
point toward an infinite, though computationally unreachable, solution set.
\item \emph{Can the admissibility sieve be sharpened locally?} No: Proposition~\ref{d1:prop:nosharp}
shows that prime-power moduli yield nothing further, so any improvement must be global.
\item \emph{Is every $m\in S_{-3}$ squarefree?} This is open in general;
Corollary~\ref{d1:cor:sqfree} settles it at every prime below $1093$, and a uniform treatment of
large primes, together with a separate analysis of the Wieferich primes $1093,3511$, would be
needed for the full statement.
\item \emph{The same questions for $a\in\{-5,-7,\dots\}$.} The proof of Theorem~\ref{d1:thm:minus1}
uses critically that $-1$ is a root of unity; for general $a$ the $2$-adic argument fails and the
qualitative behaviour of $S_{a}$ is unclear.
\end{enumerate}


\subsubsection{Computational programs}\label{d1:app:programs}
The programs used for the computations reported here---the direct search, the two-prime
factoring search via $N_{p}=2^{p-1}+3$, and the supporting order and discrete-logarithm
routines---are available at
\begin{center}
\url{https://github.com/RobinCodes/Projects/tree/main/Divisibility%20problem}.
\end{center}
The factorisations of Table~\ref{d1:tab:search}, the discrete-logarithm verification underlying
Corollary~\ref{d1:cor:sqfree}, and the sieve data of Table~\ref{d1:tab:sieve} were carried out with
\texttt{sympy}.

\subsubsection{Note on the use of artificial intelligence}\label{d1:app:ai}
This work was prepared with substantial assistance from an artificial-intelligence system, and
most of the results presented here---the admissibility sieve, the interlocking-congruence
description, the two-prime reduction, the squarefreeness analysis, and the computations of
Section~\ref{d1:sec:comp}---were created with such assistance. We do not specify the particular
system used; readers seeking further detail about the methods are invited to contact the authors:
Robin Rovenszky (\texttt{robincodes} on Discord) and notxxdog (\texttt{.notxxdog} on Discord).

No result here is accepted on the basis of machine assertion. Every definition, lemma, and proof
is stated so as to be checkable by hand, and every computational claim is reproducible by
re-running the programs of Appendix~\ref{d1:app:programs}.

\begin{thebibliography}{9}
\bibitem{d1:oeis} OEIS Foundation Inc.\ (2024), Entry A245728 in \emph{The On-Line Encyclopedia of
Integer Sequences}, \url{https://oeis.org/A245728}.
\end{thebibliography}


\subsection{Paper 2, Master paper}

\begin{abstract}
For an integer $a$ let $\Sa=\{\,n\ge 1:\ (n+1)\dvd 2^{n}-a\,\}$; equivalently, with $m=n+1$, the set of $m\ge2$ with $m\dvd 2^{m-1}-a$. We classify these sets according to the shift $a$. Three regimes occur. For $a=0$ and for every $a=2^{j}$ with $j\ge0$ the set is infinite, and we exhibit an explicit one--parameter family of solutions; the mechanism is that the discrete--logarithm target $\log_2 a$ is the \emph{same} at every prime, which makes a Carmichael--type construction available. For $a=-1$ the set is empty, by a two--adic argument that uses only that $-1$ is a root of unity of bounded order; we show that no other $a$ admits an emptiness proof of this kind, and that $a=-1$ is in fact the only even--free value that is empty. For all remaining $a$ the basic question---finiteness, infinitude, and in some cases mere non--emptiness---is open, and we argue that every such case is exactly as intractable as the shift $a=-3$. We develop the conditional structure theory uniformly in $a$: a local admissibility sieve on the prime divisors of a solution, an interlocking--congruence description of squarefree solutions, the resulting reduction of the two--prime problem to the factorisation of a single integer $2^{p-1}-a$, and a criterion for repeated prime factors. The last shows in particular that the squarefreeness statement known for $a=-3$ is special to that shift and fails for, e.g., $a=9$, where $1715=5\cdot7^{3}$ is a solution. We finish with the heuristics, the computational evidence for $-20\le a\le 20$, and the open problems. We also record solutions for two open shifts---$S_{-7}$ contains $449\cdot14779$ and $7459\cdot23459$, and $S_5$ contains $701\cdot34851539\approx2.4\times10^{10}$---which leaves $a=-15$ as the only shift in $[-20,20]$ whose non--emptiness is undecided, the natural successor to the $a=-3$ question.
\end{abstract}


\subsubsection{Introduction}

Fix $a\in\Z$ and put
\[
\Sa=\{\,n\ge1:\ (n+1)\dvd 2^{n}-a\,\}.
\]
Setting $m=n+1$ turns membership $n\in\Sa$ into the congruence
\begin{equation}\label{d2:eq:main}
m\dvd 2^{m-1}-a,\qquad m\ge2.
\end{equation}
We work throughout with $m$ and identify $\Sa$ with the set of $m\ge2$ satisfying \eqref{d2:eq:main}. The case $a=1$ is the familiar Fermat condition $2^{m-1}\equiv1$: every odd prime satisfies it, and the composite solutions are the base--$2$ Fermat pseudoprimes. The cases $a=-1$ and $a=-3$ are treated in the companion work \cite{d2:prior}, where $S_{-1}$ is shown to be empty and, for $a=-3$, only two solutions are known, each a product of two primes, the larger exceeding $10^{18}$. The sequence of $n$ with $(n+1)\dvd2^{n}+3$ is \cite[A245728]{d2:oeis}.

The purpose of this paper is to treat \eqref{d2:eq:main} for all $a$ at once and to explain why a handful of shifts are trivial---in one direction or the other---while every other shift inherits the full difficulty of the $a=-3$ problem. The outcome is a clean division of $\Z$ into three regions.

\begin{theorem}[Classification]\label{d2:thm:main}
Let $a\in\Z$.
\begin{enumerate}
\item\emph{(Soluble, infinite.)} If $a=0$ then $\Sa=\{2^{t}:t\ge1\}$. If $a=2^{j}$ for some integer $j\ge0$, then $\Sa$ is infinite; indeed, writing $c=j+1=2^{s}c'$ with $c'$ odd and $e=\ord_{c'}(2)$, every prime $p$ with $p\nmid 2c$ and $p\equiv1\pmod e$ gives a solution $m=cp$.
\item\emph{(Empty.)} $S_{-1}=\varnothing$. Moreover no even $a$ has $\Sa=\varnothing$, and the two--adic argument that proves $S_{-1}=\varnothing$ applies to no shift other than $a=\pm1$.
\item\emph{(Open.)} For every other $a$ the basic questions are open: $\Sa$ is conjecturally infinite but extraordinarily sparse, no infinitude is known, no finiteness is known, and for the shifts $a$ with
 $a-1=\pm2^{k}$ (the family of $a=-3$) even non--emptiness can be open.
\end{enumerate}
\end{theorem}

The proof of (1) and (2) is elementary and occupies Sections \ref{d2:sec:reductions}--\ref{d2:sec:soluble}. Part (3) is not a theorem in the literal sense but a description of the state of knowledge; Sections \ref{d2:sec:sieve}--\ref{d2:sec:heuristics} make it precise by developing, uniformly in $a$, the structure theory that any solution must obey and by pinning down where the obstruction to a complete solution lies.

A single elementary fact organises much of the discussion. For an \emph{odd prime} $p$, Fermat's little theorem reduces \eqref{d2:eq:main} at $m=p$ to $1\equiv a\pmod p$, so
\begin{equation}\label{d2:eq:primes}
p\in\Sa\iff p\dvd a-1\qquad(p\text{ an odd prime}).
\end{equation}
Thus the prime members of $\Sa$ are governed by the single number $a-1$: they are finite in number unless $a=1$, in which case $a-1=0$ and every odd prime qualifies. This is why $a=1$ is infinite for a trivial reason, why every member of $S_{-3}$ is composite (as $a-1=-4$ has no odd prime factor), and why the shifts with $a-1=\pm2^{k}$ are the hardest: they offer no prime solution to start from.

Two features distinguish the soluble shifts of part (1) from everything else. At the level of the construction, a solution of the shape $m=cp$ forces, modulo the varying prime $p$, the relation $a\equiv 2^{c-1}$, so a one--parameter family of this kind can only produce a power of $2$. At the level of the gluing problem of Section \ref{d2:sec:gluing}, a solution $m=p_1\cdots p_\omega$ requires each prime $p_i$ to meet a congruence whose target is $\log_2 a$ taken modulo $\ord_{p_i}(2)$; these targets are uniform in $i$ precisely when $a$ is a power of $2$, and it is this uniformity that lets one assemble infinitely many solutions, just as in the construction of Fermat pseudoprimes and Carmichael numbers \cite{d2:cipolla,d2:agp}. For a shift that is not a power of $2$ the targets genuinely vary from prime to prime, and no construction is known. The empty shift $a=-1$ escapes in the opposite manner: $-1$ is a root of unity of order $2$ modulo every odd prime, and this rigidity produces a global two--adic contradiction. No other shift is a root of unity of bounded order, so the contradiction is unavailable.

\paragraph{Notation.} For an odd prime $p$ with $p\nmid2$ we write $d_p=\ord_p(2)$, the multiplicative order of $2$ in $\F_p^{\times}$. For $a$ with $\gcd(a,p)=1$ and $a\in\langle2\rangle\subseteq\F_p^{\times}$ we write $\log_2a$ for the least non--negative $k$ with $2^{k}\equiv a\pmod p$, and use the same notation modulo $p^{2}$ when needed. We write $v_2(N)$ for the two--adic valuation, $\omega(m)$ for the number of distinct prime factors, and reserve the letter $p$ for primes.

\subsubsection{Elementary reductions}\label{d2:sec:reductions}

The following constraints hold for every shift and are used repeatedly.

\begin{lemma}\label{d2:lem:parity}
Let $m\in\Sa$.
\begin{enumerate}
\item Every odd prime dividing $m$ is coprime to $a$; equivalently $\gcd(a,m)$ is a power of $2$.
\item If $a$ is odd then $m$ is odd.
\item If $a$ is even then $v_2(m)\le v_2(a)$.
\end{enumerate}
\end{lemma}

\begin{proof}
(1) If an odd prime $p$ divides both $a$ and $m$, then $p\dvd 2^{m-1}-a$ and $p\dvd a$ give $p\dvd2^{m-1}$, which is impossible. (2) For $m\ge2$ the number $2^{m-1}$ is even, so $2^{m-1}-a$ is odd when $a$ is odd, forcing $m$ odd. (3) Put $s=v_2(m)$. Since $m\ge 2^{s}$ we have $m-1\ge 2^{s}-1\ge s$, hence $2^{s}\dvd2^{m-1}$, and $2^{s}\dvd m\dvd 2^{m-1}-a$ gives $2^{s}\dvd a$.
\end{proof}

\begin{proposition}[Prime members]\label{d2:prop:primes}
For an odd prime $p$ one has $p\in\Sa$ iff $p\dvd a-1$. For $p=2$ one has $2\in\Sa$ iff $a$ is even. Consequently $\Sa$ contains infinitely many primes if and only if $a=1$; for $a\ne1$ the prime members are exactly the odd prime divisors of $a-1$, together with $2$ when $a$ is even.
\end{proposition}

\begin{proof}
For odd $p$, $2^{p-1}\equiv1\pmod p$, so \eqref{d2:eq:main} at $m=p$ reads $1\equiv a\pmod p$. For $p=2$ the condition is $2\dvd 2-a$. The corollary about infinitude is immediate, since $a-1=0$ only for $a=1$.
\end{proof}

\begin{corollary}\label{d2:cor:nonempty}
If $a$ is even then $2\in\Sa$, and if $a$ is odd with an odd prime factor $p\dvd a-1$ then $p\in\Sa$. Hence $\Sa=\varnothing$ forces $a$ to be odd with $a-1=\pm2^{k}$, i.e.
\[
a\in\{\,2^{k}+1:k\ge1\,\}\cup\{\,1-2^{k}:k\ge1\,\}=\{3,5,9,17,\dots\}\cup\{-1,-3,-7,-15,\dots\}.
\]
\end{corollary}

Corollary \ref{d2:cor:nonempty} isolates the only candidates for emptiness. We shall see that exactly one of them, $a=-1$, is provably empty, that several ($-7,-3,3,5,9,17,\dots$) are non--empty, and that for the remainder the question is open.

\subsubsection{The soluble shifts}\label{d2:sec:soluble}

\subsubsection{The shift $a=0$}

\begin{proposition}\label{d2:prop:zero}
$S_0=\{2^{t}:t\ge1\}$.
\end{proposition}

\begin{proof}
The condition $m\dvd 2^{m-1}$ holds iff $m$ is a power of $2$. If $m=2^{t}$ then $t\le 2^{t}-1$, so $2^{t}\dvd2^{2^{t}-1}$ and $m\in S_0$. Conversely a divisor of a power of $2$ is a power of $2$.
\end{proof}

\subsubsection{Powers of two}

\begin{theorem}\label{d2:thm:powers}
Let $j\ge0$ and $a=2^{j}$. Put $c=j+1$, write $c=2^{s}c'$ with $c'$ odd, and let $e=\ord_{c'}(2)$ (with $e=1$ if $c'=1$). For every prime $p$ with $p\nmid 2c$ and $p\equiv1\pmod e$, the integer $m=cp$ lies in $\Sa$. In particular $\Sa$ is infinite.
\end{theorem}

\begin{proof}
We check $2^{m-1}\equiv a\pmod m$ on each of the coprime factors $p$, $2^{s}$, $c'$ of $m=cp$.

Modulo $p$: since $cp-1\equiv c-1\pmod{p-1}$, Fermat gives $2^{m-1}=2^{cp-1}\equiv2^{c-1}=2^{j}=a$.

Modulo $2^{s}$: both sides vanish. Indeed $m-1=cp-1\ge s$, so $2^{m-1}\equiv0$; and $a=2^{c-1}$ with $c-1\ge 2^{s}-1\ge s$, so $a\equiv0$.

Modulo $c'$: as $2$ is a unit, $2^{m-1}\cdot 2^{-(c-1)}=2^{c(p-1)}$, and $e=\ord_{c'}(2)$ divides $p-1$ by hypothesis, hence divides $c(p-1)$; thus $2^{m-1}\equiv2^{c-1}=a\pmod{c'}$.

By the Chinese remainder theorem $2^{m-1}\equiv a\pmod m$. There are infinitely many primes $p\equiv1\pmod e$ by Dirichlet, all but finitely many coprime to $2c$.
\end{proof}

\begin{example}
The construction specialises as follows: $a=1$ ($c=1$) gives every odd prime; $a=2$ ($c=2$) gives $m=2p$ for all odd $p$; $a=4$ ($c=3$, $e=2$) gives $m=3p$ for all odd $p\ne3$; $a=8$ ($c=4$) gives $m=4p$; $a=16$ ($c=5$, $e=4$) gives $m=5p$ for $p\equiv1\pmod4$. Thus, for example, $6,10,14,\dots\in S_2$ and $15,21,33,\dots\in S_4$.
\end{example}

\begin{remark}\label{d2:rem:onlypowers}
A family of the shape $m=cp$ with $c$ fixed and $p\to\infty$ can only produce a power of $2$: reducing \eqref{d2:eq:main} modulo $p$ gives $2^{c-1}\equiv a\pmod p$ for all $p$ in the family, whence $a=2^{c-1}$. The deeper reason that powers of $2$ are constructible while other shifts are not appears in Section \ref{d2:sec:gluing}.
\end{remark}

\subsubsection{The empty shift $a=-1$}

We prove $S_{-1}=\varnothing$ and isolate the property of $-1$ on which the argument depends.

\begin{theorem}\label{d2:thm:minusone}
$S_{-1}=\varnothing$.
\end{theorem}

\begin{proof}
Suppose $m>1$ satisfies $m\dvd2^{m-1}+1$. Then $2^{m-1}+1$ is odd, so $m$ is odd. Fix a prime $p\dvd m$ and put $d=d_p$. From $2^{m-1}\equiv-1\pmod p$ the order of $2^{m-1}$ in $\F_p^{\times}$ is $2$, that is $d/\gcd(d,m-1)=2$. Hence $g:=\gcd(d,m-1)=d/2$, so $v_2(g)=v_2(d)-1$. On the other hand $g=\gcd(d,m-1)$ has $v_2(g)=\min(v_2(d),v_2(m-1))$, and since this minimum is strictly less than $v_2(d)$ it equals $v_2(m-1)$. Therefore
\[
v_2(d)=v_2(m-1)+1.
\]
As $d\dvd p-1$, we get $v_2(p-1)\ge v_2(m-1)+1$. Writing $\nu=v_2(m-1)+1$, every prime $p\dvd m$ satisfies $p\equiv1\pmod{2^{\nu}}$; a product of such primes is $\equiv1\pmod{2^{\nu}}$, so $m\equiv1\pmod{2^{\nu}}$, i.e.\ $v_2(m-1)\ge\nu=v_2(m-1)+1$, a contradiction.
\end{proof}

\begin{remark}\label{d2:rem:whyminusone}
The argument needs $2^{2(m-1)}\equiv1\pmod p$ for every $p\dvd m$, which holds because $(-1)^{2}=1$ as an integer; the only shifts with $a^{2}=1$ are $a=\pm1$. For $a=1$ the same steps give $d\dvd m-1$ with no contradiction, consistent with $S_1$ being infinite. A mild generalisation---that the argument applies whenever $a\equiv-1\pmod p$ for all $p\dvd m$---constrains only those $m$ built from the finitely many prime divisors of $a+1$, and so yields no global emptiness when $a\ne-1$. Combined with Corollary \ref{d2:cor:nonempty} (no even shift is empty), this is the precise sense in which $a=-1$ is the unique empty shift accessible by elementary means.
\end{remark}

\subsubsection{The local sieve}\label{d2:sec:sieve}

From here on we develop the structure that any solution of \eqref{d2:eq:main} must obey, valid for every $a$. We assume $a$ is odd in this section, so that all solutions are odd by Lemma \ref{d2:lem:parity}; the even case is similar with the obvious modifications at the prime $2$.

\begin{lemma}[Local condition]\label{d2:lem:local}
Let $m\in\Sa$ and let $p$ be an odd prime with $p\dvd m$ and $p\nmid a$. Put $d=d_p$ and $t=m/p^{v_p(m)}$. Then $a\in\langle2\rangle\subseteq\F_p^{\times}$, and with $k=\log_2a$,
\begin{equation}\label{d2:eq:local}
t\equiv k+1\pmod d.
\end{equation}
\end{lemma}

\begin{proof}
Since $d\dvd p-1$ we have $p\equiv1\pmod d$, hence $m\equiv t\pmod d$ and $m-1\equiv t-1\pmod d$. Then $2^{m-1}\equiv2^{t-1}\pmod p$, and as $2^{m-1}\equiv a$ this gives $2^{t-1}\equiv a\pmod p$, so $a\in\langle2\rangle$ and $t-1\equiv k\pmod d$.
\end{proof}

\begin{definition}\label{d2:def:admissible}
A prime $p\nmid 6a$ is \emph{$a$--admissible} if $a\in\langle2\rangle\pmod p$ and, with $d=d_p$ and $k=\log_2a$,
\begin{itemize}
\item[(A)] $d$ even $\Rightarrow$ $k$ even;
\item[(B)] $3\dvd d$ $\Rightarrow$ $k\not\equiv2\pmod3$.
\end{itemize}
\end{definition}

\begin{theorem}\label{d2:thm:admissible}
Let $a$ be odd. Every prime $p\nmid3a$ dividing a solution $m\in\Sa$ is $a$--admissible. Moreover, if $a\not\equiv1\pmod3$ then $3\nmid m$, and in that case condition \textnormal{(B)} is active.
\end{theorem}

\begin{proof}
The cofactor $t=m/p^{v_p(m)}$ is odd, since $m$ is. Reducing \eqref{d2:eq:local} modulo $2$: if $d$ is even then $t\equiv k+1\pmod2$, and $t$ odd forces $k$ even. For (B), first note that $m$ odd gives $2^{m-1}\equiv1\pmod3$, so if $3\dvd m$ then $a\equiv1\pmod3$; contrapositively, $a\not\equiv1\pmod3$ implies $3\nmid m$, hence $3\nmid t$. Reducing \eqref{d2:eq:local} modulo $3$ when $3\dvd d$ gives $t\equiv k+1\pmod3$ with $t\not\equiv0$, so $k+1\not\equiv0$, i.e.\ $k\not\equiv2\pmod3$.
\end{proof}

For $a=-3$ one has $-3\equiv0\not\equiv1\pmod3$, so both conditions are active; Definition \ref{d2:def:admissible} then reduces to the sieve of \cite{d2:prior}. For shifts with $a\equiv1\pmod3$ condition (B) is vacuous and $3$ may divide $m$.

\begin{remark}[No prime--power sharpening]\label{d2:rem:nosharpen}
Conditions (A) and (B) cannot be strengthened by replacing the moduli $2,3$ with higher prime powers. The constraints available on the cofactor $t$ in isolation are precisely $t$ odd and (when $a\not\equiv1\pmod3$) $3\nmid t$. If $2^{e}\dvd d$ then \eqref{d2:eq:local} forces $t\equiv k+1\pmod{2^{e}}$; among these residues the requirement $t$ odd excludes a value of $k$ only through its parity, which is (A). Likewise $3^{f}\dvd d$ forces $t\equiv k+1\pmod{3^{f}}$, and $3\nmid t$ excludes a value of $k$ only through $k+1\bmod3$, which is (B). Any genuine sharpening must therefore come from the interaction of distinct prime factors, treated next.
\end{remark}

The sieve removes a positive proportion of primes as possible factors. Under an Artin--type heuristic \cite{d2:hooley} the admissible primes have a positive logarithmic density; empirically this density is near $0.15$. The smallest admissible primes, however, depend erratically on $a$, and this is what governs how small the smallest solution can be (Section \ref{d2:sec:heuristics}).

\subsubsection{Gluing: interlocking congruences}\label{d2:sec:gluing}

The local conditions of Section \ref{d2:sec:sieve} combine through the Chinese remainder theorem into a description of all solutions of a given shape.

\begin{proposition}[Squarefree solutions]\label{d2:prop:squarefree}
Let $m=p_1\cdots p_\omega$ be squarefree with distinct odd primes $p_i\nmid a$. Then $m\in\Sa$ if and only if, for every $i$,
\[
a\in\langle2\rangle\pmod{p_i}\qquad\text{and}\qquad \prod_{j\ne i}p_j\equiv k_i+1\pmod{d_i},
\]
where $d_i=d_{p_i}$ and $k_i=\log_2a\bmod p_i$. In particular each $p_i$ is admissible.
\end{proposition}

\begin{proof}
By the Chinese remainder theorem $m\dvd2^{m-1}-a$ iff $p_i\dvd2^{m-1}-a$ for all $i$, i.e.\ $2^{m-1}\equiv a\pmod{p_i}$, which requires $a\in\langle2\rangle$ and $m-1\equiv k_i\pmod{d_i}$. Since $p_i\equiv1\pmod{d_i}$ we have $m\equiv\prod_{j\ne i}p_j\pmod{d_i}$, giving the stated form.
\end{proof}

Thus a squarefree solution is a set of admissible primes whose product lands simultaneously in the prescribed class $k_i+1$ modulo each $d_i$. The targets $k_i=\log_2a\bmod d_i$ are the crux: they are uniform across $i$ exactly when $a$ is a power of $2$ (then $k_i\equiv j$), and it is this uniformity that underlies the construction of Theorem \ref{d2:thm:powers} and, more generally, the classical constructions of pseudoprimes and Carmichael numbers \cite{d2:cipolla,d2:agp}. For a shift that is not a power of $2$ the targets vary with $i$, and no method is known to assemble infinitely many solutions.

Specialising to two primes recovers, and slightly clarifies, the main computational tool.

\begin{corollary}[Two primes]\label{d2:cor:twoprime}
Let $p<q$ be odd primes with $\gcd(pq,a)=1$. Then $pq\in\Sa$ iff
\[
q\dvd 2^{p-1}-a\qquad\text{and}\qquad p\dvd 2^{q-1}-a.
\]
Fixing $p$ and writing $N_p=2^{p-1}-a$, $d_p=\ord_p2$, $k_p=\log_2a\bmod p$ (assuming $a\in\langle2\rangle\bmod p$), this is equivalent to
\[
q\dvd N_p\qquad\text{and}\qquad q\equiv k_p+1\pmod{d_p}.
\]
\end{corollary}

\begin{proof}
Both primes must divide $2^{pq-1}-a$. Since $pq-1\equiv p-1\pmod{q-1}$ we have $2^{pq-1}\equiv2^{p-1}\pmod q$, so the condition at $q$ is $2^{p-1}\equiv a\pmod q$, i.e.\ $q\dvd N_p$. Symmetrically $2^{pq-1}\equiv2^{q-1}\pmod p$, so the condition at $p$ is $2^{q-1}\equiv a\pmod p$, i.e.\ $q-1\equiv k_p\pmod{d_p}$.
\end{proof}

\begin{remark}\label{d2:rem:onesided}
The condition at the larger prime is the divisibility $q\dvd N_p$, while the condition at the smaller prime is a single residue test modulo $d_p<p$. Hence, for a fixed smaller prime $p$, the two--prime solutions are obtained by factoring the one integer $N_p=2^{p-1}-a$ and testing which of its prime factors $q>p$ satisfy $q\equiv k_p+1\pmod{d_p}$. This is the basis of the search in Section \ref{d2:sec:computation}.
\end{remark}

The same reduction works for any number of primes.

\begin{corollary}[Recursive criterion]\label{d2:cor:recursive}
Let $p_1<\dots<p_\omega$ be distinct admissible primes, set $P=p_1\cdots p_{\omega-1}=m/p_\omega$ and $N_P=2^{P-1}-a$. Then $m=Pp_\omega\in\Sa$ iff
\[
p_\omega\dvd N_P\qquad\text{and}\qquad m\equiv k_i+1\pmod{d_i}\quad(1\le i\le\omega-1).
\]
The congruence at $p_\omega$ is automatic once $p_\omega\dvd N_P$.
\end{corollary}

\begin{proof}
The congruence at $p_\omega$ reads $P\equiv k_\omega+1\pmod{d_\omega}$, i.e.\ $2^{P-1}\equiv a\pmod{p_\omega}$, i.e.\ $p_\omega\dvd N_P$; this holds for every prime factor of $N_P$. The remaining $\omega-1$ congruences are those of Proposition \ref{d2:prop:squarefree}.
\end{proof}

To search for $\omega$--prime solutions extending a fixed tuple $(p_1,\dots,p_{\omega-1})$ one factors the single integer $N_P=2^{P-1}-a$ and residue--tests its prime factors. Two points limit the method. First, solutions do not nest: the smaller primes of a solution need not themselves form a solution, so one cannot grow long solutions from short ones. Second, $N_P$ has about $0.301\,P$ decimal digits, and $P$ is a product of admissible primes, hence grows super--exponentially in $\omega$. Even the smallest admissible primes make $N_P$ unfactorable beyond $\omega=3$. This computational wall is independent of $a$ and is the reason that only two--prime solutions are ever found by search.

\subsubsection{Repeated prime factors}\label{d2:sec:repeated}

Whether a solution can fail to be squarefree is, by contrast with the previous sections, a property that depends sharply on $a$.

Recall that $p$ is a \emph{Wieferich prime} (base $2$) if $2^{p-1}\equiv1\pmod{p^{2}}$; the only such primes below $4\times10^{3}$ are $1093$ and $3511$. For non--Wieferich $p$ one has $\ord_{p^{2}}(2)=p\,d_p$.

\begin{lemma}\label{d2:lem:wieferich}
Let $m\in\Sa$ and let $p$ be a non--Wieferich prime with $p^{2}\dvd m$. Then $a\in\langle2\rangle\subseteq(\Z/p^{2})^{\times}$, and writing $\widetilde k=\log_2a\bmod p^{2}$,
\[
\widetilde k\equiv-1\pmod p.
\]
\end{lemma}

\begin{proof}
From $p^{2}\dvd2^{m-1}-a$ we get $2^{m-1}\equiv a\pmod{p^{2}}$, so $a\in\langle2\rangle$ modulo $p^{2}$ and $m-1\equiv\widetilde k\pmod{D}$ with $D=\ord_{p^{2}}(2)=p\,d_p$. As $p\dvd D$, reduce modulo $p$: $m-1\equiv\widetilde k\pmod p$. But $p^{2}\dvd m$ gives $m\equiv0\pmod p$, so $m-1\equiv-1$, whence $\widetilde k\equiv-1\pmod p$.
\end{proof}

For $a=-3$ a direct computation of $\widetilde k$ for each admissible $p<1093$ shows $\widetilde k\not\equiv-1\pmod p$ in every case; with Lemma \ref{d2:lem:wieferich} this shows that, for $a=-3$, every solution is squarefree at all primes below $1093$. The next remark shows that this conclusion is special to the shift $-3$.

\begin{remark}[Failure of squarefreeness]\label{d2:rem:sqfree-fail}
The congruence $\widetilde k\equiv-1\pmod p$ does occur for other shifts, and then $p^{2}\dvd m$ is permitted. Examples, verified directly:
\[
1715=5\cdot7^{3}\in S_{9},\qquad 49=7^{2}\in S_{15},\qquad 9=3^{2},\ 27=3^{3}\in S_{-5}.
\]
For $a=9$ the primes $p=5$ and $p=7$ both satisfy $\widetilde k\equiv-1$; for $a=15$ the prime $7$ does; for $a=-5$ the prime $3$ does. Thus the squarefreeness of solutions at small primes is not a general phenomenon but a coincidence of the discrete logarithms attached to a particular $a$.
\end{remark}

\subsubsection{Heuristics and the gradient of difficulty}\label{d2:sec:heuristics}

For shifts outside the soluble and empty regimes the only guide is heuristic, and two heuristics are worth recording.

\paragraph{Global density.} Modelling $2^{m-1}-a\bmod m$ as uniform makes $m$ a solution with ``probability'' about $1/m$, and $\sum_{m\le X}1/m\sim\log X\to\infty$. The expectation diverges, predicting infinitely many solutions. This is the same divergent count that governs the Fermat pseudoprimes $a=1$, which are provably infinite \cite{d2:cipolla,d2:pomerance}. The shift $-a$ changes how prime factors glue but not the divergence.

\paragraph{Count by number of prime factors.} Restricting to squarefree $m$ with exactly $\omega$ admissible prime factors and modelling the interlocking congruences of Proposition \ref{d2:prop:squarefree} as independent gives, for the number of $\omega$--prime solutions up to $X$,
\[
E_\omega(X)\approx\frac1{\omega!}\Bigl(\sum_{\substack{p\le X\\ \text{admissible}}}\tfrac1p\Bigr)^{\!\omega}\approx\frac{(\delta\log\log X)^{\omega}}{\omega!},
\]
with $\delta$ the logarithmic density of admissible primes. Summing over $\omega\ge2$ gives $(\log X)^{\delta}-1-\delta\log\log X\to\infty$, again predicting infinitely many but only as a minuscule power of $\log X$. At every reachable height $\delta\log\log X<1$, so $E_\omega$ decreases with $\omega$ and two--prime solutions dominate; this is why the only solutions ever exhibited have two prime factors.

These heuristics must be read with care, because they are overridden by algebra in exactly the soluble and empty cases. The global count would predict $S_{-1}$ infinite, yet $S_{-1}=\varnothing$; and it gives no hint that powers of $2$ are infinite for a structural reason. Where no algebraic structure is present---every shift that is neither a power of $2$ nor $-1$---the heuristics are all we have, and they favour an infinite but unreachable solution set.

\paragraph{A gradient, not a dichotomy, among the hard shifts.} The hard shifts are uniform in their \emph{logical} difficulty: none is closer to being settled than $a=-3$. They differ only in how accessible a single solution is, and this accessibility is controlled by the size of the smallest admissible prime. By Corollary \ref{d2:cor:nonempty} they split into two kinds.

If $a$ is even, or $a-1$ has an odd prime factor, then a prime or the prime $2$ is already a solution; here non--emptiness is free and only infinitude is open.

If $a$ is odd with $a-1=\pm2^{k}$, then there is no prime solution, and the situation mirrors $a=-3$. Among these shifts the smallest admissible primes vary erratically (Table \ref{d2:tab:smallest}), and with them the size of the smallest solution: $a=9$ has admissible primes $5,7,11$ and the small solution $35=5\cdot7$, whereas $a=-3$ has admissible primes $13,19,61,67$ and its smallest solution exceeds $10^{10}$. Solutions are known for two further shifts of this family (Section \ref{d2:sec:computation}): $S_{-7}$ contains $449\cdot14779=6635771$ and $7459\cdot23459=174980681$, and $S_5$ contains $701\cdot34851539=24430928839\approx2.4\times10^{10}$. This leaves $a=-15$ as the only shift in $[-20,20]$ with no known solution, and whether $S_{-15}$ is empty is open---the direct continuation of the $a=-3$ question.

\begin{table}[H]
\centering
\begin{tabular}{rl l}
\toprule
$a$ & smallest admissible primes & smallest known solution\\
\midrule
$9$ & $5,7,11,23,\dots$ & $35=5\cdot7$\\
$3$ & $11,13,23,\dots$ & $10669=47\cdot227$\\
$17$ & $19,47,53,\dots$ & $150281=67\cdot2243$\\
$-3$ & $13,19,61,67,\dots$ & $61\cdot228806497\approx1.4\times10^{10}$\\
$5$ & $11,19,29,\dots$ & $701\cdot34851539\approx2.4\times10^{10}$\\
$-7$ & $11,23,29,\dots$ & $449\cdot14779\approx6.6\times10^{6}$\\
$-15$ & $23,31,47,\dots$ & none known\\
\bottomrule
\end{tabular}
\caption{Shifts with $a-1=\pm2^{k}$, $a$ odd. The smallest admissible prime, an arithmetic accident, governs how small a solution can be; logical difficulty is the same throughout.}
\label{d2:tab:smallest}
\end{table}

\subsubsection{Computational evidence}\label{d2:sec:computation}

\paragraph{Method.} For a single pass we compute $r_m=2^{m-1}\bmod m$ for each $m$ up to a bound; then $m\in\Sa$ for $|a|<m$ precisely when $a\equiv r_m\pmod m$, i.e.\ $a=r_m$ or $a=r_m-m$. This finds every solution of every shift $a\in[-20,20]$ in one sweep (small $m$ with $m\le|a|$ are checked separately). At $m\le2\times10^{6}$ the sweep is a quick \texttt{sympy} computation; for the frontier shifts $a\in\{5,-7,-15\}$ it was pushed to $m\le10^{10}$, and beyond $2.4\times10^{10}$ for $a=5$, with a dedicated multi--threaded implementation using Montgomery modular exponentiation, wheel sieving and small--prime filtering. For the two--prime search we use Remark \ref{d2:rem:onesided}: for each admissible smaller prime $p$ we factor $N_p=2^{p-1}-a$ and test its prime factors $q>p$ against $q\equiv k_p+1\pmod{d_p}$. The factorisations were carried out with trial division, Pollard $\rho$, and ECM in \texttt{sympy}.

\paragraph{Findings.} The prime members observed for every $a\in[-20,20]$ agree exactly with Proposition \ref{d2:prop:primes}. The two--prime search recovers both known solutions of $S_{-3}$, namely $61\cdot228806497$ and $67\cdot44210291368986343$, and also produces $47\cdot227\in S_3$, $67\cdot2243$ and $19\cdot262127\in S_{17}$, and small solutions of $S_9$ and $S_{-5}$. It returns nothing for $a\in\{5,-7,-15\}$ with smaller prime up to $110$; the extended sweep, however, settles two of these three frontier shifts:
\[
S_{-7}\ni 449\cdot14779=6635771\ \text{and}\ 7459\cdot23459=174980681,
\qquad
S_5\ni 701\cdot34851539=24430928839.
\]
Both $S_{-7}$ solutions lie below $10^{10}$ and are the only members there; $S_5$ has no member below $10^{10}$, its smallest known one sitting at $\approx2.44\times10^{10}$. All three are products of two primes whose smaller prime ($449$, $7459$, $701$) exceeds the $110$ cutoff of the two--prime search---which is why only the direct sweep finds them---and each satisfies the criterion of Corollary \ref{d2:cor:twoprime}. For $a=-15$ both searches return nothing. The complete classification for $-20\le a\le20$ is recorded in Table \ref{d2:tab:full}.

\begin{remark}[$a=-7$ mirrors $a=-3$]\label{d2:rem:minus7}
The two known members of $S_{-7}$ reproduce, in miniature, the profile of $S_{-3}$: the search returns exactly two solutions, each a product of two primes, and then nothing below the bound. They are smaller by several orders of magnitude---$449\cdot14779=6635771\approx6.6\times10^{6}$ and $7459\cdot23459=174980681\approx1.7\times10^{8}$, against $61\cdot228806497\approx1.4\times10^{10}$ and $67\cdot44210291368986343\approx3\times10^{18}$ for $S_{-3}$. Both shifts lie in the family $a-1=-2^{k}$ ($-3=1-2^{2}$, $-7=1-2^{3}$) and offer no prime solution, so each is a genuine instance of the hard regime; the order--of--magnitude gap between them is precisely the arithmetic accident of Section \ref{d2:sec:heuristics}, not a difference in logical difficulty. The ``two semiprimes, then silence'' shape is exactly what leaves finiteness versus infinitude inaccessible to computation.
\end{remark}

\subsubsection{Summary: reductions and regimes}\label{d2:sec:summary}

By a \emph{reduction} we mean a result that decides, for at least one shift $a$, either its regime or its (non)emptiness. The proofs of this paper rest on exactly five such reductions; the lemmas of Sections \ref{d2:sec:sieve}--\ref{d2:sec:repeated} constrain the shape of solutions but decide no shift's fate on their own.

\begin{itemize}
\item[\textbf{R1}] (Proposition \ref{d2:prop:primes}, prime members.) For an odd prime $p$, $p\in\Sa$ iff $p\dvd a-1$. \emph{Decides:} $a=1$ is infinite, since $a-1=0$ puts every odd prime in $S_1$.
\item[\textbf{R2}] (Corollary \ref{d2:cor:nonempty}, non--emptiness.) If $a$ is even, or $a$ is odd with an odd prime factor of $a-1$, then $\Sa\ne\varnothing$. \emph{Decides:} non--emptiness of every such $a$. In the range $|a|\le20$ this leaves undecided only the eight shifts with $a-1=\pm2^{k}$, namely $-15,-7,-3,-1,3,5,9,17$.
\item[\textbf{R3}] (Proposition \ref{d2:prop:zero}.) $S_0=\{2^{t}:t\ge1\}$. \emph{Decides:} $a=0$ completely (infinite).
\item[\textbf{R4}] (Theorem \ref{d2:thm:powers}, powers of two.) For every $j\ge0$, $S_{2^{j}}$ is infinite via $m=(j+1)p$. \emph{Decides:} $a\in\{1,2,4,8,16,\dots\}$ infinite.
\item[\textbf{R5}] (Theorem \ref{d2:thm:minusone}.) $S_{-1}=\varnothing$. \emph{Decides:} $a=-1$ empty.
\end{itemize}

R1 and R4 overlap at $a=1$, and R3, R4 between them settle the entire soluble regime $\{0\}\cup\{2^{j}\}$; R5 settles the empty regime; R2 settles non--emptiness for all but the eight $a-1=\pm2^{k}$ shifts. For six of those eight ($-3,3,9,17,5,-7$) non--emptiness is established not by a reduction but by an explicit solution found through the search of Corollary \ref{d2:cor:twoprime} and its extended sweep; the smallest solutions for $-3$ and $5$ exceed $10^{10}$ ($\approx1.4\times10^{10}$ and $\approx2.4\times10^{10}$), while $-7$, though it too offers no prime solution, has its smallest solution already at $\approx6.6\times10^{6}$. This leaves a single shift in range, $a=-15$, with non--emptiness undecided: no reduction applies and no solution is known. Were one found, every $S_a$ with $|a|\le20$ would be non--empty, and the open question for the whole range would shift from existence to distribution---where the heuristics of Section \ref{d2:sec:heuristics} predict each open $S_a$ is infinite. Everything in Table \ref{d2:tab:full} is an instance of this analysis.

\begin{table}[H]
\centering
\footnotesize
\renewcommand{\arraystretch}{1.0}
\begin{tabular}{r l l r r l}
\toprule
$a$ & regime & non--empty? & smallest known $m$ & $\#\{m\le2\!\times\!10^{6}\}$ & odd primes $\mid a-1$\\
\midrule
$-20$ & open & yes & 2 & 10 & $3,\,7$ \\
$-19$ & open & yes & 5 & 3 & $5$ \\
$-18$ & open & yes & 2 & 6 & $19$ \\
$-17$ & open & yes & 3 & 6 & $3$ \\
$-16$ & open & yes & 2 & 15 & $17$ \\
$-15$ & open & \textbf{undecided} & none known & 0 & --- \\
$-14$ & open & yes & 2 & 10 & $3,\,5$ \\
$-13$ & open & yes & 7 & 3 & $7$ \\
$-12$ & open & yes & 2 & 13 & $13$ \\
$-11$ & open & yes & 3 & 5 & $3$ \\
$-10$ & open & yes & 2 & 6 & $11$ \\
$-9$ & open & yes & 5 & 5 & $5$ \\
$-8$ & open & yes & 2 & 29 & $3$ \\
$-7$ & open & yes & $6635771$ & 0 & --- \\
$-6$ & open & yes & 2 & 5 & $7$ \\
$-5$ & open & yes & 3 & 6 & $3$ \\
$-4$ & open & yes & 2 & 30 & $5$ \\
$-3$ & open & yes (huge) & $\approx1.4\times10^{10}$ & 0 & --- \\
$-2$ & open & yes & 2 & 4 & $3$ \\
$-1$ & empty & no & --- & 0 & --- \\
$0$ & soluble & yes ($\infty$) & 2 & $2^{t}$ only & --- \\
$1$ & soluble & yes ($\infty$) & 3 & $149286$ & all odd primes \\
$2$ & soluble & yes ($\infty$) & 2 & $78967$ & --- \\
$3$ & open & yes & 10669 & 1 & --- \\
$4$ & soluble & yes ($\infty$) & 2 & $54616$ & $3$ \\
$5$ & open & yes (huge) & $\approx2.4\times10^{10}$ & 0 & --- \\
$6$ & open & yes & 2 & 5 & $5$ \\
$7$ & open & yes & 3 & 1 & $3$ \\
$8$ & soluble & yes ($\infty$) & 2 & $42152$ & $7$ \\
$9$ & open & yes & 35 & 14 & --- \\
$10$ & open & yes & 2 & 3 & $3$ \\
$11$ & open & yes & 5 & 1 & $5$ \\
$12$ & open & yes & 2 & 6 & $11$ \\
$13$ & open & yes & 3 & 5 & $3$ \\
$14$ & open & yes & 2 & 8 & $13$ \\
$15$ & open & yes & 7 & 4 & $7$ \\
$16$ & soluble & yes ($\infty$) & 2 & $17365$ & $3,\,5$ \\
$17$ & open & yes & 150281 & 1 & --- \\
$18$ & open & yes & 2 & 7 & $17$ \\
$19$ & open & yes & 3 & 3 & $3$ \\
$20$ & open & yes & 2 & 8 & $19$ \\
\bottomrule
\end{tabular}
\caption{Every integer shift $-20\le a\le20$. ``Regime'' is one of \emph{soluble} (proved infinite, R3--R4), \emph{empty} (R5), or \emph{open}. The ``non--empty?'' column is decided by R2 for every shift with an odd prime factor of $a-1$ (last column non--empty) and for every even shift (then $2\in\Sa$, smallest known $m=2$); the six open shifts $-3,3,9,17,5,-7$ are non--empty by an exhibited solution; and exactly one in range, $a=-15$, has non--emptiness undecided. Counts are exact for $m\le2\times10^{6}$; the soluble counts grow with the bound, while open shifts show only a handful, the signature of a sparse set fed by a positive--density sieve. The shifts $a=-3$ and $a=-7$ have no solution below the bound but two known above it; $a=5$ has its single known solution near $2.4\times10^{10}$.}
\label{d2:tab:full}
\end{table}

\medskip

Table \ref{d2:tab:full} records what is known about each set; Table \ref{d2:tab:verified} records how hard each was looked at. By a search bound we mean the largest power of ten below which the set has been certified complete, and by the method that produced it we mean the reduction or search responsible for the smallest known solution. The single--pass sweep was carried to $2\times10^{6}$ for every shift, so all $m\le10^{6}$ are certified throughout; it was extended to $10^{10}$ for the frontier shifts $a\in\{5,-7,-15\}$---and just past the smallest solution, at $\approx2.4\times10^{10}$, for $a=5$---and to $10^{15}$ for $a=-3$, in agreement with \cite{d2:oeis}. The two--prime search of Corollary \ref{d2:cor:twoprime} was run independently for every admissible smaller prime $p\le110$; because its reach is set by the \emph{smaller} prime and not by the magnitude of $m$, it certifies two--prime solutions of any size---it recovers the $\approx3\times10^{18}$ member of $S_{-3}$---but it does not see $S_5$ or $S_{-7}$, whose smaller primes ($701$, and $449,7459$) exceed the cutoff and which the extended sweep alone resolves. For the soluble shifts ($a=0$ and $a=2^{j}$) and the empty shift ($a=-1$) the answer holds for all $n$ by R3--R5, so no finite bound is relevant; the sweep was nonetheless run over them to $2\times10^{6}$ as a consistency check.

\begin{table}[H]
\centering
\footnotesize
\renewcommand{\arraystretch}{1.0}
\begin{tabular}{r l c l}
\toprule
$a$ & verified up to & two--prime reach & key result\\
\midrule
$-20$ & $10^{6}$ & $p\le110$ & R2 \\
$-19$ & $10^{6}$ & $p\le110$ & R2 \\
$-18$ & $10^{6}$ & $p\le110$ & R2 \\
$-17$ & $10^{6}$ & $p\le110$ & R2 \\
$-16$ & $10^{6}$ & $p\le110$ & R2 \\
$-15$ & $10^{10}$ & $p\le110$ & open \\
$-14$ & $10^{6}$ & $p\le110$ & R2 \\
$-13$ & $10^{6}$ & $p\le110$ & R2 \\
$-12$ & $10^{6}$ & $p\le110$ & R2 \\
$-11$ & $10^{6}$ & $p\le110$ & R2 \\
$-10$ & $10^{6}$ & $p\le110$ & R2 \\
$-9$ & $10^{6}$ & $p\le110$ & R2 \\
$-8$ & $10^{6}$ & $p\le110$ & R2 \\
$-7$ & $10^{10}$ & $p\le110$ & sweep \\
$-6$ & $10^{6}$ & $p\le110$ & R2 \\
$-5$ & $10^{6}$ & $p\le110$ & R2 \\
$-4$ & $10^{6}$ & $p\le110$ & R2 \\
$-3$ & $10^{15}$ & $p\le110$ & two--prime \\
$-2$ & $10^{6}$ & $p\le110$ & R2 \\
$-1$ & all $n$ & --- & R5 (empty) \\
$0$ & all $n$ & --- & R3 \\
$1$ & all $n$ & --- & R1, R4 \\
$2$ & all $n$ & --- & R4 \\
$3$ & $10^{6}$ & $p\le110$ & two--prime \\
$4$ & all $n$ & --- & R4 \\
$5$ & $10^{10}$ & $p\le110$ & sweep \\
$6$ & $10^{6}$ & $p\le110$ & R2 \\
$7$ & $10^{6}$ & $p\le110$ & R2 \\
$8$ & all $n$ & --- & R4 \\
$9$ & $10^{6}$ & $p\le110$ & two--prime \\
$10$ & $10^{6}$ & $p\le110$ & R2 \\
$11$ & $10^{6}$ & $p\le110$ & R2 \\
$12$ & $10^{6}$ & $p\le110$ & R2 \\
$13$ & $10^{6}$ & $p\le110$ & R2 \\
$14$ & $10^{6}$ & $p\le110$ & R2 \\
$15$ & $10^{6}$ & $p\le110$ & R2 \\
$16$ & all $n$ & --- & R4 \\
$17$ & $10^{6}$ & $p\le110$ & two--prime \\
$18$ & $10^{6}$ & $p\le110$ & R2 \\
$19$ & $10^{6}$ & $p\le110$ & R2 \\
$20$ & $10^{6}$ & $p\le110$ & R2 \\
\bottomrule
\end{tabular}
\caption{How far each $S_a$, $-20\le a\le20$, has been searched and by what means. \emph{Verified up to}: the largest power of ten $10^{k}$ such that every $m\le10^{k}$ has been certified by the single--pass sweep, so that the only solutions below this bound are those in Table \ref{d2:tab:full}; ``all $n$'' marks the shifts settled for every $n$ by R3--R5 (the $2\times10^{6}$ sweep over them is only a check). \emph{Two--prime reach}: the largest smaller prime $p$ tried in the factoring search of Corollary \ref{d2:cor:twoprime} and Remark \ref{d2:rem:onesided}; this locates $pq\in\Sa$ with $p\le110$ irrespective of the size of $q$. \emph{Key result}: the reduction or search that settles the shift---\textbf{R2} a prime (or $2$) member, \textbf{R3}/\textbf{R4} the soluble construction, \textbf{R5} the emptiness proof, \emph{two--prime} a solution from the factoring search ($p\le110$, smaller primes here $47,67,5,61,67$ for $3,17,9,-3$), \emph{sweep} a solution found only by the extended single--pass search because its smaller prime exceeds $110$ ($701$ for $5$; $449,7459$ for $-7$), and \emph{open} non--emptiness still undecided ($-15$).}
\label{d2:tab:verified}
\end{table}

\subsubsection{Open problems}\label{d2:sec:open}

\begin{enumerate}
\item \emph{Is any open $\Sa$ finite?} This is open for every shift outside $\{0,-1\}\cup\{2^{j}\}$. A proof of finiteness would require an upper bound, uniform in $p$, on the number of prime factors of $2^{p-1}-a$ lying in prescribed residue classes, which is beyond present analytic number theory.
\item \emph{Is any open $\Sa$ infinite?} Equally open. The constructions that settle the powers of $2$ (and the classical pseudoprime constructions) rely on the discrete--logarithm targets being uniform across primes; for a shift that is not a power of $2$ these targets vary, and no construction is known.
\item \emph{The last frontier shift.} Is $S_{-15}$ empty? With the solutions exhibited here for $S_{-7}$ and $S_5$ (Section \ref{d2:sec:computation}), $a=-15$ is the unique shift in $[-20,20]$ whose non--emptiness is undecided, and the immediate successor of the $a=-3$ question. A single exhibited solution would settle it; a proof of emptiness would require a mechanism beyond the root--of--unity argument, which is unavailable here. Should such a solution be found, every $S_a$ with $|a|\le20$ would be non--empty---closing the existence question across the whole range and redirecting the problem toward distribution, where the heuristics of Section \ref{d2:sec:heuristics} predict each open $S_a$ is infinite. Proving infinitude for even a single hard shift would then be the natural next goal.
\item \emph{Squarefreeness.} By Lemma \ref{d2:lem:wieferich} a repeated prime factor $p$ forces $\widetilde k\equiv-1\pmod p$ at non--Wieferich $p$; this fails for $a=-3$ at all small primes but holds for $a=9,15,-5$ at small primes. For which $a$ are solutions squarefree at small primes, and is there a clean description in terms of $a$ of the set of $p$ with $\widetilde k\equiv-1$?
\item \emph{General even shifts.} The even shifts admit both even and odd solutions and the prime $2$ is always present; a systematic treatment of the even regime, parallel to Sections \ref{d2:sec:sieve}--\ref{d2:sec:gluing}, remains to be written out.
\end{enumerate}


\subsubsection{Computational programs}\label{d2:app:programs}
The programs used for the computations of Section~\ref{d2:sec:computation}---the single--pass residue sweep, the dedicated multi--threaded extended sweep (Montgomery modular exponentiation, wheel sieving, and small--prime filtering), and the two--prime factoring search---are available at
\begin{center}
\url{https://github.com/RobinCodes/Projects/tree/main/Divisibility%20problem}.
\end{center}
The factorisations, the discrete--logarithm verifications, and the sieve data were carried out with \texttt{sympy}.

\subsubsection{Note on the use of artificial intelligence}\label{d2:app:ai}
This work was prepared with substantial assistance from an artificial-intelligence system, and most of the results presented here---both the structural theory of Sections~\ref{d2:sec:reductions}--\ref{d2:sec:repeated} and the computations of Section~\ref{d2:sec:computation} (the single--pass residue sweep, the extended multi--threaded sweep, and the two--prime factoring search)---were created with such assistance. We do not specify the particular system used; readers seeking further detail about the methods are invited to contact the authors: Robin Rovenszky (\texttt{robincodes} on Discord) and notxxdog (\texttt{.notxxdog} on Discord).

No result here is accepted on the basis of machine assertion. Every definition, lemma, and proof is stated so as to be checkable by hand, and every computational claim is reproducible by re-running the programs of Appendix~\ref{d2:app:programs}.

\begin{thebibliography}{9}

\bibitem{d2:prior}
R.~Rovenszky and notxxdog, \emph{On the congruences $2^{n}\equiv-1$ and $2^{n}\equiv-3\pmod{n+1}$}, companion manuscript.

\bibitem{d2:oeis}
OEIS Foundation Inc., entry A245728 in \emph{The On--Line Encyclopedia of Integer Sequences}, \url{https://oeis.org/A245728}.

\bibitem{d2:cipolla}
M.~Cipolla, \emph{Sui numeri composti $P$ che verificano la congruenza di Fermat $a^{P-1}\equiv1\pmod P$}, Annali di Matematica Pura ed Applicata \textbf{9} (1904), 139--160.

\bibitem{d2:pomerance}
C.~Pomerance, J.~L.~Selfridge, S.~S.~Wagstaff, Jr., \emph{The pseudoprimes to $25\cdot10^{9}$}, Math.\ Comp.\ \textbf{35} (1980), 1003--1026.

\bibitem{d2:agp}
W.~R.~Alford, A.~Granville, C.~Pomerance, \emph{There are infinitely many Carmichael numbers}, Ann.\ of Math.\ \textbf{139} (1994), 703--722.

\bibitem{d2:hooley}
C.~Hooley, \emph{On Artin's conjecture}, J.\ Reine Angew.\ Math.\ \textbf{225} (1967), 209--220.

\end{thebibliography}


\section{Jackopedia}

A growing collection of self-contained problems studied in the same spirit as those above.
The first concerns how long a doubling sequence can run before it is forced to display a fixed decimal pattern.

\subsection{Avoidance}

\begin{abstract}
Start at $1$ and repeatedly apply $x\mapsto bx+c$ with a digit $c\in\{0,\dots,b-1\}$; this generates every positive integer exactly once, along its base-$b$ expansion. Fix a target base $B$ and a digit-string $n$. We ask how long such a generation can proceed while the \emph{base-$B$} representation of every value produced avoids $n$ as a contiguous factor. For $(b,B)=(2,10)$ this is precisely the longest binary string, with no leading zero, all of whose prefixes read in decimal avoid $n$.

We separate three quantities. The literal \emph{shortest} string reaching a value that contains $n$ is shown to be a triviality of base conversion (Proposition~\ref{av:prop:short}). The substantive objects are the longest avoider $\Lbb(n)$ and its true dual, the shortest \emph{maximal} avoider $\hbb(n)$. We give exact data for $(2,10)$ with single-digit $n$, and our main structural result is a dichotomy governed by whether $\log b/\log B$ is rational. When $b$ and $B$ are multiplicatively dependent we exhibit explicit eventually-periodic words proving $\Lbb=\infty$ (Theorem~\ref{av:thm:dep}); when they are independent we prove the opposite---$\Lbb(n)<\infty$ for every $n$ (Theorem~\ref{av:thm:dichotomy})---by showing that the leading base-$B$ digits of any chain equidistribute and hence display every pattern. This completes a dichotomy governed entirely by whether $\log b/\log B$ is rational, and exposes the finiteness as an \emph{arithmetic} phenomenon (incommensurable bases), in deliberate contrast to Thue--Morse-style word avoidance.

Finally we study a self-referential variant in which each term must avoid all of its own predecessors (Section~\ref{av:sec:variant}). A short reduction (Proposition~\ref{av:prop:variant}) ties its finiteness to the same avoidance question, single-digit starts are solved exactly, and we isolate the accumulating self-constraints---rather than avoidance of the start---as the mechanism that bounds the sequence.
\end{abstract}

\subsubsection{The generation tree}

Fix integers $b,B\ge 2$. We call $b$ the \emph{generation base} and $B$ the \emph{reading base}.

\begin{definition}
The \emph{generation tree} $T_b$ is the rooted infinite tree with root $1$ in which each node $x$ has the $b$ children $bx,\,bx+1,\dots,bx+(b-1)$. A \emph{generation string} of length $k$ is a path $x_1,x_2,\dots,x_k$ from the root, equivalently a word $d_1d_2\cdots d_k$ over $\{0,\dots,b-1\}$ with $d_1\neq 0$ and $x_i=\val(d_1\cdots d_i)=\sum_{j\le i} d_j\,b^{\,i-j}$.
\end{definition}

Thus generation strings are exactly the base-$b$ representations of positive integers, and a length-$k$ path corresponds to a $k$-digit base-$b$ number. For $b=2$ the two operations are $x\mapsto 2x$ (append $0$) and $x\mapsto 2x+1$ (append $1$): the ``doubling'' process of the title.

Let $\rho_B(x)\in\{0,\dots,B-1\}^{*}$ denote the base-$B$ digit word of $x$ (most significant digit first, no leading zero for $x\ge 1$). A \emph{pattern} is a nonempty word $n=n_1\cdots n_\ell$ over $\{0,\dots,B-1\}$ with $n_1\neq 0$; we identify it with the integer $\nu:=\base{n}{B}=\sum_{i}n_i B^{\,\ell-i}$. Patterns are exactly the base-$B$ representations of positive integers, matching the convention $n\in\ZZ_{\ge1}$.

\begin{definition}
An integer $x\ge 1$ is \emph{good} (for $n,B$) if $n$ is not a factor of $\rho_B(x)$, and \emph{bad} otherwise. A generation string $x_1,\dots,x_k$ is \emph{valid} if every $x_i$ is good.
\end{definition}

Because the parent of $x$ in $T_b$ is $\lfloor x/b\rfloor$, validity of a path is equivalent to all of its ancestors being good. Hence a node $x$ lies on some valid path iff $x$ and all of $\lfloor x/b\rfloor,\lfloor x/b^2\rfloor,\dots,1$ are good. The valid nodes form a subtree of $T_b$ rooted at $1$ (assuming $1$ is good), which we call the \emph{avoidance tree} $A=A_{b,B}(n)$.

\begin{definition}
We study three quantities.
\begin{enumerate}[label=\textup{(\roman*)},leftmargin=2.2em]
\item $\Lbb(n)=\sup\{k:\text{a valid string of length }k\text{ exists}\}\in\NN\cup\{\infty\}$, the \emph{longest avoider}: the height of $A$.
\item $\Sbb(n)=\min\{k:\text{a generation string of length }k\text{ ends at a bad value}\}$, the \emph{shortest reach}.
\item $\hbb(n)=\min\{k:\text{a maximal valid string of length }k\text{ exists}\}$, the \emph{shortest maximal avoider}, where a valid string is \emph{maximal} if none of its $b$ one-letter extensions is valid (equivalently the terminal node is a leaf of $A$).
\end{enumerate}
\end{definition}

For the motivating instance $(b,B)=(2,10)$, $\Lbb(n)$ is the length of the longest binary string with no leading zero whose every prefix, written in decimal, avoids $n$. For example $\Lbb(1)=0$ (the prefix $\rho_{10}(1)=\texttt{1}$ already contains $\texttt{1}$), while $\Lbb(2)=6$ is witnessed by $\texttt{111100}=\val^{-1}(60)$ with prefix values $1,3,7,15,30,60$, none containing the digit $\texttt{2}$.

\subsubsection{The shortest problem is base conversion}

We first dispatch the shortest-reach quantity, which turns out to carry no arithmetic content; the genuinely dual notion is $\hbb$, treated in Section~\ref{av:sec:dual}.

\begin{lemma}[Reachability]\label{av:lem:reach}
Every integer $m\ge 1$ is the value of exactly one generation string, namely $\rho_b(m)$, whose length is $\lfloor\log_b m\rfloor+1$.
\end{lemma}

\begin{proof}
Immediate: $T_b$ realizes the standard base-$b$ numeration, and the number of base-$b$ digits of $m$ is $\lfloor\log_b m\rfloor+1$.
\end{proof}

\begin{proposition}[Shortest reach]\label{av:prop:short}
For every admissible pattern $n$,
\[
\Sbb(n)=\lfloor\log_b \nu\rfloor+1,\qquad \nu=\base{n}{B}.
\]
Moreover $\Sbb(n)$ is exactly the depth of the shallowest bad node, and that node is reachable through good ancestors.
\end{proposition}

\begin{proof}
The smallest bad integer is $\nu$ itself. Indeed, if $m$ is bad then $\rho_B(m)$ contains $n$, so either $\len_B(m)=\ell$ and $\rho_B(m)=n$, giving $m=\nu$, or $\len_B(m)>\ell$ and $m\ge B^{\ell}>\nu$ (as $\nu<B^{\ell}$). By Lemma~\ref{av:lem:reach} the unique generation string reaching $\nu$ has length $\lfloor\log_b\nu\rfloor+1$, and any value at depth $d\le \lfloor\log_b\nu\rfloor$ is $<b^{d}\le\nu$, hence good. Thus no bad node occurs above this depth, the bound is attained, and every proper ancestor of $\nu$ (being $<\nu$) is good.
\end{proof}

So ``reach a value containing $n$ in the fewest bits'' is solved by writing $n$ in base $B$, taking its integer value $\nu$, and counting its base-$b$ digits: for $(2,10)$ one simply has $\Sbb(n)=\lfloor\log_2 n\rfloor+1$, the bit-length of $n$. Proposition~\ref{av:prop:short} also pinpoints the contrast that makes Section~\ref{av:sec:dual} interesting: $\Sbb$ is the first depth at which one \emph{can} hit $n$, whereas $\hbb$ is the first depth at which one is \emph{forced} to.

\subsubsection{The longest avoider and the role of the bases}

\subsubsection{The same base: infinite by word combinatorics}

\begin{theorem}[Same base]\label{av:thm:same}
Let $b=B$. A generation string is valid iff its underlying base-$b$ word avoids $n$ as a factor. Consequently $\Lbb[b,b](n)=\infty$ for every admissible pattern $n$, with the single exception $(b,n)=(2,\texttt{1})$, where $\Lbb[2,2](\texttt{1})=0$.
\end{theorem}

\begin{proof}
When $B=b$ the value $x_i=\val(d_1\cdots d_i)$ has base-$b$ word exactly $d_1\cdots d_i$ (no leading zero, as $d_1\neq0$). Hence $x_i$ is good iff the prefix $d_1\cdots d_i$ avoids $n$, and the path is valid iff every prefix---equivalently the whole word---avoids $n$. Thus $\Lbb[b,b](n)=\infty$ iff there is an infinite word over $\{0,\dots,b-1\}$ with nonzero first letter avoiding the factor $n$.

Such a word exists in all cases but one. If $n$ contains two distinct letters, the constant word $a^{\infty}$ for any $a\in\{1,\dots,b-1\}$ avoids it. If $n=d^{\ell}$ is a power of a single letter $d$: when some nonzero letter $a\neq d$ exists (always, unless $b=2$ and $d=1$) the word $a^{\infty}$ contains no $d$ at all and so avoids $n$; when $b=2,d=1$ and $\ell\ge 2$, the word $\texttt{1}\,\texttt{0}^{\infty}$ has a single $\texttt{1}$ and avoids $\texttt{1}^{\ell}$. The remaining case $b=2,\,n=\texttt{1}$ forces the (necessarily nonzero) leading digit to be the forbidden letter, so no nonempty valid string exists and $\Lbb[2,2](\texttt{1})=0$.
\end{proof}

\subsubsection{Data for $(b,B)=(2,10)$}

For single-digit $n$ the avoidance tree $A_{2,10}(n)$ is finite, and an exhaustive breadth-first traversal to extinction both computes $\Lbb(n)$ and \emph{certifies} finiteness for that $n$. Table~\ref{av:tab:2-10} lists the results together with the (trivial) shortest reach, the shortest maximal avoider of Section~\ref{av:sec:dual}, and a witnessing longest string.

\begin{table}[h]
\centering
\begin{tabular}{cccccl}
\toprule
$n$ & $\Sbb(n)$ & $\hbb(n)$ & $\Lbb(n)$ & terminal value of a longest avoider \\
\midrule
$1$ & $1$ & $0$  & $0$  & --- (root already bad) \\
$2$ & $2$ & $4$  & $6$  & $60=\val(\texttt{111100})$ \\
$3$ & $2$ & $5$  & $8$  & $160$ \\
$4$ & $3$ & $5$  & $11$ & $2000$ \\
$5$ & $3$ & $5$  & $10$ & $768$ \\
$6$ & $3$ & $5$  & $18$ & $182857$ \\
$7$ & $3$ & $6$  & $32$ & $3\,360\,819\,200$ \\
$8$ & $4$ & $6$  & $22$ & $4\,095\,999$ \\
$9$ & $4$ & $6$  & $56$ & $70\,465\,683\,456\,000\,000$ \\
\bottomrule
\end{tabular}
\caption{The three quantities for $(b,B)=(2,10)$ and single-digit $n$. Here $\Sbb(n)=\lfloor\log_2 n\rfloor+1$ is the bit-length of $n$ (Proposition~\ref{av:prop:short}). The longest avoider for $n=8$ rides the ladder $\dots\,7999\to15999\to31999\to\dots\to4{,}095{,}999$, after which every double contains an $\texttt{8}$.}
\label{av:tab:2-10}
\end{table}

The sequence $\bigl(\Lbb[2,10](n)\bigr)_{n\ge1}=0,6,8,11,10,18,32,22,56,\dots$ does not appear in the OEIS. We stress that each entry is a finiteness statement for one specific $n$; the traversal says nothing uniform, a point we return to in Section~\ref{av:sec:indep}.

\subsubsection{Multiplicative dependence forces infinitude}

We now turn to $b\neq B$. The decisive invariant is whether $\log b/\log B$ is rational. For integers $b,B\ge2$ this is equivalent to the existence of $p,q\ge1$ with $b^{q}=B^{p}$, and by unique factorization to
\[
b=c^{s},\quad B=c^{t}\qquad\text{for some integer }c\ge 2\text{ and }s,t\ge1;
\]
we then call $b,B$ \emph{multiplicatively dependent}. Otherwise they are \emph{independent} (e.g.\ $2$ and $10$, since $\log_{10}2\notin\QQ$).

In the dependent case both numerations are block recodings of base $c$: a base-$b$ digit is a block of $s$ base-$c$ digits, and a base-$B$ digit is a block of $t$ base-$c$ digits. This is enough to construct explicit infinite avoiders.

\begin{lemma}[Repunit avoider]\label{av:lem:repunit}
Let $b=c^{s}$, $B=c^{t}$. Consider the generation string $u^{\infty}$ where $u\in\{0,\dots,b-1\}$ is the base-$b$ digit whose base-$c$ block is $\underbrace{1\cdots1}_{s}$; that is, the value after $j$ generation steps is the base-$c$ repunit $R_{js}=\frac{c^{js}-1}{c-1}$. Then every base-$B$ digit occurring in any of these values lies in the finite set
\[
\mathcal R_t:=\Bigl\{\,R_r=\tfrac{c^{r}-1}{c-1}\ :\ 1\le r\le t\,\Bigr\}.
\]
\end{lemma}

\begin{proof}
$R_{js}$ has base-$c$ word $\underbrace{1\cdots1}_{js}$. Grouping it into $t$-blocks from the least significant end yields full blocks $\underbrace{1\cdots1}_{t}$, of base-$B$ digit value $R_t$, and (if $t\nmid js$) one leading partial block $\underbrace{1\cdots1}_{r}$ with $r=js\bmod t\in\{1,\dots,t-1\}$, of value $R_r$. Hence all base-$B$ digits lie in $\{R_1,\dots,R_t\}=\mathcal R_t$.
\end{proof}

\begin{theorem}[Dependent $\Rightarrow$ infinite]\label{av:thm:dep}
If $b,B\ge2$ are multiplicatively dependent, then there is an admissible pattern $n$ with $\Lbb(n)=\infty$.
\end{theorem}

\begin{proof}
Write $b=c^{s},B=c^{t}$. Suppose first $c^{t}>t+1$ (this excludes only $c=2,t=1$). The set $\mathcal R_t$ of Lemma~\ref{av:lem:repunit} has at most $t$ elements, all positive, while there are $B-1=c^{t}-1>t$ nonzero base-$B$ digits; hence some nonzero digit $d\notin\mathcal R_t$ exists. Take $n=[d]$, the single digit $d$. By Lemma~\ref{av:lem:repunit} the generation string $u^{\infty}$ never produces the digit $d$, so every prefix is good---in particular the root value $R_s\in\mathcal R_t$ is good---and $u^{\infty}$ is an infinite valid string. Thus $\Lbb(n)=\infty$.

The sole remaining case is $c=2,t=1$, i.e.\ $B=2$ and $b=2^{s}$. Take $n=\texttt{11}$. The binary word $(\texttt{10})^{\infty}$ contains no factor $\texttt{11}$; it is realizable as a base-$2^{s}$ generation string (its $s$-bit blocks form a fixed periodic base-$2^{s}$ digit sequence), and its root prefix $\texttt{10}$ is good. Hence $\Lbb(\texttt{11})=\infty$.
\end{proof}

\begin{remark}[Two flavours of infinitude]
Theorem~\ref{av:thm:dep} produces a \emph{thin} infinite avoider---a single periodic path. The full avoidance tree may be thin or thick. For $(b,B,n)=(2,4,\texttt{2})$ the only survivors are essentially the all-ones path $\texttt{1}^{\infty}$, whose values $2^{k}-1$ have base-$4$ digits in $\{\texttt{1},\texttt{3}\}=\mathcal R_2$ (here $c=2,t=2$), so the breadth-first frontier collapses to a single node at every large depth. For $(2,4,\texttt{3})$, by contrast, validity is the condition ``no aligned binary block $\texttt{11}$,'' the surviving strings are exponentially many (the frontier passes $3.5\times10^{6}$), and the avoider $(\texttt{10})^{\infty}$ of Theorem~\ref{av:thm:dep} is just one path among the throng.
\end{remark}

This thin-versus-thick contrast is not accidental: in the dependent case the whole problem is finite-state, and both flavours are read off a single automaton.

\begin{theorem}[The dependent case is automatic]\label{av:thm:auto}
Let $b=g^{s}$ and $B=g^{t}$ for integers $g\ge2$, $s,t\ge1$. Then $\{v\ge1:\rho_B(v)\text{ avoids }n\}$ is a $b$-recognizable set of integers. Consequently:
\begin{enumerate}\itemsep1pt
\item[(i)] the avoidance tree $A_{b,B}(n)$ is generated by a finite automaton, and $\Lbb(n)\in\{0,1,2,\dots\}\cup\{\infty\}$ is \emph{computable};
\item[(ii)] $\Lbb(n)=\infty$ iff that automaton has a directed cycle of live states reachable from the root;
\item[(iii)] the number of valid strings of length $k$ grows polynomially or exponentially in $k$, according as the live sub-automaton is cycle-disjoint (a thin path) or contains a branching cycle (an exponential frontier).
\end{enumerate}
\end{theorem}

\begin{proof}[Proof sketch]
A base-$b$ digit is a block of $s$ base-$g$ digits and a base-$B$ digit is a block of $t$ of them, so appending a base-$b$ digit ($v\mapsto bv+c$) and grouping into base-$B$ are both finite-state recodings of the base-$g$ digit stream. The pattern $n$ corresponds to a fixed base-$g$ word of length $\ell t$ to be matched at positions $\equiv0\ (\mathrm{mod}\ t)$; an automaton reading base-$g$ digits and tracking (a) the position modulo $t$, (b) the Aho--Corasick state of that word, and (c) a ``seen'' flag is finite. This is the standard inter-recognizability of bases that are powers of a common base, so the avoiding set is $g$-, hence $b$-, recognizable. A recognizable set's defining DFA makes the longest root-path, its finiteness, and its growth type decidable by inspection of the reachable live strongly connected components.
\end{proof}

This is exactly why the multiplicatively dependent regime is word-combinatorial---and why finiteness there is decidable, in contrast to the genuinely arithmetic independent case below.

\subsubsection{Independence forces finiteness: the dichotomy theorem}\label{av:sec:indep}

Every multiplicatively \emph{independent} instance is finite (Table~\ref{av:tab:cross} shows samples), in sharp contrast to the dependent cases. We now \emph{prove} this, completing the dichotomy. The mechanism is the distribution of \emph{leading} digits.

\begin{table}[h]
\centering
\begin{tabular}{llc l}
\toprule
$(b,B)$ & pattern $n$ & dependent? & $\Lbb(n)$ \\
\midrule
$(2,3)$  & \texttt{2}    & no ($\log_3 2\notin\QQ$) & $2$ \\
$(2,3)$  & \texttt{22}   & no & $12$ \\
$(3,10)$ & \texttt{7}    & no & $20$ \\
$(2,4)$  & \texttt{2}    & yes ($4=2^{2}$) & $\infty$ \quad(avoider $\texttt{1}^{\infty}$) \\
$(2,4)$  & \texttt{3}    & yes & $\infty$ \quad(avoider $(\texttt{10})^{\infty}$) \\
$(2,8)$  & \texttt{5}    & yes ($8=2^{3}$) & $\infty$ \\
$(3,9)$  & \texttt{5}    & yes ($9=3^{2}$) & $\infty$ \\
\bottomrule
\end{tabular}
\caption{Finiteness tracks multiplicative independence of the bases.}
\label{av:tab:cross}
\end{table}

The key is that the choices in a chain cannot steer its \emph{leading} digits, which are governed by an irrational rotation and therefore equidistribute.

\begin{theorem}[Leading-block occurrence]\label{av:thm:leading}
Let $b,B\ge2$ with $\theta:=\log_B b\notin\QQ$, and let $w$ be a nonempty base-$B$ word with nonzero first digit, length $\ell$ and value $\nu=(w)_B$. Then for \emph{every} chain $x_0,x_1,\dots$ with $x_{k+1}\in\{bx_k,\dots,bx_k+(b-1)\}$, the base-$B$ representation of $x_k$ has $w$ as its leading $\ell$ digits for infinitely many $k$. In particular every infinite such chain contains $w$, so $L^{*}_{b,B}(w)<\infty$, and a fortiori $\Lbb(n)<\infty$ for $n=(w)_B$.
\end{theorem}

\begin{proof}
\emph{A chain is a single real number.} Since $x_{k+1}/b^{k+1}=x_k/b^{k}+c_k/b^{k+1}$ with $c_k\in\{0,\dots,b-1\}$, the sequence $x_k/b^{k}$ is nondecreasing and bounded above by $x_0+1$, hence converges to some real $\alpha\ge x_0\ge1$. Writing $\rho_k:=\log_B(x_k/b^{k})-\log_B\alpha$, we have $\rho_k\to0$ and
\[
\log_B x_k \;=\; k\theta+\log_B\alpha+\rho_k .
\]

\emph{Leading digits are a mantissa condition.} An integer $N$ with at least $\ell$ base-$B$ digits has leading $\ell$ digits equal to $w$ iff $N\in[\nu B^{j},(\nu+1)B^{j})$ for $j=\lfloor\log_B N\rfloor-(\ell-1)\ge0$, i.e.\ iff
\[
\{\log_B N\}\in I_w:=\bigl[\log_B\nu-(\ell-1),\ \log_B(\nu+1)-(\ell-1)\bigr)\pmod1,
\]
an arc of length $|I_w|=\log_B\frac{\nu+1}{\nu}>0$ (interpreted cyclically when $\nu+1=B^{\ell}$).

\emph{Equidistribution.} As $\theta\notin\QQ$, Weyl's theorem makes $(k\theta+\log_B\alpha)_{k\ge0}$ equidistributed mod $1$. Fix a closed arc $J\subset\operatorname{int}I_w$ of positive length and $\eta>0$ with the $\eta$-neighbourhood of $J$ contained in $I_w$. Because $\rho_k\to0$, there is $K_0$ with $|\rho_k|<\eta$ for all $k\ge K_0$; for such $k$, $\{k\theta+\log_B\alpha\}\in J$ implies $\{\log_B x_k\}\in I_w$. By equidistribution the former holds for a set of $k$ of density $|J|>0$, hence for infinitely many $k\ge K_0$. For all large such $k$ the integer $x_k$ has at least $\ell$ digits (as $x_k\ge b^{k}x_0\to\infty$), so its leading $\ell$ digits are exactly $w$. Thus $w$ occurs in $x_k$ for infinitely many $k$.

Consequently no infinite chain avoids $w$. The avoidance tree is finitely branching and has no infinite path, so by K\H{o}nig's lemma it is finite; that is, $L^{*}_{b,B}(w)<\infty$.
\end{proof}

The phase $\log_B\alpha$ is the only freedom a chain has, and an irrational rotation equidistributes from \emph{every} phase: there is no escaping starting point. Combined with Theorem~\ref{av:thm:dep}, this settles the dichotomy.

\begin{theorem}[Dichotomy]\label{av:thm:dichotomy}
$\Lbb(n)<\infty$ for every admissible pattern $n$ \emph{if and only if} $\log b/\log B\notin\QQ$.
\end{theorem}

\begin{proof}
If $\log b/\log B\notin\QQ$, Theorem~\ref{av:thm:leading} gives $\Lbb(n)\le L^{*}_{b,B}(n)<\infty$ for every $n$. If $\log b/\log B\in\QQ$ then $b,B$ are multiplicatively dependent and Theorem~\ref{av:thm:dep} exhibits an $n$ with $\Lbb(n)=\infty$.
\end{proof}

The argument is effective: the leading block visits $I_w$ with density $|I_w|$, so a chain can dodge $w$ in its most significant digits only until the rotation by $\theta$ has swept the circle.

\begin{corollary}[Quantitative bound]\label{av:cor:quant}
$L^{*}_{b,B}(w)\le K^{*}(w)$, the covering time of $I_w$ under rotation by $\theta$, i.e.\ the least $K$ with $\bigcup_{0\le k\le K}(I_w-k\theta)\equiv[0,1)\ (\mathrm{mod}\ 1)$. One has $K^{*}(w)=O_\theta\!\bigl(1/|I_w|\bigr)$ and $1/|I_w|=\bigl(\log_B\tfrac{\nu+1}{\nu}\bigr)^{-1}\sim\nu\ln B$, so $\Lbb(w)=O_\theta(\nu)$: at most exponential in the length of $w$.
\end{corollary}

For $(2,10)$ this is borne out sharply. The single-digit data satisfy $L_{2,10}(w)\le K^{*}(w)$ throughout,
\[
\begin{array}{c|cccccccc}
w & 2&3&4&5&6&7&8&9\\\hline
L_{2,10}(w) & 6&8&11&10&18&32&22&56\\[1pt]
K^{*}(w) & 9&9&12&32&42&52&62&62
\end{array}
\]
The two-digit covering times are likewise finite---$72,82,92,288$ for $w=\texttt{10},\texttt{12},\texttt{24},\texttt{99}$---which is what makes those (computationally unreachable) avoidance trees finite. The bound $K^{*}(w)\asymp\nu$ is, however, far from the true size of $\Lbb$: Section~\ref{av:sec:bounds} gives evidence that $\Lbb$ grows only like $\sqrt\nu$, so $K^{*}$ overshoots by a further factor $\sim\sqrt\nu$. Already among single digits the ratio $L_{2,10}(w)/K^{*}(w)$ drifts down from $0.90$ at $w=\texttt{9}$ to $0.31$ at $w=\texttt{5}$.

\subsubsection{How large is the avoider?}\label{av:sec:bounds}

The covering time $K^{*}(w)$ is itself a clean Diophantine quantity, and admits an explicit description through the continued fraction of $\theta$.

\begin{proposition}[Effective covering time]\label{av:prop:cover}
Write $\theta=\log_{10}2=[0;3,3,9,2,2,4,6,2,1,1,3,\dots]$, with convergent denominators $q_j=1,3,10,93,196,485,2136,\dots$, and for a word $w$ of value $\nu$ let $q_j<1/|I_w|\le q_{j+1}$. Then
\[
1/|I_w|\ \le\ K^{*}(w)\ \le\ q_j+q_{j+1},\qquad\text{so}\qquad K^{*}(w)=\Theta(q_{j+1}).
\]
Since $\theta=\log2/\log10$ has \emph{finite} irrationality measure $\mu$ (Baker's theorem on linear forms in logarithms), $q_{j+1}=O\!\bigl(q_j^{\,\mu-1+o(1)}\bigr)$, whence
\[
\Lbb[2,10](w)\ \le\ K^{*}(w)\ =\ O\!\bigl(\nu^{\,\mu-1+o(1)}\bigr).
\]
Conjecturally $\mu=2$, giving the near-linear bound $\nu^{1+o(1)}$.
\end{proposition}

\begin{proof}
The lower bound $1/|I_w|\le K^{*}(w)$ is the volume count $(K^{*}+1)|I_w|\ge1$. For the upper bound, $\bigcup_{k\le K}(I_w-k\theta)=[0,1)$ holds once the $K+1$ points $\{k\theta\}_{k\le K}$ leave no gap exceeding $|I_w|$. By the three-distance theorem the gaps among $\{k\theta\}_{k<N}$ take at most three values, and once $N\ge q_j+q_{j+1}$ every gap is at most $\lVert q_j\theta\rVert\le1/q_{j+1}\le|I_w|$; hence $K^{*}(w)\le q_j+q_{j+1}$. The relation $|\theta-p/q|>q^{-\mu-o(1)}$ gives $q_{j+1}\asymp(a_{j+1})q_j$ with $a_{j+1}=O(q_j^{\mu-2+o(1)})$, i.e.\ $q_{j+1}=O(q_j^{\mu-1+o(1)})$; with $q_j<1/|I_w|\asymp\nu$ this is $O(\nu^{\mu-1+o(1)})$.
\end{proof}

The single large partial quotient $a_3=9$ is visible in the data: it makes $q_3=93$ govern every $w$ with $10<1/|I_w|\le93$ (values $5\le\nu\lesssim37$), which is why $K^{*}$ sits on a plateau near $52$--$102$ there, and $w=\texttt{73}$ saturates the bound exactly, $169\le288\le196+93=289$.

A matching lower bound on $\Lbb$ itself is harder; a first-moment count gives the square root of the truth.

\begin{proposition}[Lower bound]\label{av:prop:lower}
$\Lbb[2,10](w)=\Omega(\sqrt{\nu})$.
\end{proposition}

\begin{proof}
Work in $[1,2)$, so $x_0=1$, and set $x_k=\lfloor2^{k}\alpha\rfloor$. As $\alpha$ ranges over $[1,2)$ the value $x_k$ takes each integer in $[2^{k},2^{k+1})$ on a subinterval of length $2^{-k}$, so
\[
\bigl|\{\alpha\in[1,2): w\subseteq\rho_{10}(x_k)\}\bigr|=2^{-k}\,\#\{m\in[2^{k},2^{k+1}):w\subseteq\rho_{10}(m)\}.
\]
The integer $m$ has $D_k=O(k)$ decimal digits; fixing the $\ell=\len(w)$ digits of any one window leaves the rest constrained to a band of length $2^{k}$, so each of the $\le D_k$ window positions admits at most $2^{k}10^{-\ell}+1$ values of $m$. For $k$ with $2^{k}\ge10^{\ell}$ this is $\le2\cdot2^{k}10^{-\ell}$, and since $10^{-\ell}<1/\nu$ the displayed measure is $<2D_k/\nu$. Summing over the steps $k\le K$ that can contain $w$ at all,
\[
\bigl|\{\alpha\in[1,2):\ \exists k\le K,\ w\subseteq\rho_{10}(x_k)\}\bigr|\ \le\ \sum_{k\le K}\frac{2D_k}{\nu}\ \le\ \frac{C\,K^{2}}{\nu}.
\]
Taking $K=\lfloor\sqrt{\nu/2C}\rfloor$ makes this $<1$, so some $\alpha\in[1,2)$ keeps every $x_k$ ($k\le K$) free of $w$; the corresponding length-$K$ string from $1$ avoids $w$, giving $\Lbb[2,10](w)\ge K=\Omega(\sqrt\nu)$.
\end{proof}

Thus $\Omega(\sqrt{\nu})\le\Lbb[2,10](w)\le O(\nu^{1+o(1)})$, a gap of a square root. Numerically the truth sits at the \emph{lower} end. Across four orders of magnitude ($\nu$ from $5$ to $5\times10^{4}$), beam and greedy searches---both of which only \emph{under}estimate $\Lbb$, as the exact single-digit maximum $56$ already exceeds what greedy finds---give $\Lbb[2,10](w)/\sqrt\nu$ essentially constant near $18$, with fitted exponent $0.50$--$0.55$; widening the beam $250$-fold barely changes the length, so the search is not the bottleneck. We accordingly revise the expectation.

\begin{conjecture}\label{av:conj:sqrt}
$\Lbb[2,10](w)=\nu^{1/2+o(1)}$. The first-moment bound of Proposition~\ref{av:prop:lower} is sharp in the exponent, while the covering bound $K^{*}\asymp\nu$ overestimates by a factor $\sim\sqrt\nu$.
\end{conjecture}

The matching upper bound is the open half, and the right picture is a coupon-collector one: a length-$K$ chain displays $\sim0.15K^{2}$ length-$\ell$ digit-windows, and once these exhaust the $M=9\cdot10^{\ell-1}\asymp\nu$ possible strings---around $K\asymp\sqrt{\nu\log\nu}$---the pattern $w$ can no longer be dodged. The naive alternative, treating the $2^{K}$ chains as \emph{independently} dodging $w$ with per-window odds $10^{-\ell}$, predicts survival to $K\asymp\nu$; this is what an earlier draft of ours mistook for the truth, and it fails because the digit-windows of a doubling chain are strongly correlated---successive terms share all but their last digit---so the chains carry far fewer than $2^{K}$ effective degrees of freedom. Making the coupon-collector count rigorous needs an equidistribution input we do not have: that the windows of \emph{every} surviving path cover the $\ell$-strings.

\begin{remark}[The Thue--Morse contrast, resolved]
In pure combinatorics on words, avoiding a pattern forever is the rule: a single factor over an alphabet of size $\ge2$ is avoidable by an eventually periodic word, and the Thue--Morse word $t=0110100110010110\cdots$ avoids even cubes over $\{0,1\}$ indefinitely. The same-base and multiplicatively-dependent cases (Theorems~\ref{av:thm:same}, \ref{av:thm:dep}) are exactly the regimes in which our problem \emph{is} word-combinatorial, and they are infinite for that reason; the avoiders $\texttt{1}^{\infty}$ and $(\texttt{10})^{\infty}$ are the humblest eventually-periodic relatives of $t$.

Theorem~\ref{av:thm:leading} explains why incommensurable bases behave oppositely. One might expect the $\pm c$ choices to grant control of the leading digits, but they do not: the children $bx,\dots,bx+(b-1)$ share their leading block, so the most significant digits evolve by the rigid rotation $m\mapsto\{m+\theta\}$ irrespective of the choices, and an infinite chain commits to a single phase $\log_B\alpha$. Irrationality of $\theta$ then forces (Weyl, Benford) that block through every value, $w$ included. The finiteness is genuinely \emph{arithmetic}---it is the incommensurability $\theta\notin\QQ$, not any word-combinatorial obstruction, that defeats avoidance---and it acts already in the leading digits, sidestepping the far harder open questions about the \emph{full} digit strings of $b^{k}$ (such as whether $2^{k}$ omits a fixed digit for only finitely many $k$, which asks for occurrences at \emph{all} large $k$ rather than the infinitely-many we need).
\end{remark}

\subsubsection{The shortest maximal avoider}\label{av:sec:dual}

The honest dual of ``longest avoider'' is not ``shortest reach'' (Proposition~\ref{av:prop:short}, a base conversion) but the shallowest place where the process is \emph{trapped}.

\begin{definition}
A valid string is \emph{maximal} if all $b$ of its one-letter extensions are bad; equivalently its terminal node is a leaf of the avoidance tree $A$. The lengths of maximal valid strings are exactly the leaf depths of $A$, and they all lie in the interval $[\hbb(n),\Lbb(n)]$.
\end{definition}

\begin{proposition}\label{av:prop:sandwich}
For every admissible $n$ with good root, $\Sbb(n)\le \hbb(n)\le \Lbb(n)$.
\end{proposition}

\begin{proof}
$\Lbb$ is the maximum leaf depth and $\hbb$ the minimum, so $\hbb\le\Lbb$. A maximal string of length $\hbb$ has all children bad; a bad node has depth at least $\Sbb$ by Proposition~\ref{av:prop:short}, whence the children, at depth $\hbb$, satisfy $\hbb\ge\Sbb$.
\end{proof}

For $(2,10)$ the values of $\hbb$ are recorded in Table~\ref{av:tab:2-10}: $0,4,5,5,5,5,6,6,6$ for $n=1,\dots,9$. They are witnessed by short trapped strings, e.g.\ for $n=2$ the string $\texttt{1101}=\val^{-1}(13)$ with prefix values $1,3,6,13$ (all avoiding $\texttt{2}$) whose children $26,27$ both contain $\texttt{2}$; and for $n=9$ the string $\val^{-1}(45)$ with prefix values $1,2,5,11,22,45$ whose children $90,91$ both contain $\texttt{9}$.

The two endpoints behave very differently. The shortest maximal avoider $\hbb$ grows slowly---it hugs $\Sbb=\Theta(\log\nu)$---while the longest avoider $\Lbb$ is erratic and large (already $56$ at $n=9$). Their gap $\Lbb-\hbb$ measures how much room the pattern leaves for steering: how long the $\pm c$ choices can postpone the inevitable once it becomes possible. The slow growth of $\hbb$ can be pinned down exactly.

\begin{theorem}[Order of the dual]\label{av:thm:dual}
$\hbb[2,10](n)=\Theta(\log\nu)$; more precisely $\hbb[2,10](n)=\Sbb(n)+O(1)$.
\end{theorem}

\begin{proof}
The lower bound $\hbb\ge\Sbb=\Theta(\log\nu)$ is Proposition~\ref{av:prop:sandwich}. For the upper bound, consider the five consecutive values $v\in\{5n,5n+1,\dots,5n+4\}$. Each has $2v\in\{10n,10n+2,\dots,10n+8\}$ and $2v+1\in\{10n+1,\dots,10n+9\}$, all of which begin with the digits of $n$; so for every one of these $v$, \emph{both} children contain $n$ as a leading block. Their common prefix---all digits but the last---is $\lfloor 5n/10\rfloor=\lfloor n/2\rfloor$, and $n$ is never a factor of $\lfloor n/2\rfloor$ (which is $<n$ and has no more digits), so an occurrence of $n$ inside any of the five must use the unit digit; this pins that digit, so at most one of the five contains $n$. Hence at least four of $\{5n,\dots,5n+4\}$ avoid $n$ while having both children trapped. Such a $v$ has $\lfloor\log_2 v\rfloor+1=\Sbb(n)+O(1)$ bits (indeed $\log_2 5\approx2.32$). It remains to route to $v$ along an $n$-free path: the binary ancestors $\lfloor v/2^{i}\rfloor$ with at least $\ell$ digits are the $O(1)$ values $\approx\tfrac52n,\tfrac54n$ (all smaller ones are too short to contain $n$), and for leading digit of $n$ at most $3$ these have exactly $\ell$ digits and exceed $n$, so cannot contain it; the choice of $r\in\{0,\dots,4\}$ clears the remaining leading-digit cases. One checks this succeeds for every $n\le99$, with $\hbb-\Sbb\in\{2,3\}$ throughout (mean $2.23$), matching the $\log_2 5$ signature.
\end{proof}

So the two endpoints have honestly different orders: $\hbb=\Theta(\log\nu)$ against $\Lbb=\nu^{\Theta(1)}$ (Section~\ref{av:sec:bounds}). The witness for $\hbb$ is the universal trapped block $\{5n,\dots,5n+4\}$, whose children all open with $n$.

\subsubsection{A self-referential variant: avoiding one's own past}\label{av:sec:variant}

We now change the rules so that the forbidden set is not fixed in advance but is \emph{generated by the sequence itself}.

\begin{definition}
Fix $n\ge 1$ and set $s_0=n$. A \emph{self-avoiding doubling sequence} is a sequence $s_0,s_1,\dots,s_{m}$ with $s_{k+1}\in\{2s_k,\,2s_k+1\}$ such that for every $k\ge 1$, the decimal word $\rho_{10}(s_k)$ contains \emph{none} of $\rho_{10}(s_0),\dots,\rho_{10}(s_{k-1})$ as a factor. Write $V(n)$ for the maximal length $m+1$ of such a sequence.
\end{definition}

Two features distinguish this from the fixed-pattern problem. First, the constraint tightens over time: the $k$-th term must dodge all $k$ of its predecessors. Second, since $s_{k+1}\ge 2s_k>s_k$ the sequence is strictly increasing, so a predecessor (being smaller, hence no longer) can in principle occur inside a successor; it is exactly these occurrences that are banned. As before this is a longest-path problem in the binary tree rooted at $n$, but now node-validity depends on the whole ancestral path, not on a fixed word.

\subsubsection{Single-digit starts}

For $1\le n\le 9$ an exhaustive traversal terminates, proving finiteness and giving the exact values of Table~\ref{av:tab:variant}.

\begin{table}[h]
\centering
\begin{tabular}{ccl}
\toprule
$n$ & $V(n)$ & a longest self-avoiding sequence \\
\midrule
$1$ & $4$  & $1,2,4,9$ \\
$2$ & $9$  & $2,4,9,18,37,75,151,303,607$ \\
$3$ & $10$ & $3,7,14,28,56,112,225,451,902,1805$ \\
$4$ & $10$ & $4,8,16,33,66,132,265,531,1063,2127$ \\
$5$ & $10$ & $5,11,23,46,93,187,374,749,1499,2999$ \\
$6$ & $13$ & $6,13,27,55,110,221,443,887,1774,3549,7099,14199,28399$ \\
$7$ & $27$ & $7,15,31,63,126,\dots,264328493,528656986$ \\
$8$ & $10$ & $8,17,34,69,139,279,559,1119,2239,4479$ \\
$9$ & $10$ & $9,18,37,75,151,303,607,1215,2431,4863$ \\
\bottomrule
\end{tabular}
\caption{Exact lengths $V(n)$ for single-digit starts. The value $n=7$ is a pronounced outlier; its sequence is essentially the pure doubling chain $7,15,31,63,126,\dots$, which happens to avoid the digit $\texttt{7}$ unusually long---echoing $\Lbb[2,10](7)=32$ from Table~\ref{av:tab:2-10}.}
\label{av:tab:variant}
\end{table}

\subsubsection{The reduction to fixed-pattern avoidance}

The behaviour of $V$ is governed by the longest-path problem of Section~3, through a single observation.

\begin{definition}
For a nonempty decimal word $w$, let $L^{*}(w)$ be the supremum of $m$ over all doubling chains $x_1,x_2,\dots,x_m$ (with $x_{i+1}\in\{2x_i,2x_i+1\}$ and \emph{any} start $x_1\ge1$) in which no $\rho_{10}(x_i)$ contains $w$ as a factor.
\end{definition}

$L^{*}(w)$ differs from $\Lbb[2,10](w)$ only in freeing the starting point: it is the longest doubling \emph{chain} avoiding $w$, rather than the longest such chain anchored at $1$.

\begin{proposition}[Variant reduction]\label{av:prop:variant}
For every $n\ge1$ the terms $s_1,\dots,s_{V(n)-1}$ form a doubling chain no term of which contains $\rho_{10}(n)$; hence
\[
V(n)\ \le\ 1+L^{*}\!\bigl(\rho_{10}(n)\bigr).
\]
In particular $V(n)<\infty$ for all $n$ whenever $L^{*}(w)<\infty$ for all $w$.
\end{proposition}

\begin{proof}
Each $s_k$ with $k\ge1$ must avoid every predecessor, in particular $s_0=n$, so $\rho_{10}(n)$ is not a factor of $\rho_{10}(s_k)$; and $s_{k+1}\in\{2s_k,2s_k+1\}$. Thus $(s_k)_{1\le k\le V(n)-1}$ is an admissible chain for $L^{*}(\rho_{10}(n))$, of length $V(n)-1$.
\end{proof}

\begin{corollary}[Finiteness of the variant]\label{av:cor:variant}
$V(n)<\infty$ for every $n$. Indeed $L^{*}(w)<\infty$ for every nonempty decimal word $w$ by Theorem~\ref{av:thm:leading} (the case $b=2$, $B=10$, $\theta=\log_{10}2\notin\QQ$), so Proposition~\ref{av:prop:variant} gives $V(n)\le 1+L^{*}(\rho_{10}(n))\le 1+K^{*}(\rho_{10}(n))=O(n)$ by Corollary~\ref{av:cor:quant}.
\end{corollary}

This is the start-free reading of Theorem~\ref{av:thm:dichotomy} for $(2,10)$: no doubling chain, from any start, avoids a fixed decimal string forever, because it cannot avoid that string in its own leading digits.

\subsubsection{The scale of \texorpdfstring{$V$}{V}, and what actually bounds it}

Proposition~\ref{av:prop:variant} predicts that $V(n)$ should grow with the length $\ell$ of $\rho_{10}(n)$: a longer start is a \emph{weaker} constraint, so $L^{*}$ is larger and there is more room. This is borne out. The single-digit maxima reach only $27$; two-digit starts run into the sixties (e.g.\ $V(24)\ge 67$, $V(13)\ge 63$); three-digit starts reach the hundreds ($V(250)\ge 203$, $V(777)\ge 274$); four-digit starts are larger still. The values listed for $\ell\ge2$ are lower bounds from beam search, in which the beam consistently \emph{dies out}---evidence that the true maxima, while large, are finite.

\begin{remark}[What bounds $V$: the start gives a ceiling, not the scale]
Every term after the first avoids $\rho_{10}(n)$, so by Theorem~\ref{av:thm:leading} the sequence cannot outlast the return of $\rho_{10}(n)$ as a \emph{leading} block, which the rotation by $\theta$ forces within $K^{*}(\rho_{10}(n))=O_\theta(1/|I_{\rho_{10}(n)}|)\sim n\ln10$ steps; this guarantees $V(n)\le1+K^{*}=O(n)$. The ceiling is tight only for the smallest starts---the longest chain avoiding $\texttt{10}$ alone has length $66$ against $K^{*}=72$, and $V(10)$ halts near $63$---but as $n$ grows the actual length falls far below it (the brackets below), held down instead by the same interior-occurrence effect that keeps $\Lbb$ near $\sqrt\nu$. The contrast $V(9)=10\ll L_{2,10}(9)=56$ is a start effect of a different kind: fixing $s_0=9$ pins the phase near $\log_{10}9.5$, from which the leading-$9$ arc returns after only ten steps, whereas a chain begun at $1$ may choose a phase postponing it to step $56$.
\end{remark}

The two competing forces are visible in two-digit brackets, $V(n)$ caught between a beam-search lower bound and $1+K^{*}(\rho_{10}(n))$:
\[
\begin{array}{c|ccccccc}
n & \texttt{10}&\texttt{17}&\texttt{24}&\texttt{37}&\texttt{50}&\texttt{73}&\texttt{99}\\\hline
\text{lower} & 63&72&69&65&83&93&134\\[1pt]
1+K^{*} & 73&83&93&103&196&289&289\\[1pt]
\text{ratio} & .88&.88&.75&.64&.43&.32&.47
\end{array}
\]
For small two-digit $n$ the bracket is tight and $V$ tracks the start's return time $K^{*}$ (start-avoidance dominates); as $n$ grows the ratio falls, the room $K^{*}\sim n\ln10$ outstripping the actual length, so the accumulating self-constraints begin to bite first. In particular $V$ is \emph{not} insensitive to $n$ within a digit-length: it rises with $n$, roughly as $1/|I_{\rho_{10}(n)}|$ until self-avoidance overtakes.

\begin{remark}[The single-digit outlier $V(7)=27$]
Among one-digit starts the leading-$d$ return time from the forced phase $\log_{10}(d{.}\,\cdots)$ rises with $d$: maximised over the start phase it is $4,10,10,13,33,43,53,63,63$ for $d=1,\dots,9$, and $V(d)$ never exceeds it. What makes $7$ exceptional is that it nearly \emph{achieves} its ceiling: the doubling orbit $7,15,31,63,126,\dots$ happens to be unusually free of both the digit $7$ and its own earlier terms, so the self-avoidance does little and $V(7)=27$ rides just under the leading-$7$ return time $53$. By contrast $8$ and $9$, with the larger ceiling $63$, are cut down to $V=10$ by interior collisions long before the leading block can return. The outlier is thus an arithmetic accident of the orbit of $7$, not a structural feature.
\end{remark}

\subsubsection{What remains open}

The structural picture is now complete: the dichotomy is a theorem (Theorems~\ref{av:thm:leading}, \ref{av:thm:dichotomy}), the dependent case is decidable and its growth dichotomy explicit (Theorem~\ref{av:thm:auto}), the dual has order $\Theta(\log\nu)$ (Theorem~\ref{av:thm:dual}), the covering time is pinned to the continued fraction of $\theta$ (Proposition~\ref{av:prop:cover}), and the self-referential variant is finite with $V(n)=O(n)$ (Corollary~\ref{av:cor:variant}). Three questions of differing character remain.

\begin{openproblem}[The square-root upper bound]
Prove the open half of Conjecture~\ref{av:conj:sqrt}: $\Lbb[2,10](w)=O\!\bigl(\sqrt\nu\,\mathrm{polylog}\,\nu\bigr)$, matching the first-moment lower bound. This means improving the covering bound $K^{*}\asymp\nu$ by a full factor $\sqrt\nu$; the natural route is to show that the length-$\ell$ digit-windows of any avoiding path longer than $\sim\sqrt{\nu\log\nu}$ are forced to exhaust the $\ell$-strings---a coupon-collector statement requiring control of how the windows of a doubling chain equidistribute.
\end{openproblem}

\begin{openproblem}[Exact values and the OEIS]
Compute $\Lbb[2,10](n)$ exactly for $\ell\ge2$. The avoidance tree is finite (depth $\le K^{*}(n)\le 288$ for two digits) but its frontier peaks beyond $10^{9}$, out of reach of the bitset traversal; a smarter certificate of extinction is needed. The single-digit row $0,6,8,11,10,18,32,22,56$ and its dual $0,4,5,5,5,5,6,6,6$ await a companion two-digit row before submission. The same holds for the exact $V(n)$, $\ell\ge2$, currently only bracketed.
\end{openproblem}

\begin{openproblem}[Full rigour for the dual and the variant scale]
Theorem~\ref{av:thm:dual} routes to a trapped value in $\{5n,\dots,5n+4\}$ and clears the $O(1)$ ancestor obstructions $\approx\tfrac52n,\tfrac54n$ rigorously for leading digit $\le3$ and by check for $n\le99$; give a uniform argument for all leading digits. Likewise determine the growth of $V(n)$: the data place it well below the start-return ceiling $K^{*}(\rho_{10}(n))\asymp n$, consistent with the same square-root-type law as $\Lbb$, and one would like both the exponent and the crossover at which intermediate self-avoidance, rather than the start, becomes the binding constraint.
\end{openproblem}

\bigskip
\noindent\textit{Computational note.}\ All finite values were obtained by traversing the avoidance tree level by level until a level is empty; an empty level is a constructive certificate that the tree is finite for that $n$. The bottleneck is the frontier size, which is erratic ($n=8$ dies with $58$ live nodes; $n=9$ peaks near $2.6\times10^{6}$). Representing a depth-$k$ frontier as a bitset over the offsets $o\in[0,b^{\,k-1})$ turns the transition into ``stretch each set bit into its $b$ children'' followed by masking against the bad-value set; this is efficient while the live set is a constant fraction of its band, whereas a sparse representation is preferable when the live set stays small but deep.


\appendix
\section{Computation}
All codes available at Github:\\

Divisibility problem: \url{https://github.com/RobinCodes/Projects/tree/main/Divisibility%20problem}\\

Binary-string problem: \url{https://github.com/RobinCodes/Projects/tree/main/Collatz-type%20problem/Collatz-type%20problem%20v3}\\

\section{Next Steps}
\begin{enumerate}
    \item bbchallenge
    \begin{itemize}
        \item Longitudinal analysis
\item Backwards longitudinal??? (decider)
\item New website
\\note feature ideas like Results > Papers, Results > Deciders
\item What are Laver tables and Goodstein sequences?
\item BB Bignum Bakeoff
\item Collect BB(2,5) holdouts that I started above somewhere in this Discord, then compile BB(6) in a spreadsheet similarly. I’m imagining a spreadsheet with the 3 peripheral domains (including BB(3,3)) as tabs and having these there.
\item Make a wiki page for these bb-like functions. different section for the ones for programming languages
    \end{itemize}
    \item Jack O' Pedia
    \begin{itemize}
        \item Also put avoidance into main.tex
    \end{itemize}
    \item Make a site for Project Ladislaus
\end{enumerate}

\end{document}
